\chapter{Brave-New Infinitesimal Symmetry and Noether Theory}
\label{ch:symmetry-noether}

\section{Purpose, continuity, and scope}
\label{sec:symmetry-noether-purpose-continuity-and-scope}

The preceding variational chapter constructed brave-new variational geometry before
any symmetry datum was introduced.  Its primary invariant was the complete
Cartan--vertical bifiltration, from which it constructed the coherent variational
differential, full and one-stage first variations, the universal Euler family,
derived critical loci, horizontal realization towers, relative Lepage realizations,
the universal presymplectic transgression, finite cubical variations, and, on the
represented smooth scalar sector, the critical Hessian and Jacobi linearization.
The differential action functional, where available, was constructed only after
those intrinsic variational objects were already in place.

The formal-groupoid, differentiation, Atiyah--Spencer, and native Lie-integration
chapters provide the independent infinitesimal-symmetry geometry used here.  A
brave-new formal Lie groupoid has a canonical nonlinear anchor into the maximal
relative infinitesimal relation, an effective formal quotient, formal isotropy, and
descended adjoint transport.  Spectral differentiation constructs intrinsic stable
coefficients, actual-isotropy controllers, homotopy anchor fibres, spectral
Artinian formal moduli, and the spectral partition-Lie algebroid controller.  The
Atiyah--Spencer chapter supplies the formalized relation, stable Schlessinger
reflection, tangent representations, and the higher operation calculus needed
below.  The native Lie-integration chapter supplies comparison and materialization
maps for those differentiated objects.  Those maps are retained with their proved
scope, but no inverse integration theorem is used in the construction of either
Noether theorem below.

The purpose of the present chapter is deliberately narrower than a theory of gauge
fields.  It develops the infinitesimal symmetry of a variational problem far enough
to state and prove brave-new forms of Noether's First and Second Theorems and to
construct the associated critical, presymplectic, Hamiltonian, and reduction
geometry.  The word ``symmetry'' below always refers to the formal-groupoid symmetry
of the descended field geometry.  No principal connection or curvature theory of a gauge field, finite-order
differential parameter operator, formal adjoint, integration-by-parts package,
superpotential, externally chosen reducibility complex, gauge fixing, BRST complex,
Koszul--Tate resolution, or BV structure is part of this chapter.  The differentiated
controller does, however, construct its intrinsic coefficient connection and, on the
locally coherent differentiated branch, the intrinsic pro-coherent realization space
of anchor splittings together with its curvature and Bianchi obstruction theory.  On
the represented finite-locally-free overlap, this intrinsic realization is identified
with the ordinary represented anchor-splitting space.  These are infinitesimal
controller constructions, not a theory of gauge fields.  Specific gauge theories and a general brave-new theory
of gauge fields are deferred to later chapters.

Throughout, fix
\[
x:X\longrightarrow S
\]
in the native big spectral Nisnevich $\infty$-topos and put
\[
Q:=Q_{X/S}.
\]
Let $(A,\rho_A)$ be a Cartan object for the brave-new jet comonad and write
\[
a:\overline A\longrightarrow Q
\]
for its constructed descent object.  Let
\[
M\in E_Q
\]
be a stable coefficient and let
\[
L:\mathbf1_Q\longrightarrow \operatorname{Sec}^0(A;M)[n]
\]
be a degree-$n$ Lagrangian.  Write
\[
\ell_L:\mathbf1_{\overline A}\longrightarrow a^*M[n]
\]
for its adjoint local section.

Retain from the variational chapter
\[
\operatorname{Var}_{X/S}(A;M),\qquad
V_M^\infty(A),\qquad
V_M^{(1)}(A),\qquad
V_M^{(k)}(A),
\]
the local Euler section
\[
E_L:\mathbf1_{\overline A}\longrightarrow V_M^{(1)}(A)[n],
\]
the full critical object $\operatorname{Crit}^{\infty}(L)$, the one-stage derived
critical object
\[
j_L:\operatorname{Crit}(L)\longrightarrow\overline A,
\]
the relative Lepage objects $\operatorname{Lep}^q(L)$, and the universal
presymplectic transgressions $\omega_L^{(q)}$.  Whenever the represented smooth
scalar comparison is available, retain also $\operatorname{Jac}_L$ and
$\operatorname{Hess}^{\mathrm{crit}}_L$.

The only new symmetry input is
\[
\mathcal G_\bullet\in
\operatorname{FormLieGpd}(\overline A/Q).
\]
Its anchor, quotient, first-order represented shadow, Artinian controller, isotropy,
and higher operations are all outputs of the preceding chapters.  In particular,
no gauge-parameter bundle or differential operator is introduced in order to state
either Noether theorem.

The chapter has four conceptual stages:
\[
\text{formal symmetry relation}
\longrightarrow
\text{symmetry--horizontal variational calculus}
\longrightarrow
\text{Noether current and native coherence},
\]
\[
\text{formalized Artinian symmetry}
\longrightarrow
\text{Schlessinger reflection}
\longrightarrow
\text{partition-Lie differentiation and Noether II},
\]
followed by the critical and mixed presymplectic/Hamiltonian constructions.  The
raw, formalized, represented first-order, and differentiated branches meet only by
the comparison morphisms actually constructed below.

\subsection{Chapter-level target}
\label{subsec:symmetry-noether-chapter-level-target}

The constructions below provide:
\begin{enumerate}
\item the canonical formal symmetry anchor, effective quotient, and constructed
      effective infinitesimal arrow-image subrelation;
\item the complete symmetry--horizontal bifiltration and its coherent double
      differential;
\item structured one- and finite-cubical symmetry variations, their right-Kan
      coefficients, and restriction from unrestricted variational coefficients;
\item the represented first-order symmetry-lift classification and its tangent
      shadow, used only on its represented theorem domain;
\item complete symmetry-realization objects at every horizontal stage and the
      universal brave-new Noether current;
\item the intrinsic off-shell First Noether homotopy, its one-stage Euler
      consequence, and the full Euler-nullification transgression object;
\item the complete native symmetry-coherence obstruction tower;
\item the quotient-formalized comparison, stable Schlessinger-reflected first
      Noether coefficient, relative unstable reflection of complete symmetry
      parameters on coherent formal leaves, and the spectral-\'etale exchange theorem
      for that relative unstable reflector;
\item the differentiated retractive partition-Lie algebroid, zero and Noether
      splittings, first Noether derivative, first syzygy, directional Taylor
      classes, and complete morphism corolla--latching hierarchy;
\item the differentiated Brave-New Second Noether theorem, both relative to the
      reflected symmetry parameter and in its universal strong section form;
\item the intrinsic Noether--Spencer splitting space, the independent tangent
      representation bar, the canonical Noether coefficient connection and its
      flatness/Bianchi hierarchy, the intrinsic pro-coherent anchor-splitting
      realization with curvature/Bianchi hierarchy, and its represented
      finite-locally-free comparison;
\item raw critical restriction followed by canonical formalization, the effective
      symmetry-reduced critical quotient, and its parameterized descent calculus;
\item on the represented smooth scalar degree-zero sector, the affine-point linear
      critical tangent, chartwise Jacobi fibre, and critical-Hessian radical
      factorization;
\item the mixed symmetry--vertical--horizontal calculus, ordered presymplectic
      contraction, Hamiltonian defect, Hamiltonian realization, residual
      transgression, and presymplectic-radical realization;
\item the exact represented first-order $(1,1)$-linearization of the ordered mixed
      transgression, its intrinsic contraction target, and the resulting represented
      presymplectic-radical factorization without an all-degree HKR identification;
\item complete presymplectic \v{C}ech descent to the reduced critical calculus,
      with Cartan basicness and null contraction derived as consequences, together
      with reduced presymplectic forms, reduced currents, and the Cartan-diagram
      realization needed to descend current transgressions with their selected
      nullhomotopies;
\item the exact functoriality, raw base-change, formal-\'etale locality, finite
      Nisnevich descent, and separately scoped spectral-\'etale functoriality of the
      formalized Artinian branch.
\end{enumerate}
\section{Audit of retained distinctions}
\label{sec:symmetry-noether-audit-of-retained-distinctions}

The constructions of this chapter use several coefficient theories which the
preceding chapters deliberately kept distinct.  Those distinctions remain in force.

On the variational side, the intrinsic full first variation and its split
square-zero restriction are related by
\[
\lambda^1_{\mathrm{loc}}:
V_M^\infty(A)\longrightarrow V_M^{(1)}(A),
\]
but they are not identified in general.  Likewise
$\operatorname{Crit}^{\infty}(L)$ kills the full local Euler section whereas
$\operatorname{Crit}(L)$ kills only the one-stage Euler section; there is a
canonical morphism
\[
\operatorname{Crit}^{\infty}(L)\longrightarrow\operatorname{Crit}(L),
\]
not an unconditional equivalence.

The global one-stage Euler class and its local adjoint are kept distinct.  The
global class lies in a dependent product over $Q$; the local class lies over
$\overline A$.  Their comparison is always through the constructed adjunction
$a^*\dashv\Pi_a$.

The horizontal stage $q$ and the symmetry-coherence depth $r$ have different
meanings.  The horizontal stage records a tail of the complete Cartan totalization
filtration.  Symmetry coherence records partial totalization in the cosimplicial
formal-symmetry direction.  Neither index is identified with Artinian Taylor weight
or one-marked tangent weight.

On the Lie side, retain separately
\[
I^{\mathrm{for}}_{\mathcal G},\qquad
i_{\mathcal G}^{\pi,E_\infty},\qquad
F_{\mathcal G}^{\pi,E_\infty},\qquad
\mathbb A_{\mathcal G}^{\pi,E_\infty}.
\]
The differentiated actual-isotropy controller and the homotopy anchor fibre are
related by
\[
\gamma_{\mathcal G}^{\pi,E_\infty}:
i_{\mathcal G}^{\pi,E_\infty}
\longrightarrow
F_{\mathcal G}^{\pi,E_\infty},
\]
and remain distinct unless a constructed inverse locus is inhabited.

The raw relation, quotient-formalized relation, directional Artinian cross-effects,
formal-moduli representation, and partition-Lie controller are likewise not
identified by notation.  The native symmetry-coherence obstruction tower is a
cosimplicial-totalization construction.  The differentiated Second Noether theorem
is instead obtained only after Artinian formalization, Schlessinger reflection, and
differentiation.

Finally, Hamiltonianity, nullification of $d_VJ$, and presymplectic radicality are
three different conditions.  Hamiltonianity plus a nullhomotopy of $d_VJ$ kills the
image of the presymplectic transgression in the horizontal tail; the residual loop
class must still be nullified by the separately constructed radical-realization
object.

\begin{convention}[Non-conflation]
\label{conv:symmetry-noether-non-conflation}
In this chapter:
\begin{enumerate}
\item ``zero symmetry action'' means factorization through the constructed homotopy
      anchor fibre;
\item ``actual infinitesimal isotropy'' means the independently differentiated
      actual formal isotropy;
\item ``native Noether obstruction'' means an obstruction in the raw symmetry
      totalization tower;
\item ``formalized Noether obstruction'' means the corresponding obstruction on the
      quotient-formalized Artinian branch;
\item ``differentiated Noether splitting'' means the partition-Lie algebroid
      splitting obtained by differentiating the reflected relative infinite-loop
      section;
\item ``higher differentiated Noether coherence'' means the morphism
      corolla--latching relations of that splitting, not the bare fibres of its
      directional Taylor coefficients;
\item ``presymplectic radical realization'' means a compatible nullification of the
      residual transgression left after Hamiltonianity and $d_VJ$-nullification;
\item ``represented presymplectic radical'' means the homotopy fibre of the separately
      constructed exact first-order $(1,1)$ contraction morphism; it is not the fibre
      of an unconstructed exterior-form model of the raw mixed coefficient;
\item ``complete presymplectic descent'' means a point of the symmetry-\v{C}ech
      totalization object $\operatorname{PreSymDesc}$ constructed below;
\item ``presymplectic basicness'' means its derived Cartan/null-contraction shadow,
      represented independently by $\operatorname{BasicReal}$.
\end{enumerate}
\end{convention}

\begin{convention}[Theorem domains]
\label{conv:symmetry-noether-theorem-domains}
Raw formal-groupoid, finite-limit, stable-normalization, right-Kan, realization,
and effective-descent constructions are formed in the native big spectral
Nisnevich $\infty$-topos.  Quasi-coherent cotangent and ordinary tangent formulas
are used only on represented theorem domains.  The global partition-Lie and
formalized Spencer branch is formed on the locally coherent spectral coefficient
sector and has the varying-object spectral-\'etale functoriality already
constructed there.  Rationality or explicit dg-model hypotheses occur only in the
corresponding rectification statements.

The critical linearization theorem below is specifically the representably smooth
scalar degree-zero sector
\[
M=\mathcal O_Q^{\mathrm{add}},\qquad n=0,
\]
with the Chapter 12 finite-locally-free affine-point cotangent refinement.  Ordinary
quasi-coherent tangent complexes are asserted only after pullback to represented
affine charts.  No analogous linear-space or tangent formula is asserted for a
general stable coefficient.

The represented first-order comparison with the partition-Lie controller additionally
uses the represented locally coherent effective-quotient domain of the differentiation
theorem; outside it the formal-leaf tangent coefficient is retained intrinsically.

The relative unstable Schlessinger reflector is formed on coherent spectral Artinian
formal leaves.  Its coherent one- and two-variable spectral-\'etale exchange is proved
below from the Chapter 14 Artinian transport; for an arbitrary spectral base change
only the canonical mate comparison is asserted.

The intrinsic anchor-splitting realization and its curvature/Bianchi hierarchy are
formed in the pro-coherent anchored category on the locally coherent differentiated
branch.  Identification with the ordinary represented splitting space uses the
represented locally coherent quotient comparison together with finite-locally-free
duality for the relevant cotangent coefficients.  Outside that overlap the intrinsic
pro-coherent and represented first-order splitting spaces are retained separately.

The represented presymplectic-radical construction below uses a separately
constructed exact first-order $(1,1)$-linearization of the complete mixed
transgression.  It does not identify the raw vertical degree-two cochain coefficient
with an exterior square of a cotangent complex, and it does not infer a bilinear
form from the unlinearized complete mixed coefficient.
\end{convention}

\begin{convention}[Realization instead of compatibility hypotheses]
\label{conv:symmetry-noether-realization-instead-of-compatibility-hypotheses}
Whenever an additional compatibility between already constructed objects is needed,
the default construction is an internal mapping object, a zero locus, an inverse
locus, or the homotopy fibre of the relevant forgetful or defect morphism.  Such an
object may be empty.  Its inhabitation is a property; a point contains the required
compatibility homotopy and its higher coherences.  No Cartan action on an arbitrary
coefficient, basicness datum, critical-formality datum, or reduction datum is
retained as input.
\end{convention}
\section{Parameterized realization objects}
\label{sec:symmetry-noether-parameterized-realization-objects}

Several compatibility problems below vary over a geometric parameter object and
are subsequently used as bases for further pullbacks.  We therefore internalize
them once and for all.

\begin{construction}[Parameterized mapping object]
\label{constr:symmetry-noether-parameterized-mapping-object}
Let $P$ be an object of the native spectral Nisnevich $\infty$-topos and suppose
\[
T\longmapsto\mathcal C(T)
\]
is a contravariant coefficient system of $\infty$-categories on objects over $P$
which satisfies effective descent.  Let $X,Y\in\mathcal C(P)$ be objects whose
pullbacks are defined on this coefficient system.  Define the presheaf
\[
\underline{\operatorname{Map}}_P(X,Y)(T\to P)
:=
\operatorname{Map}_{\mathcal C(T)}(X_T,Y_T).
\]
Whenever a subscript $P$ is omitted below, the relevant parameter base is the base
displayed immediately before the construction.
\end{construction}

\begin{lemma}[Descent of parameterized mapping objects]
\label{lem:symmetry-noether-descent-of-parameterized-mapping-objects}
The preceding presheaf is a spectral Nisnevich sheaf and therefore defines an
object
\[
\underline{\operatorname{Map}}_P(X,Y)\longrightarrow P.
\]
The same is true after taking finite or complete limits, equivalence loci, and
homotopy fibres over specified parameterized sections.
\end{lemma}

\begin{proof}
Let $T_\bullet\to T$ be an effective descent diagram.  By hypothesis
\[
\mathcal C(T)\simeq\lim_{[m]\in\Delta}\mathcal C(T_m).
\]
Mapping spaces in a limit of $\infty$-categories are the limits of the corresponding
mapping spaces, so
\[
\operatorname{Map}_{\mathcal C(T)}(X_T,Y_T)
\simeq
\lim_m
\operatorname{Map}_{\mathcal C(T_m)}(X_{T_m},Y_{T_m}).
\]
Thus the mapping presheaf satisfies descent.  Limits and homotopy fibres of sheaves
are again sheaves.
\end{proof}



\begin{construction}[Stable restriction along a relation morphism]
\label{constr:symmetry-noether-stable-restriction-along-a-relation-morphism}
Let
\[
\lambda:R\longrightarrow R'
\]
be a morphism over a parameter base $P$, with structural maps
\[
r:R\longrightarrow P,
\qquad
r':R'\longrightarrow P,
\qquad
r=r'\lambda.
\]
Let $K\in E_P$.  Whenever the stable pullback/right-adjoint pair
\[
\lambda^*:E_{R'}\rightleftarrows E_R:\Pi_\lambda
\]
is defined, let
\[
\eta_\lambda:\operatorname{id}\longrightarrow\Pi_\lambda\lambda^*
\]
be its unit.  Define the stable restriction morphism
\[
\boxed{
\operatorname{res}_{\lambda,K}:
\Pi_{r'}(r')^*K
\longrightarrow
\Pi_r r^*K
}
\]
by the composite
\[
\Pi_{r'}(r')^*K
\xrightarrow{\ \Pi_{r'}\eta_\lambda\ }
\Pi_{r'}\Pi_\lambda\lambda^*(r')^*K
\simeq
\Pi_r r^*K.
\]
Equivalently, $\operatorname{res}_{\lambda,K}$ is the right-adjoint mate of the
canonical identity
\[
\lambda^*(r')^*K\simeq r^*K.
\]
In the native stable coefficient system this construction is unconditional.  On a
formalized coherent coefficient stack it is formed chartwise wherever the same
stable pullback/right-adjoint pair has been constructed.
\end{construction}

\begin{lemma}[Composition, simplicial compatibility, and base change of stable restriction]
\label{lem:symmetry-noether-stable-restriction-composition-base-change}
Stable restriction has the following canonical coherences.
\begin{enumerate}
\item For composable relation morphisms
      \[
      R\xrightarrow{\lambda}R'\xrightarrow{\mu}R''
      \]
      over $P$,
      \[
      \operatorname{res}_{\mu\lambda,K}
      \simeq
      \operatorname{res}_{\lambda,K}
      \operatorname{res}_{\mu,K}.
      \]
\item If $\lambda_\bullet:R_\bullet\to R'_\bullet$ is simplicial, the degreewise
      restriction morphisms assemble to a morphism of cosimplicial stable
      coefficients.
\item In every Cartesian square for which stable right Beck--Chevalley is available,
      pullback carries $\operatorname{res}_{\lambda,K}$ to the corresponding
      restriction morphism after base change.
\end{enumerate}
All three assertions include the induced higher coherence data.
\end{lemma}

\begin{proof}
For composition, both displayed composites are right-adjoint mates of the same
canonical identification
\[
(\mu\lambda)^*(r'')^*K
\simeq
\lambda^*\mu^*(r'')^*K
\simeq
r^*K.
\]
The contractible uniqueness of mates gives the comparison and all associativity
coherence.  If $\lambda_\bullet$ is simplicial, every face or degeneracy square is a
commutative square of relation maps, so naturality of the units and the same mate
calculus identify the two possible composites.  Finally, stable right
Beck--Chevalley transports the unit of $\lambda^*\dashv\Pi_\lambda$ to the unit of
the pulled-back adjunction; applying the structural dependent products gives the
base-change comparison.  The Beck--Chevalley pasting laws supply the higher
coherences.
\end{proof}

\begin{construction}[Parameterized equivalence locus]
\label{constr:symmetry-noether-parameterized-equivalence-locus}
Let $P$ and $\mathcal C(-)$ be as above and let
\[
f:X\longrightarrow Y
\]
be a morphism in $\mathcal C(P)$.  Let $\mathbb E$ denote the nerve of the
walking-isomorphism groupoid with objects $0,1$ and mutually inverse arrows
\[
0
\mathrel{\overset{u}{\underset{u^{-1}}{\rightleftarrows}}}
1,
\]
and let
\[
j:\Delta^1\longrightarrow\mathbb E
\]
select the arrow $u$.  For a test $T\to P$, define the endpoint-fixed walking-equivalence
space
\[
\operatorname{Eq}_{\mathcal C(T)}(X_T,Y_T)
:=
\operatorname{hofib}_{(X_T,Y_T)}
\left(
\operatorname{Fun}(\mathbb E,\mathcal C(T))^\simeq
\longrightarrow
\mathcal C(T)^\simeq\times\mathcal C(T)^\simeq
\right).
\]
Restriction along $j$ gives
\[
j^*:
\operatorname{Eq}_{\mathcal C(T)}(X_T,Y_T)
\longrightarrow
\operatorname{Map}_{\mathcal C(T)}(X_T,Y_T).
\]
Define the parameterized equivalence locus by
\[
\boxed{
\operatorname{Inv}_P(f)(T\to P)
:=
\operatorname{hofib}_{f_T}(j^*).
}
\]
Thus a point is an extension of the already constructed arrow $f_T$ to the
walking equivalence.  Equivalently it consists of inverse data together with the
unit, counit, triangle, and all higher coherences forced by the groupoid
$\mathbb E$.  No inverse is retained independently as input.
\end{construction}

\begin{lemma}[Recognition and descent of the equivalence locus]
\label{lem:symmetry-noether-recognition-and-descent-of-the-equivalence-locus}
For every test $T\to P$,
\[
\operatorname{Inv}_P(f)_T\simeq
\begin{cases}
\varnothing,& f_T\text{ is not an equivalence},\\
*,& f_T\text{ is an equivalence}.
\end{cases}
\]
Consequently $\operatorname{Inv}_P(f)\to P$ is already the derived
$(-1)$-truncated equivalence locus of $f$; no separate image construction is
required.  The construction satisfies effective descent and arbitrary base change
whenever the coefficient system does.
\end{lemma}

\begin{proof}
The inclusion
\[
j:\Delta^1\longrightarrow\mathbb E
\]
exhibits $\mathbb E$ as the localization of $\Delta^1$ obtained by inverting its
unique nonidentity edge.  Hence restriction
\[
\operatorname{Fun}(\mathbb E,\mathcal C(T))
\longrightarrow
\operatorname{Fun}(\Delta^1,\mathcal C(T))
\]
is fully faithful with essential image the equivalence arrows.  After fixing the
endpoints, the fibre over $f_T$ is therefore empty if $f_T$ is not an equivalence
and contractible if it is.

For descent, $\operatorname{Fun}(\mathbb E,-)$ preserves limits of
$\infty$-categories, maximal subgroupoids preserve limits, and the endpoint and
arrow fibres are finite limits of spaces.  Thus a descent equivalence
\[
\mathcal C(T)\simeq\lim_m\mathcal C(T_m)
\]
induces the corresponding equivalence of inverse-locus spaces.  The same
limit argument gives base change.
\end{proof}

\begin{construction}[Parameterized zero locus of a stable section]
\label{constr:symmetry-noether-parameterized-zero-locus-of-a-stable-section}
Let $K\in E_P$ and let
\[
\sigma:\mathbf1_P\longrightarrow K
\]
be a stable section.  Under the parameterized sphere/relative-zeroth-space
adjunction it determines a section
\[
\bar\sigma:P\longrightarrow\Omega_P^\infty K.
\]
Let
\[
z_K:P\longrightarrow\Omega_P^\infty K
\]
be the zero section.  Define the geometric nullification object
\[
\boxed{
Z_P(\sigma)
:=
P\times_{\Omega_P^\infty K}P,
}
\]
where the two maps are $\bar\sigma$ and $z_K$.  Conversely, when a
section
\[
s:P\longrightarrow\Omega_P^\infty K
\]
is given directly, write $Z_P(s)$ for the same zero locus formed from the stable
section adjoint to $s$.
\end{construction}

\begin{lemma}[Recognition of stable nullification]
\label{lem:symmetry-noether-recognition-of-stable-nullification}
For every test $T\to P$, a $T$-point of $Z_P(\sigma)$ is precisely a specified
coherent nullhomotopy
\[
\sigma_T\simeq0
\]
in
$\operatorname{Map}_{E_T}(\mathbf1_T,K_T)$.  The construction commutes with
parameterized pullback and descent.  The same formula applies in any stable
coefficient system satisfying descent, including the stable Cartan coefficient
systems and the coherent spectral-\'etale coefficient stacks used below.
\end{lemma}

\begin{proof}
Mapping $T$ into the displayed homotopy pullback gives the homotopy fibre of
\[
\operatorname{Map}_{E_T}(\mathbf1_T,K_T)
\]
over the pair of points represented by $\sigma_T$ and $0$.  This is the path
space from $\sigma_T$ to zero.  Pullback compatibility follows from pullback
compatibility of the relative zeroth-space functor and finite limits; descent is
the same limit argument.
\end{proof}

\begin{construction}[Compatible nullification over a prescribed image homotopy]
\label{constr:symmetry-noether-compatible-nullification-over-a-prescribed-image-homotopy}
Let
\[
u:K\longrightarrow L
\]
be a morphism in $E_P$, let $\sigma:\mathbf1_P\to K$, and suppose a specified
nullhomotopy of the image has already been constructed, equivalently a section
\[
\eta:P\longrightarrow Z_P(u\sigma).
\]
Functoriality gives
\[
u_*:Z_P(\sigma)\longrightarrow Z_P(u\sigma).
\]
Define
\[
\boxed{
Z_P(\sigma;u,\eta)
:=
P\times_{Z_P(u\sigma)}Z_P(\sigma),
}
\]
where $P\to Z_P(u\sigma)$ is $\eta$.
A point is a nullhomotopy of $\sigma$ whose image under $u$ is the already
selected nullhomotopy $\eta$, including the required higher compatibility.
\end{construction}

\begin{construction}[Coherent spectral-\'etale mapping and equivalence sheaves]
\label{constr:symmetry-noether-coherent-spectral-etale-mapping-and-equivalence-sheaves}
On the locally coherent formalized coefficient sector, let
$\mathcal C_{\mathrm{\acute et}}(-)$ be one of the spectral formal-moduli,
partition-Lie algebroid, tangent-representation, or pro-coherent coefficient
stacks of Chapters 14--15.  For objects $X,Y$ over a coherent affine-\'etale chart,
define
\[
\underline{\operatorname{Map}}_{\mathrm{\acute et}}(X,Y)(U/V)
:=
\operatorname{Map}_{\mathcal C_{\mathrm{\acute et}}(U/V)}(X_{U/V},Y_{U/V}).
\]
The coherent spectral-\'etale descent equivalences make this an \'etale sheaf of
spaces.  Finite limits, homotopy fibres, and the finite-weight mapping towers below
are formed in this \'etale sheaf when the coefficient theory is not a native
Nisnevich coefficient system.

For a morphism $f$ in such a coefficient stack, define
\[
\operatorname{Inv}_{\mathrm{\acute et}}(f)
\]
chartwise by the same walking-equivalence construction as above: fix the two
endpoints in $\operatorname{Fun}(\mathbb E,-)^\simeq$ and take the fibre of
restriction along $j:\Delta^1\to\mathbb E$ over $f$.  Its chartwise fibres are
empty or contractible according as $f$ is not or is an equivalence, and the
coherent spectral-\'etale transition maps make these loci descend.
\end{construction}

\begin{proof}
The formalized coefficient categories are limits of their coherent affine-\'etale
restrictions.  Mapping spaces and functor categories out of the fixed small
$\infty$-groupoid $\mathbb E$ commute with these limits, as do maximal subgroupoids
and homotopy fibres.  The recognition statement is therefore the chartwise
walking-equivalence recognition lemma, and the remaining assertions follow because
limits of sheaves of spaces are computed objectwise.
\end{proof}

\begin{convention}[Pullback to a spectral-\'etale realization locus]
\label{conv:symmetry-noether-pullback-to-spectral-etale-realization-locus}
Let $\mathscr X$ be an \'etale sheaf of spaces on the coherent affine-\'etale
coefficient site.  When an object or morphism of a coherent coefficient stack is
said to be pulled back ``over $\mathscr X$'', this means pullback of the coefficient
stack to the category of elements
\[
\int \mathscr X.
\]
Thus an object of $\int\mathscr X$ is a coherent affine-\'etale chart together with
a point of $\mathscr X$ on that chart, and transition maps carry the selected point
by sheaf restriction.  In particular, if
\[
\mathscr X=\operatorname{Inv}_{\mathrm{\acute et}}(f),
\]
then the tautological point on $\int\mathscr X$ is a walking-equivalence extension
of the pulled-back arrow $f$.  It therefore supplies its inverse, unit, counit,
triangle, and all higher coherences, compatibly under coherent spectral-\'etale
transition.  This convention does not assert that $\mathscr X$ is represented by a
native geometric object.
\end{convention}

\begin{convention}[Internal and external realization notation]
\label{conv:symmetry-noether-internal-and-external-realization-notation}
Every native realization object below which is later used in a fibre product over
$Q$, $\overline A$, a critical parameter object, or another native parameter base
is understood as the internal object supplied by the preceding descent lemma.  Its
global space of choices is obtained only afterwards by applying global sections.
On the formalized pro-coherent/partition-Lie branch, the corresponding notation
denotes an object of the coherent spectral-\'etale mapping sheaf constructed above.
In particular, symbols such as
\[
\operatorname{CartLift},\qquad
\operatorname{BasicReal},\qquad
\operatorname{CartMorReal}
\]
denote parameterized objects whenever a geometric base is present.
\end{convention}
\section{Formal field symmetry and the ambient infinitesimal relation}
\label{sec:symmetry-noether-formal-field-symmetry-and-the-ambient-infinitesimal-relation}

Apply the relative de Rham effective-image construction to
\[
a:\overline A\longrightarrow Q.
\]
Write
\[
\epsilon_A:
\overline A\twoheadrightarrow G_A:=Q_{\overline A/Q}
\]
for the effective relative infinitesimal quotient and
\[
R^A_\bullet:=R_{\overline A/Q,\bullet}
=
\check C_\bullet(\epsilon_A).
\]
Thus
\[
R^A_1
\simeq
\overline A\times_{G_A}\overline A
\simeq
\overline A\times_{\overline A_{\mathrm{dR}/Q}}\overline A.
\]
The ordinary pair relation
$\overline A\times_Q\overline A$ is not identified with $R^A_1$.

\begin{definition}[Formal field-symmetry groupoid]
\label{def:symmetry-noether-formal-field-symmetry-groupoid}
A formal field-symmetry groupoid of the descended Cartan field geometry is an
object
\[
\mathcal G_\bullet\in
\operatorname{FormLieGpd}(\overline A/Q).
\]
\end{definition}

The formal-groupoid chapter constructs a contractibly unique simplicial anchor
\[
\phi_{\mathcal G}:
\mathcal G_\bullet\longrightarrow R^A_\bullet.
\]
No anchor is part of the definition of $\mathcal G_\bullet$.

\begin{definition}[Formal symmetry quotient]
\label{def:symmetry-noether-formal-symmetry-quotient}
Let
\[
q_{\mathcal G}:
\overline A\twoheadrightarrow Q_{\mathcal G}
:=
|\mathcal G_\bullet|
\]
be the effective quotient.  Realization of the canonical anchor gives a canonical
morphism
\[
r_{\mathcal G}:
Q_{\mathcal G}\longrightarrow G_A
\]
with
\[
r_{\mathcal G}q_{\mathcal G}\simeq\epsilon_A.
\]
Equivalently, the canonical quotient factorization is displayed by
\begin{tikzcd}[column sep=huge,row sep=large]
\overline A
  \arrow[r,two heads,"q_{\mathcal G}"]
  \arrow[dr,two heads,"\epsilon_A"'] &
Q_{\mathcal G}
  \arrow[d,two heads,"r_{\mathcal G}"] \\
& G_A.
\end{tikzcd}
\end{definition}

\begin{proposition}[Quotient and relation presentation]
\label{prop:symmetry-noether-quotient-and-relation-presentation}
For every $p\ge0$ there is a canonical equivalence
\[
\mathcal G_p
\simeq
\underbrace{
\overline A\times_{Q_{\mathcal G}}\cdots\times_{Q_{\mathcal G}}\overline A
}_{p+1\ \mathrm{factors}},
\]
and the anchor in degree $p$ is induced by postcomposition of the quotient
factorization
\[
\overline A
\xrightarrow{q_{\mathcal G}}
Q_{\mathcal G}
\xrightarrow{r_{\mathcal G}}
G_A.
\]
All simplicial identities are the corresponding finite-limit identities.
\end{proposition}

\begin{proof}
Formal groupoids are effective groupoid objects.  Hence the groupoid is canonically
the \v{C}ech nerve of its realization map
$q_{\mathcal G}:\overline A\twoheadrightarrow Q_{\mathcal G}$.
The canonical anchor realizes to $r_{\mathcal G}$ and is the identity in degree
zero.  Substitution in the two \v{C}ech nerves gives the degreewise maps.  Since
both nerves are formed by repeated fibre products, deletion and repetition of
vertices give the complete simplicial coherence.
\end{proof}

\begin{remark}[Not a monomorphism]
\label{rem:symmetry-noether-not-a-monomorphism}
The canonical anchor is not required to be a monomorphism.  Its homotopy fibres
contain isotropy.  In particular, one-object formal groups and completed symmetry
groupoids would be excluded by imposing a monomorphism condition.
\end{remark}

\begin{construction}[Effective infinitesimal arrow image]
\label{constr:symmetry-noether-effective-infinitesimal-arrow-image}
Factor the degree-one anchor by effective image:
\[
\mathcal G_1
\xrightarrow{e_1}
\operatorname{Im}^{\mathrm{for}}_1(\mathcal G)
\xrightarrow{m_1}
R^A_1,
\]
with $e_1$ an effective epimorphism and $m_1$ a monomorphism.  Write
\[
I_1:=\operatorname{Im}^{\mathrm{for}}_1(\mathcal G).
\]
The source and target maps of $R^A_1$ restrict uniquely to maps
$s_I,t_I:I_1\rightrightarrows\overline A$, because their composites with $e_1$
are the source and target of $\mathcal G_1$.  The groupoid unit and inversion
likewise factor uniquely through $I_1$.
\end{construction}

\begin{theorem}[The effective arrow image is an infinitesimal subrelation]
\label{thm:symmetry-noether-effective-arrow-image-is-an-infinitesimal-subrelation}
There is a contractibly unique composition
\[
m_I:
I_1\times_{t_I,\overline A,s_I}I_1
\longrightarrow I_1
\]
for which $m_1m_I$ is the restriction of the ambient relation composition.  With
this composition, the restricted unit and inversion, $I_1\rightrightarrows\overline A$
is an internal groupoid and
\[
\operatorname{Im}^{\mathrm{for}}_\bullet(\mathcal G)
\hookrightarrow R^A_\bullet
\]
is a simplicial subrelation, where
\[
\operatorname{Im}^{\mathrm{for}}_p(\mathcal G)
:=
\underbrace{
I_1\times_{\overline A}\cdots\times_{\overline A}I_1
}_{p\ \mathrm{composable\ arrows}}
\qquad (p\ge1),
\qquad
\operatorname{Im}^{\mathrm{for}}_0(\mathcal G):=\overline A.
\]
For every $p\ge0$, the degree-$p$ anchor has the canonical effective-image
factorization
\[
\mathcal G_p\twoheadrightarrow
\operatorname{Im}^{\mathrm{for}}_p(\mathcal G)
\hookrightarrow R^A_p.
\]
The passage to the effective image still forgets isotropy and therefore does not
replace the formal groupoid itself.
\end{theorem}

\begin{proof}
Define $s_I:=s_Rm_1$ and $t_I:=t_Rm_1$; compatibility of the anchor gives
$s_Ie_1\simeq s_{\mathcal G}$ and $t_Ie_1\simeq t_{\mathcal G}$.  The unit
factors through $I_1$ because $\phi_{\mathcal G,1}u_{\mathcal G}=u_R$ and the
image of the left side lies in $I_1$.  For inversion, the composite
$\iota_Rm_1:I_1\to R^A_1$ factors through $I_1$ after the effective cover $e_1$,
because
\[
\iota_R\phi_{\mathcal G,1}
\simeq
\phi_{\mathcal G,1}\iota_{\mathcal G}.
\]
Locality of factorization through the monomorphism $m_1$ therefore constructs the
restricted inversion through a contractible space.  Effective epimorphisms are
stable under pullback and finite products in an
$\infty$-topos, so
\[
e_1\times e_1:
\mathcal G_1\times_{\overline A}\mathcal G_1
\twoheadrightarrow
I_1\times_{\overline A}I_1
\]
is an effective epimorphism.  On this effective cover, compatibility of the
canonical anchor with groupoid composition gives
\[
m_R(m_1e_1\times m_1e_1)
\simeq
m_1e_1m_{\mathcal G}.
\]
Factorization through the monomorphism $m_1$ is local for an effective
epimorphism, so this identity descends through a contractible space to the required
$m_I$.  The construction is summarized by the commutative diagram
\begin{tikzcd}[column sep=large,row sep=large]
\mathcal G_1\times_{\overline A}\mathcal G_1
  \arrow[r,two heads,"e_1\times e_1"]
  \arrow[d,"m_{\mathcal G}"'] &
I_1\times_{\overline A}I_1
  \arrow[r,hook,"m_1\times m_1"]
  \arrow[d,dashed,"m_I"] &
R^A_1\times_{\overline A}R^A_1
  \arrow[d,"m_R"] \\
\mathcal G_1
  \arrow[r,two heads,"e_1"'] &
I_1
  \arrow[r,hook,"m_1"'] &
R^A_1.
\end{tikzcd}
Associativity, the two unit identities, and the inverse identities may be checked
after pullback along the corresponding effective covers from $\mathcal G$; they
then follow from the groupoid identities there.  Equivalently, after composition
with the monomorphism into $R^A_\bullet$ they are the ambient relation identities,
and monomorphisms detect the required equalities.  Iterated composable-arrow fibre
products therefore give a simplicial internal subgroupoid of the ambient relation.  The degreewise map
from $\mathcal G_p$ is an effective epimorphism by pullback stability, and the map
to $R^A_p$ is a monomorphism by finite-limit stability of monomorphisms.  Hence it
is the effective-image factorization in every degree.
\end{proof}

\begin{remark}[Later gauge-theoretic applications]
\label{rem:symmetry-noether-later-gauge-theoretic-applications}
Action groupoids and formal gauge groupoids constructed in the preceding chapters
are examples of formal field-symmetry groupoids.  No specialization to those
examples is made here.  A later theory of brave-new gauge fields may supply such a
groupoid together with additional geometric structure; the present chapter uses
only the formal-groupoid symmetry itself.
\end{remark}
\section{The symmetry--horizontal double relation}
\label{sec:symmetry-noether-the-symmetry-horizontal-double-relation}

The variational chapter uses the horizontal Cartan relation
\[
H_q
:=
\underbrace{
X\times_Q\cdots\times_Q X
}_{q+1\ \mathrm{factors}}.
\]
We now replace only the vertical maximal infinitesimal relation by the chosen
formal symmetry relation.

\begin{definition}[Symmetry--horizontal double relation]
\label{def:symmetry-noether-symmetry-horizontal-double-relation}
For $p,q\ge0$ define
\[
R^{\mathcal G}_{p,q}(A)
:=
H_q\times_Q\mathcal G_p.
\]
Write
\[
r^{\mathcal G}_{p,q}:
R^{\mathcal G}_{p,q}(A)\longrightarrow Q
\]
for the structural morphism.
\end{definition}

\begin{proposition}[Bisimplicial symmetry--horizontal geometry]
\label{prop:symmetry-noether-bisimplicial-symmetry-horizontal-geometry}
The assignment
\[
([p],[q])
\longmapsto
R^{\mathcal G}_{p,q}(A)
\]
extends canonically to a bisimplicial object
\[
R^{\mathcal G}_{\bullet,\bullet}(A):
\Delta^{\mathrm{op}}\times\Delta^{\mathrm{op}}
\longrightarrow
(\mathsf{XSp})/Q.
\]
The formal anchor induces a morphism of bisimplicial objects
\[
\Phi_{\mathcal G}:
R^{\mathcal G}_{\bullet,\bullet}(A)
\longrightarrow
R^{\mathrm{var}}_{\bullet,\bullet}(A).
\]
\end{proposition}

\begin{proof}
The symmetry simplicial operators act only on the $\mathcal G_p$ factor and the
horizontal operators act only on the $H_q$ factor.  Both are finite-limit
constructions over the common base $Q$, so the two actions commute with their full
pasting coherence.  The maps
$\phi_{\mathcal G,p}:\mathcal G_p\to R^A_p$ commute with every simplicial operator;
taking their product with $H_q$ gives $\Phi_{\mathcal G}$.
\end{proof}

\section{Symmetry variational cochains and the complete bifiltration}
\label{sec:symmetry-noether-symmetry-variational-cochains-and-the-complete-bifiltration}

\begin{definition}[Symmetry variational cochains]
\label{def:symmetry-noether-symmetry-variational-cochains}
For $M\in E_Q$ define
\[
C_{\mathcal G}^{p,q}(A;M)
:=
\Pi_{r^{\mathcal G}_{p,q}}
(r^{\mathcal G}_{p,q})^*M
\in E_Q.
\]
The bisimplicial relation makes
\[
C_{\mathcal G}^{\bullet,\bullet}(A;M)
:
\Delta\times\Delta\longrightarrow E_Q
\]
bicosimplicial.
\end{definition}

Contravariance of stable coefficient restriction along
$\Phi_{\mathcal G}$ gives a canonical morphism
\[
\operatorname{res}_{\mathcal G}:
C_{\mathrm{var}}^{\bullet,\bullet}(A;M)
\longrightarrow
C_{\mathcal G}^{\bullet,\bullet}(A;M).
\]

\begin{definition}[Symmetry and horizontal normalization]
\label{def:symmetry-noether-symmetry-and-horizontal-normalization}
Normalize first in the symmetry direction while retaining the horizontal cosimplicial
direction:
\[
B_{\mathcal G}^{p,\bullet}(A;M)
:=
N_{\mathcal G}^p
C_{\mathcal G}^{\bullet,\bullet}(A;M).
\]
Then put
\[
\Omega_{\mathcal G}^{p,q}(A;M)
:=
N_H^q
B_{\mathcal G}^{p,\bullet}(A;M).
\]
\end{definition}

\begin{definition}[Symmetry secondary coefficients]
\label{def:symmetry-noether-symmetry-secondary-coefficients}
For every $p\ge0$ define
\[
\operatorname{Sec}_{\mathcal G}^{p}(A;M)
:=
\operatorname{Tot}_H
B_{\mathcal G}^{p,\bullet}(A;M).
\]
For $q\ge1$ let
\[
F_H^q
\operatorname{Sec}_{\mathcal G}^{p}(A;M)
:=
\operatorname{fib}
\left(
\operatorname{Sec}_{\mathcal G}^{p}(A;M)
\longrightarrow
\operatorname{Tot}_{H,\le q-1}
B_{\mathcal G}^{p,\bullet}(A;M)
\right).
\]
\end{definition}

For a cosimplicial object $C^\bullet$ write
\[
D^m(C):=\operatorname{Tot}_{\Delta_{\le m}}C^\bullet,
\qquad
D^\infty(C):=\operatorname{Tot}_{\Delta}C^\bullet,
\]
and set $D^{-1}(C)=0$.

\begin{definition}[Complete symmetry--horizontal bifiltration]
\label{def:symmetry-noether-complete-symmetry-horizontal-bifiltration}
Form first the complete symmetry totalization filtration in the stable category of
horizontal cosimplicial objects,
\[
F^p_{\mathcal G}C_{\mathcal G}
:=
\operatorname{fib}
\left(
\operatorname{Tot}_{\mathcal G}C_{\mathcal G}
\longrightarrow
\operatorname{Tot}_{\mathcal G,\le p-1}C_{\mathcal G}
\right).
\]
For $p,q\ge0$ define
\[
F^{p,q}_{\mathcal G,H}(A;M)
:=
\operatorname{fib}
\left(
\operatorname{Tot}_H F^p_{\mathcal G}C_{\mathcal G}
\longrightarrow
\operatorname{Tot}_{H,\le q-1}F^p_{\mathcal G}C_{\mathcal G}
\right).
\]
Write
\[
\operatorname{Var}_{\mathcal G}(A;M)
:=
F^{\bullet,\bullet}_{\mathcal G,H}(A;M).
\]
\end{definition}

Let
\[
\partial_{\mathcal G}:
\Omega_{\mathcal G}^{p,q}
\longrightarrow
\Omega_{\mathcal G}^{p+1,q},
\qquad
\partial_H:
\Omega_{\mathcal G}^{p,q}
\longrightarrow
\Omega_{\mathcal G}^{p,q+1}
\]
denote the unsigned coherent boundaries supplied by stable Dold--Kan.  Define the
signed unary shadow by
\[
\boxed{
 d_{\mathcal G}:=\partial_{\mathcal G},
 \qquad
 d_H|_{\Omega_{\mathcal G}^{p,q}}
 :=(-1)^p\partial_H.
}
\]
The underlying primary datum remains the product-indexed coherent double cochain;
the signs only rephase its unary differential shadow.

\begin{theorem}[Complete symmetry--horizontal coherent calculus]
\label{thm:symmetry-noether-complete-symmetry-horizontal-coherent-calculus}
The preceding bifiltration is complete in each variable.  Its iterated associated
graded is
\[
\operatorname{gr}_{\mathcal G}^p
\operatorname{gr}_H^q
\operatorname{Var}_{\mathcal G}(A;M)
\simeq
\Omega_{\mathcal G}^{p,q}(A;M)[-p-q].
\]
The reverse order of associated graded is canonically equivalent.  The associated
coherent double cochain object has structure maps
\[
d_{\mathcal G}:
\Omega_{\mathcal G}^{p,q}
\longrightarrow
\Omega_{\mathcal G}^{p+1,q},
\qquad
d_H:
\Omega_{\mathcal G}^{p,q}
\longrightarrow
\Omega_{\mathcal G}^{p,q+1}
\]
with specified coherent relations
\[
d_{\mathcal G}^2\simeq0,
\qquad
d_H^2\simeq0,
\qquad
d_Hd_{\mathcal G}+d_{\mathcal G}d_H\simeq0,
\]
including all higher pasting data.

The restriction morphism from the ambient variational double relation induces a
natural morphism of complete bifiltered objects
\[
\operatorname{Var}_{X/S}(A;M)
\longrightarrow
\operatorname{Var}_{\mathcal G}(A;M).
\]
\end{theorem}

\begin{proof}
The proof is the complete-filtration/coherent-double-cochain argument of the
variational chapter applied to
$C_{\mathcal G}^{\bullet,\bullet}(A;M)$.  Symmetry and horizontal normalization are
finite total fibres of commuting codegeneracy cubes, so they commute.  The two
complete filtrations are formed by limits in stable functor categories; Fubini for
limits identifies the two orders.  Stable Dold--Kan identifies the successive
partial-totalization fibres with normalized cochains shifted by their
cosimplicial degrees.  The signed mixed relation is the same Koszul rephasing of
the product-indexed coherent double cochain used previously.  Functoriality of all
these operations carries $\operatorname{res}_{\mathcal G}$ to the displayed
bifiltered morphism.
\end{proof}

\begin{construction}[Cartan-lift realization for an arbitrary coefficient]
\label{constr:symmetry-noether-cartan-lift-realization-for-an-arbitrary-coefficient}
The symmetry--horizontal construction above needs no Cartan structure on an arbitrary
coefficient because $M$ is pulled back from $Q$.  For later use, however, the
existence problem for a Cartan structure on an arbitrary
$E\in E_{\overline A}$ is represented over the constructed effective quotient
$Q_{\mathcal G}$.

Let
\[
q_{\mathcal G}:\overline A\twoheadrightarrow Q_{\mathcal G}
\]
and let $J_{\mathcal G}=q_{\mathcal G}^*\Pi_{q_{\mathcal G}}$ be its jet
comonad.  For every test $t:T\to Q_{\mathcal G}$ form
\[
\overline A_T:=T\times_{Q_{\mathcal G}}\overline A,
\qquad
q_{\mathcal G,T}:\overline A_T\longrightarrow T,
\]
and put
\[
J_{\mathcal G,T}:=q_{\mathcal G,T}^*\Pi_{q_{\mathcal G,T}},
\qquad
E_T:=E|_{\overline A_T}.
\]
Define
\[
\boxed{
\operatorname{CartLift}_{\mathcal G}(E)(T\to Q_{\mathcal G})
:=
\operatorname{hofib}_{E_T}
\left(
\bigl(\operatorname{Coalg}_{J_{\mathcal G,T}}
(E_{\overline A_T})\bigr)^\simeq
\longrightarrow
(E_{\overline A_T})^\simeq
\right).
}
\]
Thus a point is precisely a coherent $J_{\mathcal G,T}$-coalgebra structure on
the already constructed coefficient $E_T$.
\end{construction}

\begin{proposition}[Descent and base change of the Cartan-lift realization]
\label{prop:symmetry-noether-descent-and-base-change-of-the-cartan-lift-realization}
The preceding presheaf is a spectral Nisnevich sheaf on
$(\mathsf{XSp})/Q_{\mathcal G}$ and therefore defines an internal object
\[
\boxed{
\operatorname{CartLift}_{\mathcal G}(E)
\longrightarrow Q_{\mathcal G}.
}
\]
It is compatible with arbitrary raw spectral base change.  Its global sections
are the space of Cartan coactions on $E$.  Pullback along
$q_{\mathcal G}:\overline A\to Q_{\mathcal G}$ gives the corresponding family
over the field object when such a family is useful.
\end{proposition}

\begin{proof}
Let $T_\bullet\to T$ be an effective descent diagram over $Q_{\mathcal G}$.
Base change of the effective quotient gives compatible Cartesian diagrams
$\overline A_{T_m}\to T_m$, and stable Beck--Chevalley identifies the pulled-back
jet comonads.  Effective descent gives
\[
E_{\overline A_T}
\simeq
\lim_m E_{\overline A_{T_m}}
\]
compatibly with those comonads.  The Eilenberg--Moore construction for a compatible
comonad diagram is computed in the same limit of coefficient categories.  Taking
maximal subgroupoids and then the homotopy fibre over the compatible coefficient
$E_T$ therefore gives
\[
\operatorname{CartLift}_{\mathcal G}(E)(T)
\simeq
\lim_m
\operatorname{CartLift}_{\mathcal G}(E)(T_m).
\]
The base-change statement is the same argument for a Cartesian square of quotient
maps.  For $T=Q_{\mathcal G}$ the fibre is exactly the space of
$J_{\mathcal G}$-coalgebra structures on $E$.
\end{proof}

\begin{construction}[Cofree Cartan comparison]
\label{constr:symmetry-noether-cofree-cartan-comparison}
Independently of whether
$\operatorname{CartLift}_{\mathcal G}(E)$ is inhabited, the cofree cobar
\[
\operatorname{Cobar}_{\mathcal G}^{\mathrm{cofr},\bullet}(E)
\]
is always defined.  Over $\operatorname{CartLift}_{\mathcal G}(E)$ the universal
Cartan coaction induces the canonical comparison from the selected Cartan cobar to
the cofree cobar.  Thus an actual Cartan coaction is neither assumed nor replaced
by the cofree resolution.
\end{construction}

\section{Structured symmetry one-variations}
\label{sec:symmetry-noether-structured-symmetry-one-variations}

Let
\[
v=(T,N,c,\widetilde c)
\in
\mathcal V_{\overline A/Q}^{(1)}
\]
be a split structured one-variation.  Thus
$T=\operatorname{Spec}B$, $N$ is connective,
\[
\widetilde c:T[N]\longrightarrow\overline A
\]
lies over $Q$, and its restriction to the central square-zero section is $c$.

The variational chapter canonically constructs
\[
\sigma_v:
T[N]\longrightarrow R^A_1
\]
with vertices
\[
cr_N,\qquad \widetilde c.
\]

\begin{definition}[Space of symmetry lifts]
\label{def:symmetry-noether-space-of-symmetry-lifts}
Let
\[
\operatorname{Lift}_{\mathcal G}(v)
\]
be the homotopy fibre over $\sigma_v$ of
\[
\operatorname{Map}(T[N],\mathcal G_1)
\longrightarrow
\operatorname{Map}(T[N],R^A_1),
\]
taken in the pointed mapping problem which sends the central section
$T\hookrightarrow T[N]$ to the groupoid unit over $c$.
A symmetry-lifted one-variation is a pair
\[
(v,\widetilde\sigma_v),
\qquad
\widetilde\sigma_v\in
\operatorname{Lift}_{\mathcal G}(v).
\]
\end{definition}

\begin{definition}[Symmetry one-variation category]
\label{def:symmetry-noether-symmetry-one-variation-category}
A structured morphism
\[
f:v'\longrightarrow v
\]
of one-variations contains the induced morphism of split square-zero thickenings
from the domain of $v'$ to the domain of $v$.  Precomposition of a symmetry lift
with that morphism, together with functoriality of the canonical relation simplex,
gives a canonical map
\[
f^*:
\operatorname{Lift}_{\mathcal G}(v)
\longrightarrow
\operatorname{Lift}_{\mathcal G}(v').
\]
These maps are coherently compatible with identities and composition.  Thus
\[
\operatorname{Lift}_{\mathcal G}^{(1)}:
\left(\mathcal V_{\overline A/Q}^{(1)}\right)^{\mathrm{op}}
\longrightarrow
\mathcal S,
\qquad
v\longmapsto\operatorname{Lift}_{\mathcal G}(v),
\]
is a presheaf of spaces.  Define
\[
\mathcal V_{\mathcal G}^{(1)}
\longrightarrow
\mathcal V_{\overline A/Q}^{(1)}
\]
to be its right-fibration unstraightening, equivalently the Grothendieck
construction of this contravariant functor.  Write
\[
u_{\mathcal G}^{(1)}:
\mathcal V_{\mathcal G}^{(1)}
\longrightarrow
\mathcal V_{\overline A/Q}^{(1)}
\]
for the forgetful functor and
\[
\pi_{\mathcal G,1}:=
\pi_1u_{\mathcal G}^{(1)}
\]
for the affine-centre functor.
\end{definition}

\begin{construction}[Symmetry first-variation coefficient]
\label{constr:symmetry-noether-symmetry-first-variation-coefficient}
Define
\[
V_{\mathcal G,M}^{(1)}(A)
:=
\operatorname{Ran}_{\pi_{\mathcal G,1}^{\mathrm{op}}}
\left(
K_M^{(1)}
\circ
(u_{\mathcal G}^{(1)})^{\mathrm{op}}
\right).
\]
Restriction of cones from all structured variations to the symmetry-lifted subcategory
gives a canonical morphism
\[
\boxed{
\iota_{\mathcal G}^{(1)}:
V_M^{(1)}(A)
\longrightarrow
V_{\mathcal G,M}^{(1)}(A).
}
\]
\end{construction}

\begin{theorem}[Symmetry first-difference descent and raw base change]
\label{thm:symmetry-noether-symmetry-first-difference-descent-and-raw-base-change}
The raw right-Kan coefficient
$V_{\mathcal G,M}^{(1)}(A)$ is already a finite spectral Nisnevich sheaf for the
centre-pullback topology and commutes with arbitrary spectral base change.
The morphism $\iota_{\mathcal G}^{(1)}$ is natural for those base changes.
\end{theorem}

\begin{proof}
The lift functor is a finite homotopy pullback of mapping spaces against the
canonical formal anchor.  The anchor, the canonical relation simplex $\sigma_v$,
and structured split square-zero extensions commute with arbitrary base change.
Thus pullback of affine centres gives Cartesian lifts in
$\mathcal V_{\mathcal G}^{(1)}$.  The coefficient
$K_M^{(1)}$ already satisfies centre-pullback Nisnevich descent and stable
Beck--Chevalley.  The Cartesian-centre right-Kan descent lemma of the variational
chapter therefore applies to the symmetry-lifted indexing category.  The same
comma-category comparison after base change identifies the right-Kan limits
pointwise.  Finally, restriction of cones is natural under these identifications.
\end{proof}

\section{Finite cubical symmetry variations}
\label{sec:symmetry-noether-finite-cubical-symmetry-variations}

For a finite labelled set $I$, write
\[
B[N_I]
:=
\bigotimes_{B,\,i\in I}(B\oplus N_i),
\qquad
T[N_I]
:=
\operatorname{Spec}(B[N_I]),
\]
where every tensor product is relative to the central algebra $B$.  These are the
split cubical algebras, face immersions, and split retractions constructed in the
variational chapter.  For $I\subseteq J$ write
\[
j_{I,J}:T[N_I]\longrightarrow T[N_J],
\qquad
p_{J,I}:T[N_J]\longrightarrow T[N_I],
\qquad
p_{J,I}j_{I,J}=\mathrm{id}.
\]

Let
\[
v\in\mathcal V_{\overline A/Q}^{(k)}
\]
be a complete structured $k$-variation and write
\[
\widetilde a_I:T[N_I]\longrightarrow\overline A
\]
for its field on the $I$-face.

\begin{definition}[Nondegenerate face-chain category]
\label{def:symmetry-noether-nondegenerate-face-chain-category}
Let
\[
\operatorname{Ch}^{\mathrm{nd}}_*([k])
\]
be the finite category whose objects are nonempty strict chains
\[
\sigma=(I_0\subsetneq I_1\subsetneq\cdots\subsetneq I_r)
\]
in the power-set poset $\mathcal P([k])$ and whose morphisms are generated by
deleting members of a chain.  Thus an arrow is oriented from a nondegenerate
simplex to one of its faces.  No degeneracy morphism is included in this indexing
category.
\end{definition}

For
\[
\sigma=(I_0\subsetneq\cdots\subsetneq I_r)
\]
put
\[
T_\sigma:=T[N_{I_r}]
\]
and define
\[
f_j^\sigma
:=
\widetilde a_{I_j}p_{I_r,I_j}
:
T_\sigma\longrightarrow\overline A,
\qquad
0\le j\le r.
\]
Let
\[
i_\sigma:T\longrightarrow T_\sigma
\]
be the central section.

\begin{lemma}[Coherent de Rham cone for a face chain]
\label{lem:symmetry-noether-coherent-de-rham-cone-for-a-face-chain}
The morphism $i_\sigma$ is a finite composite of base changes of split square-zero
thickenings and is therefore inverted by relative de Rham localization over $Q$.
The tuple
\[
(f_0^\sigma,\ldots,f_r^\sigma)
\]
has a canonical coherent common image in $\overline A_{\mathrm{dR}/Q}$.
This common-image cone is natural in structured morphisms of variations, finite-set
relabeling, face restriction, and raw spectral base change.
\end{lemma}

\begin{proof}
Choose an ordering of the finite set $I_r$.  The central section factors as a
finite composite
\[
T\longrightarrow T[N_{i_1}]
\longrightarrow T[N_{i_1},N_{i_2}]
\longrightarrow\cdots\longrightarrow T[N_{I_r}],
\]
and every arrow is a base change of the central section of a split square-zero
extension.  Relative de Rham localization inverts each arrow and therefore
inverts $i_\sigma$.  The conclusion is independent of the chosen ordering because
the split cube carries the coherent symmetric-monoidal reindexing of the
variational chapter.

All $f_j^\sigma$ restrict along $i_\sigma$ to the same central field $c:T\to
\overline A$.  Write $L_{\mathrm{dR}/Q}$ for relative de Rham localization.  The
space of inverses of the equivalence $L_{\mathrm{dR}/Q}(i_\sigma)$ is
contractible.  Using its canonical inverse datum, put
\[
h_\sigma
:=
L_{\mathrm{dR}/Q}(c)\,
L_{\mathrm{dR}/Q}(i_\sigma)^{-1}.
\]
From
\[
L_{\mathrm{dR}/Q}(f_j^\sigma)
L_{\mathrm{dR}/Q}(i_\sigma)
\simeq
L_{\mathrm{dR}/Q}(c)
\]
we obtain coherent equivalences
\[
L_{\mathrm{dR}/Q}(f_j^\sigma)\simeq h_\sigma
\]
for all $j$, all induced from the same inverse datum.  This is the required
coherent cone.  Every construction used is functorial under the displayed
operations, proving the naturality statement.
\end{proof}

\begin{construction}[Canonical relation simplex of a face chain]
\label{constr:symmetry-noether-canonical-relation-simplex-of-a-face-chain}
Using
\[
R^A_r
\simeq
\underbrace{
\overline A\times_{\overline A_{\mathrm{dR}/Q}}\cdots
\times_{\overline A_{\mathrm{dR}/Q}}\overline A
}_{r+1\ \mathrm{factors}},
\]
the coherent cone of the preceding lemma gives a canonical factorization
\[
\boxed{
\sigma_{v,\sigma}:
T_\sigma\longrightarrow R^A_r
}
\]
whose ordered vertices are $f_0^\sigma,\ldots,f_r^\sigma$.  The factorization is
canonical through a contractible space and retains the full higher coherence of
the common relative de Rham image; it is not assembled from unrelated pairwise
homotopies.
\end{construction}

\begin{lemma}[Face compatibility of the relation simplices]
\label{lem:symmetry-noether-face-compatibility-of-the-relation-simplices}
Let
\[
\partial_j\sigma
=
(I_0<\cdots<\widehat I_j<\cdots<I_r).
\]
If $j<r$, then $T_{\partial_j\sigma}=T_\sigma$ and
\[
d_j\sigma_{v,\sigma}
=
\sigma_{v,\partial_j\sigma}.
\]
If $j=r$, then
\[
T_{\partial_r\sigma}=T[N_{I_{r-1}}],
\]
and
\[
\boxed{
d_r\sigma_{v,\sigma}\circ j_{I_{r-1},I_r}
=
\sigma_{v,\partial_r\sigma}.
}
\]
These identities are compatible with all iterated face maps, structured morphisms
of variations, finite-set relabelling, and raw spectral base change.
\end{lemma}

\begin{proof}
For $j<r$, deletion of the $j$th member deletes the $j$th vertex of the tuple
\[
(f_0^\sigma,\ldots,f_r^\sigma)
\]
without changing its parameter object.  This is exactly the $j$th face of the
\v{C}ech relation.

For $j=r$, restriction along
$j_{I_{r-1},I_r}:T[N_{I_{r-1}}]\to T[N_{I_r}]$ gives
\[
p_{I_r,I_i}j_{I_{r-1},I_r}
=
p_{I_{r-1},I_i}
\qquad
(0\le i<r),
\]
so deleting the terminal vertex after restriction gives precisely the tuple
defining $\sigma_{v,\partial_r\sigma}$.  The simplicial face identities follow
from the composition identities for the retractions $p_{J,I}$.  Naturality and
base change are termwise.
\end{proof}

\begin{construction}[Face functor of symmetry-lift spaces]
\label{constr:symmetry-noether-face-functor-of-symmetry-lift-spaces}
For
\[
\sigma=(I_0<\cdots<I_r)
\]
define
\[
\mathcal L_{\mathcal G}(v;\sigma)
:=
\operatorname{hofib}_{\sigma_{v,\sigma}}
\left(
\operatorname{Map}(T_\sigma,\mathcal G_r)
\longrightarrow
\operatorname{Map}(T_\sigma,R^A_r)
\right),
\]
with the pointed condition that the central section maps to the fully degenerate
unit simplex.

For the face $\partial_j\sigma$, define
\[
\mathcal L_{\mathcal G}(v;\sigma)
\longrightarrow
\mathcal L_{\mathcal G}(v;\partial_j\sigma)
\]
as follows:
\begin{enumerate}
\item if $j<r$, apply the simplicial face
      $d_j:\mathcal G_r\to\mathcal G_{r-1}$;
\item if $j=r$, first apply $d_r$ and then restrict the resulting map along
      $j_{I_{r-1},I_r}:T_{\partial_r\sigma}\to T_\sigma$.
\end{enumerate}
The preceding lemma shows that these maps land in the required homotopy fibres.
The simplicial identities of $\mathcal G_\bullet$ and the identities among the
$j_{I,J}$ make this a functor
\[
\mathcal L_{\mathcal G}(v;-):
\operatorname{Ch}^{\mathrm{nd}}_*([k])
\longrightarrow
\mathcal S.
\]
\end{construction}

\begin{definition}[Coherent symmetry lift of a cube]
\label{def:symmetry-noether-coherent-symmetry-lift-of-a-cube}
Define
\[
\boxed{
\operatorname{Lift}^{(k)}_{\mathcal G}(v)
:=
\lim_{\sigma\in\operatorname{Ch}^{\mathrm{nd}}_*([k])}
\mathcal L_{\mathcal G}(v;\sigma).
}
\]
The indexing category is finite.
\end{definition}

\begin{proposition}[Degenerate simplices carry no independent lift data]
\label{prop:symmetry-noether-degenerate-simplices-carry-no-independent-lift-data}
Every weak chain
\[
I_0\subseteq I_1\subseteq\cdots\subseteq I_r
\]
is uniquely a simplicial degeneracy of its nondegenerate reduction.  A compatible
point of $\operatorname{Lift}^{(k)}_{\mathcal G}(v)$ therefore determines the lift
on every degenerate simplex by applying the corresponding degeneracy of
$\mathcal G_\bullet$.  Conversely those degeneracy values are forced.  Thus the
nondegenerate face limit contains exactly the full simplicial coherence and no
extra degeneracy choices.
\end{proposition}

\begin{proof}
The nerve of a poset has the usual unique decomposition of every simplex as a
degeneracy of a nondegenerate simplex.  Applying the corresponding degeneracy maps
of $\mathcal G_\bullet$ defines the degenerate lifts.  The simplicial identities
show that all face restrictions agree with the already chosen nondegenerate data.
Uniqueness follows because a degeneracy value is prescribed by the simplicial
operator itself.
\end{proof}

\begin{proposition}[Recovery of the one-stage lift]
\label{prop:symmetry-noether-recovery-of-the-one-stage-lift}
For $k=1$ there is a canonical equivalence
\[
\operatorname{Lift}^{(1)}_{\mathcal G}(v)
\simeq
\operatorname{Lift}_{\mathcal G}(v).
\]
\end{proposition}

\begin{proof}
The only nontrivial positive-dimensional nondegenerate chain is
\[
\varnothing\subsetneq\{1\}.
\]
The zero-dimensional lift spaces are contractible because
$\mathcal G_0=R^A_0=\overline A$.  The remaining lifting problem is precisely the
degree-one homotopy fibre defining $\operatorname{Lift}_{\mathcal G}(v)$.
\end{proof}

\begin{definition}[Symmetry $k$-variation category]
\label{def:symmetry-noether-symmetry-k-variation-category}
The face-chain relation simplices and the face functor of lift spaces are natural
under structured morphisms of $k$-variations.  Precomposition therefore gives,
for every structured morphism $f:v'\to v$, a natural transformation of the
finite face-lift diagrams and hence a canonical map
\[
f^*:
\operatorname{Lift}^{(k)}_{\mathcal G}(v)
\longrightarrow
\operatorname{Lift}^{(k)}_{\mathcal G}(v').
\]
The induced maps are coherently compatible with identities and composition, so
\[
\operatorname{Lift}_{\mathcal G}^{(k)}:
\left(\mathcal V_{\overline A/Q}^{(k)}\right)^{\mathrm{op}}
\longrightarrow
\mathcal S
\]
is a presheaf.  Define
\[
\mathcal V_{\mathcal G}^{(k)}
\longrightarrow
\mathcal V_{\overline A/Q}^{(k)}
\]
to be its right-fibration unstraightening, equivalently the Grothendieck
construction of this contravariant functor.  Write
\[
u_{\mathcal G}^{(k)}:
\mathcal V_{\mathcal G}^{(k)}
\longrightarrow
\mathcal V_{\overline A/Q}^{(k)}
\]
for the forgetful functor and
\[
\pi_{\mathcal G,k}:=\pi_ku_{\mathcal G}^{(k)}.
\]
\end{definition}

\begin{construction}[Symmetry cubical coefficient]
\label{constr:symmetry-noether-symmetry-cubical-coefficient}
Define
\[
V_{\mathcal G,M}^{(k)}(A)
:=
\operatorname{Ran}_{\pi_{\mathcal G,k}^{\mathrm{op}}}
\left(
K_M^{(k)}
\circ
(u_{\mathcal G}^{(k)})^{\mathrm{op}}
\right).
\]
Restriction of right-Kan cones gives
\[
\iota_{\mathcal G}^{(k)}:
V_M^{(k)}(A)
\longrightarrow
V_{\mathcal G,M}^{(k)}(A).
\]
\end{construction}

\begin{theorem}[Finite-set symmetry, descent, and internal base change]
\label{thm:symmetry-noether-finite-set-symmetry-descent-and-internal-base-change}
Permutation of the labelled directions gives a coherent action
\[
\Sigma_k\curvearrowright V_{\mathcal G,M}^{(k)}(A),
\]
and $\iota_{\mathcal G}^{(k)}$ is $\Sigma_k$-equivariant.  The internal symmetry
cubical coefficient satisfies finite centre-pullback spectral Nisnevich descent and
arbitrary spectral base change.

For an arbitrary base change, the external category
$\mathcal V_{\mathcal G}^{(k)}$ carries the canonical restriction functor.  No
blanket equivalence of external variation categories is asserted.  At each affine
centre of the pulled-back field object, however, the comma categories and
coefficient diagrams which compute the two right-Kan values are canonically
equivalent; these pointwise equivalences are what produce the displayed internal
base-change equivalence.
\end{theorem}

\begin{proof}
A permutation transports the split cube, its nondegenerate face category, the
relation-simplex diagram, and the symmetry-lift diagram.  The variational total-fibre
coefficient is already natural under finite-set relabelling, so right Kan extension
gives the action.

For descent and base change, every lift condition is a finite homotopy limit of
mapping-space lifting problems against the base-change-compatible formal anchor.
Pullback of affine centres gives the corresponding restriction functor on external
variation categories.  The Cartesian-centre argument identifies, for each affine
centre after base change, the relevant comma category with the original
comma category at the composite centre.  Stable Beck--Chevalley identifies the
coefficient diagrams.  Taking their limits gives the internal right-Kan
base-change equivalence and the descent assertion.
\end{proof}

\section{The represented first-order symmetry-variation theorem}
\label{sec:symmetry-noether-the-represented-first-order-symmetry-variation-theorem}

The intrinsic symmetry-lift definition requires no representability.  We now enter
exactly the represented first-order theorem domain constructed for formal groupoids:
the relevant local representatives of $Q$, $\overline A$, and $\mathcal G_1$ are
spectral schemes, the source map
\[
s:\mathcal G_1\longrightarrow\overline A
\]
and the field map
\[
a:\overline A\longrightarrow Q
\]
are represented there, and the cotangent base-change formulas of the preceding
chapters apply.  These are theorem-domain properties of the already constructed
objects, not fields attached to $\mathcal G_\bullet$.

Let
\[
c:T=\operatorname{Spec}B\longrightarrow\overline A
\]
be an affine centre.  Put
\[
N_c
:=
\Gamma
\left(
T,
s_c^*L_{\overline A_c/T}
\right),
\qquad
\overline A_c:=T\times_Q\overline A.
\]
The variational chapter constructs the universal affine variation
\[
u_c\in\mathcal V_{\overline A/Q}^{(1)}
\]
and proves that every variation $v$ centred over $c$ is classified, through a
contractible space, by a derivation
\[
\delta_v:N_c\longrightarrow N.
\]

On the formal-groupoid side, let
\[
a^\vee_{\mathcal G,c}
:=
\Gamma\!\left(
T,
c^*u^*L_{\mathcal G_1/\overline A,s}
\right)
\in\operatorname{Mod}_B
\]
be the affine module representing the first source coefficient at $c$.  Equivalently,
before taking affine global sections the coefficient is
$c^*u^*L_{\mathcal G_1/\overline A,s}\in\operatorname{QCoh}(T)$; the displayed
module uses the canonical equivalence
$\operatorname{QCoh}(T)\simeq\operatorname{Mod}_B$.
The first-order linearization of the canonical nonlinear anchor constructs the
cotangent morphism
\[
\alpha^\vee_{\mathcal G,c}:
N_c\longrightarrow a^\vee_{\mathcal G,c}.
\]

\begin{theorem}[First-order symmetry-lift classification]
\label{thm:symmetry-noether-first-order-symmetry-lift-classification}
For every connective $B$-module $N$ and every split structured variation $v$
classified by
\[
\delta_v:N_c\longrightarrow N,
\]
there is a canonical natural equivalence
\[
\boxed{
\operatorname{Lift}_{\mathcal G}(v)
\simeq
\operatorname{hofib}_{\delta_v}
\left(
\Omega^\infty
\operatorname{map}_B(a^\vee_{\mathcal G,c},N)
\longrightarrow
\Omega^\infty
\operatorname{map}_B(N_c,N)
\right),
}
\]
where the displayed morphism is precomposition with
$\alpha^\vee_{\mathcal G,c}$.
The equivalence is natural in affine restriction, in $N$, in structured morphisms
of variations, and under represented base change.
\end{theorem}

\begin{proof}
The represented first structured neighbourhood of the groupoid unit is
corepresented on $T$ by
$c^*u^*L_{\mathcal G_1/\overline A,s}$.  Passing through
$\operatorname{QCoh}(T)\simeq\operatorname{Mod}_B$ identifies that coefficient with
$a^\vee_{\mathcal G,c}$.  Hence a first-order source-vertical groupoid arrow with
coefficient $N$ is classified by a map
\[
\eta:a^\vee_{\mathcal G,c}\longrightarrow N.
\]
The first-order source--target theorem identifies its image under the nonlinear
anchor with the field variation classified by
\[
N_c
\xrightarrow{\alpha^\vee_{\mathcal G,c}}
a^\vee_{\mathcal G,c}
\xrightarrow{\eta}
N.
\]
On the other hand, the universal affine variation theorem identifies the given
variation $v$ with $\delta_v:N_c\to N$ through a contractible mapping space.
Therefore a symmetry lift of $v$ is exactly a choice of $\eta$ together with a
specified homotopy
\[
\eta\alpha^\vee_{\mathcal G,c}\simeq\delta_v.
\]
This is the displayed homotopy fibre.  Every construction used in this
identification is natural in affine restriction and represented base change.
\end{proof}

\begin{corollary}[Represented tangent shadow]
\label{cor:symmetry-noether-represented-tangent-shadow}
On the represented first-order theorem domain define
\[
a_{\mathcal G}
:=
\left(a_{\mathcal G}^{\vee}\right)^\vee,
\qquad
T_{\overline A/Q}^{\mathrm{rep}}
:=
L_{\overline A/Q}^{\vee}.
\]
Internal duality applied to the constructed cotangent anchor gives a canonical
morphism
\[
\boxed{
\rho_{\mathcal G}^{\mathrm{rep}}:
a_{\mathcal G}\longrightarrow
T_{\overline A/Q}^{\mathrm{rep}}.
}
\]
No smoothness, perfectness, or finite-local-freeness hypothesis is required for
this internal-dual morphism.  When the cotangent coefficients are finite locally
free, these internal duals agree with the ordinary finite-rank tangent coefficients.
Under the preceding lifting-space equivalence, a represented infinitesimal
groupoid direction maps to its represented first field variation by
$\rho_{\mathcal G}^{\mathrm{rep}}$.
\end{corollary}

\begin{remark}[No coefficient-valued tangent identification]
\label{rem:symmetry-noether-no-coefficient-valued-tangent-identification}
For a general stable coefficient $M$, no natural transformation from the
partition-Lie tangent source to $V_{\mathcal G,M}^{(1)}(A)$ is asserted.  The
coefficient-valued one-stage variation is defined by its right-Kan total-fibre
construction.  Tangent recognition is a separate represented scalar comparison.
\end{remark}

\section{Differentiated symmetry and zero-action parameters}
\label{sec:symmetry-noether-differentiated-symmetry-and-zero-action-parameters}
On the locally coherent spectral coefficient sector, retain the independently
constructed spectral partition-Lie algebroid
\[
\mathbb A_{\mathcal G}^{\pi,E_\infty}
\]
with underlying shifted anchor
\[
\rho_{\mathcal G}^{\pi,E_\infty}:
U\mathbb A_{\mathcal G}^{\pi,E_\infty}
\longrightarrow
T^{pc}_{\overline A/Q}[1].
\]

\begin{definition}[Homotopy zero-action algebroid]
\label{def:symmetry-noether-homotopy-zero-action-algebroid}
Let
\[
F_{\mathcal G}^{\pi,E_\infty}
\]
be the zero-anchor algebroid obtained from the right-adjoint anchor-fibre
construction.  Its underlying pro-coherent object is
\[
UF_{\mathcal G}^{\pi,E_\infty}
\simeq
\operatorname{fib}
(\rho_{\mathcal G}^{\pi,E_\infty}).
\]
\end{definition}

Differentiated actual isotropy carries the canonical comparison
\[
\gamma_{\mathcal G}^{\pi,E_\infty}:
i_{\mathcal G}^{\pi,E_\infty}
\longrightarrow
F_{\mathcal G}^{\pi,E_\infty}.
\]

\begin{convention}[Retained isotropy discrepancy]
\label{def:symmetry-noether-isotropy-discrepancy}
Retain from the differentiation chapter the already constructed spectral
partition-Lie isotropy discrepancy
\[
\boxed{
C_{\mathcal G}^{\mathrm{iso}}
:=
\operatorname{cofib}
\left(
U\gamma_{\mathcal G}^{\pi,E_\infty}
\right)
}
\]
in the underlying stable pro-coherent category.  The present chapter uses this
coefficient but does not redefine actual isotropy or the anchor fibre.
\end{convention}


\begin{construction}[Isotropy-exactness realization]
\label{constr:symmetry-noether-isotropy-exactness-realization}
On each coherent affine-\'etale chart, regard
\[
U\gamma_{\mathcal G}^{\pi,E_\infty}:
Ui_{\mathcal G}^{\pi,E_\infty}
\longrightarrow
UF_{\mathcal G}^{\pi,E_\infty}
\]
as a morphism in the pro-coherent coefficient category.  Apply the inverse-locus
construction in the coherent spectral-\'etale mapping sheaf and define
\[
\boxed{
\operatorname{IsoEq}_{\mathcal G}
:=
\operatorname{Inv}_{\mathrm{\acute et}}
\left(U\gamma_{\mathcal G}^{\pi,E_\infty}\right).
}
\]
A point is an extension of the comparison arrow to the walking equivalence and
therefore contains its contractibly unique coherent inverse data.
The coherent spectral-\'etale transport of $\gamma_{\mathcal G}$ makes these local
inverse spaces descend.
\end{construction}

\begin{proposition}[Universal isotropy-exact comparison]
\label{prop:symmetry-noether-universal-isotropy-exact-comparison}
Over $\operatorname{IsoEq}_{\mathcal G}$, in the sense of
Convention~\ref{conv:symmetry-noether-pullback-to-spectral-etale-realization-locus},
the canonical comparison
\[
\boxed{
\gamma_{\mathcal G}^{\pi,E_\infty}:
i_{\mathcal G}^{\pi,E_\infty}
\xrightarrow{\ \sim\ }
F_{\mathcal G}^{\pi,E_\infty}
}
\]
is an equivalence of zero-anchor spectral partition-Lie algebroids, not merely an
equivalence of their underlying pro-coherent coefficients.  Equivalently, the
pullback of the isotropy discrepancy $C_{\mathcal G}^{\mathrm{iso}}$ vanishes.  The
fibre of $\operatorname{IsoEq}_{\mathcal G}$ at a chart is empty when the
comparison is not invertible and contractible when it is invertible.
\end{proposition}

\begin{proof}
Pull the coherent coefficient stack back to the category of elements of
$\operatorname{IsoEq}_{\mathcal G}$.  By definition, its tautological point extends
the pulled-back morphism
\[
U\gamma_{\mathcal G}^{\pi,E_\infty}:
Ui_{\mathcal G}^{\pi,E_\infty}
\longrightarrow
UF_{\mathcal G}^{\pi,E_\infty}
\]
to the walking equivalence, and therefore supplies a coherent inverse together
with its unit, counit, triangle, and all higher compatibility data.  Hence
$U\gamma_{\mathcal G}^{\pi,E_\infty}$ is an equivalence in the pulled-back
pro-coherent coefficient stack.

The forgetful functor from the monadic category of spectral partition-Lie
algebroids to anchored pro-coherent coefficients is conservative and therefore
reflects equivalences.  Hence the already constructed algebroid morphism
$\gamma_{\mathcal G}^{\pi,E_\infty}$ is itself an equivalence over the same
realization locus.  In the underlying stable category a morphism is an equivalence
if and only if its cofiber is zero, giving the discrepancy statement.  The
empty-or-contractible fibre assertion is the chartwise walking-equivalence
recognition property of $\operatorname{Inv}_{\mathrm{\acute et}}$.
\end{proof}

\begin{proposition}[First-order compatibility without an operator realization]
\label{prop:symmetry-noether-first-order-compatibility-without-an-operator-realization}
Assume that the effective quotient $Q_{\mathcal G}$ lies in the represented locally
coherent quotient domain of the differentiation comparison theorem.  Then the
first structured source coefficient of the nonlinear formal anchor identifies with
\[
a_{\mathcal G}^{\vee}\simeq L_{\overline A/Q_{\mathcal G}},
\]
and its cotangent anchor is the transitivity morphism
\[
L_{\overline A/Q}
\longrightarrow
L_{\overline A/Q_{\mathcal G}}.
\]
The internal-dual represented tangent anchor and the first positive weight of the
partition-Lie Spencer action are induced by this same transitivity morphism.  The
comparison is a first-order coefficient and symbol comparison only.  No
finite-order differential operator from an ordinary parameter bundle is constructed
or required.

Outside the represented-quotient comparison domain, the formalized branch retains
the intrinsic formal-leaf ambient tangent coefficient
\[
\boxed{
T^{pc}_{L_{\mathcal G}/Q}\big|_{\overline A}
:=
\operatorname{cofib}
\left(
T^{pc}_{\overline A/L_{\mathcal G}^{\mathrm{fm},sp}}
\longrightarrow
T^{pc}_{\overline A/Q}
\right),
}
\]
canonically identified with the underlying homotopy anchor fibre.  When the formal
leaf is represented, tangent transitivity identifies this coefficient with the
pullback of the ordinary tangent of the represented leaf.
\end{proposition}

\begin{proof}
On a represented effective quotient, the represented quotient theorem identifies
the source coefficient with $L_{\overline A/Q_{\mathcal G}}$ and the first
structured anchor with cotangent transitivity.  Internal duality applies to the
transitivity morphism without requiring dualizability.  The weight-one
Atiyah--Spencer comparison on the same represented locally coherent quotient domain
identifies the first positive Spencer weight with the continuous dual of this
transitivity morphism.  Thus the three constructions agree at their common
first-order interface while their ambient categories and higher filtrations remain
distinct.

On the general coherent formalized branch, the differentiation chapter constructs
the displayed cofiber from the formal leaf and proves that it is the underlying
homotopy anchor fibre.  Its represented identification is ordinary tangent
transitivity.  No quotient presentation or additional comparison datum is used.
\end{proof}
\section{The formal symmetry variation of a Lagrangian}
\label{sec:symmetry-noether-the-formal-symmetry-variation-of-a-lagrangian}

The local Lagrangian
\[
\ell_L:\mathbf 1_{\overline A}\longrightarrow a^*M[n]
\]
may be pulled along the source and target maps of $\mathcal G_1$.

\begin{construction}[Raw symmetry difference]
\label{constr:symmetry-noether-raw-symmetry-difference}
The additive difference
\[
t^*\ell_L-s^*\ell_L
\]
is a normalized degree-one section of the raw symmetry relation coefficient.  Under
horizontal Cartan totalization it is the degree-one symmetry boundary of $L$ in the
symmetry--horizontal coherent calculus.  Write
\[
\delta_{\mathcal G}^{\infty}L
:
\mathbf1_Q\longrightarrow
\operatorname{Sec}_{\mathcal G}^{1}(A;M)[n]
\]
for the resulting full symmetry variation.
\end{construction}

\begin{proposition}[Restriction from the ambient variational differential]
\label{prop:symmetry-noether-restriction-from-the-ambient-variational-differential}
The full symmetry variation is the image of the ambient full variational derivative
under bifiltered restriction:
\[
\delta_{\mathcal G}^{\infty}L
\simeq
\operatorname{res}_{\mathcal G}
\left(
\delta_{\mathrm{var}}^{\infty}L
\right).
\]
On a symmetry-lifted split square-zero variation
$(v,\widetilde\sigma_v)$, its one-stage restriction evaluates to the alternating
first difference
\[
\Delta_v^{(1)}L.
\]
\end{proposition}

\begin{proof}
In degree one, anchor restriction is pullback of the normalized source--target
difference along
\[
\widetilde\sigma_v:T[N]\longrightarrow\mathcal G_1.
\]
Composing with the canonical anchor recovers
$\sigma_v:T[N]\to R^A_1$.  The variational chapter identifies restriction along
$\sigma_v$ with the split first cubical difference.  Naturality of the full
bifiltered restriction gives the global statement.
\end{proof}

\begin{construction}[Symmetry one-stage restriction]
\label{constr:symmetry-noether-symmetry-one-stage-restriction}
A normalized full symmetry one-cochain restricts along every symmetry lift
\[
\widetilde\sigma_v:T[N]\longrightarrow\mathcal G_1.
\]
Because the cochain is normalized, the restriction vanishes on the central
section and therefore lands in $K_M^{(1)}(v)$.  These restrictions form a cone over
the symmetry-lifted right-Kan diagram.  Hence there is a canonical morphism
\[
\lambda_{\mathcal G}^{1}:
\operatorname{Sec}_{\mathcal G}^{1}(A;M)
\longrightarrow
\Pi_aV_{\mathcal G,M}^{(1)}(A).
\]
It fits into the canonical commutative square
\begin{tikzcd}[column sep=large,row sep=large]
\operatorname{Sec}^{1}(A;M)
  \arrow[r]
  \arrow[d,"\lambda_{\mathrm{var}}^1"'] &
\operatorname{Sec}_{\mathcal G}^{1}(A;M)
  \arrow[d,"\lambda_{\mathcal G}^1"] \\
\Pi_aV_M^{(1)}(A)
  \arrow[r] &
\Pi_aV_{\mathcal G,M}^{(1)}(A).
\end{tikzcd}
where the horizontal maps are induced by restriction along the formal anchor.
\end{construction}

\begin{proof}
For every symmetry-lifted variation, composition with the canonical anchor identifies
the symmetry relation simplex with the ambient relation simplex $\sigma_v$.  Thus
restricting an ambient normalized one-cochain first to $\mathcal G_1$ and then to
the symmetry lift agrees with restricting it directly along $\sigma_v$.  The two
right-Kan cones are therefore compatible.
\end{proof}

\begin{construction}[Local mate of symmetry one-stage restriction]
\label{constr:symmetry-noether-local-mate-of-symmetry-one-stage-restriction}
Pull $\lambda_{\mathcal G}^{1}$ back along $a$ and compose with the counit of
$a^*\dashv\Pi_a$:
\[
\lambda_{\mathcal G,\mathrm{loc}}^{1}:
a^*\operatorname{Sec}_{\mathcal G}^{1}(A;M)
\xrightarrow{\,a^*\lambda_{\mathcal G}^1\,}
a^*\Pi_aV_{\mathcal G,M}^{(1)}(A)
\longrightarrow
V_{\mathcal G,M}^{(1)}(A).
\]
This is the local one-stage restriction morphism.  It is the mate of
$\lambda_{\mathcal G}^1$ and contains no additional choice.
\end{construction}

\begin{definition}[Global and local symmetry Euler classes]
\label{def:symmetry-noether-global-and-local-symmetry-euler-classes}
Define the global symmetry-restricted first Euler class
\[
\boxed{
e_{\mathcal G,L}^{(1)}
:=
\lambda_{\mathcal G}^{1}
\delta_{\mathcal G}^{\infty}L
:
\mathbf1_Q
\longrightarrow
\Pi_aV_{\mathcal G,M}^{(1)}(A)[n].
}
\]
Let
\[
\boxed{
E_{\mathcal G,L}^{(1)}
:
\mathbf1_{\overline A}
\longrightarrow
V_{\mathcal G,M}^{(1)}(A)[n]
}
\]
be its adjoint under $a^*\dashv\Pi_a$.
\end{definition}

\begin{proposition}[Symmetry Euler restriction and evaluation]
\label{prop:symmetry-noether-symmetry-euler-restriction-and-evaluation}
There is a canonical equivalence
\[
E_{\mathcal G,L}^{(1)}
\simeq
\iota_{\mathcal G}^{(1)}E_L.
\]
Equivalently, the global class is the composite
\[
\boxed{
\mathbf1_Q
\longrightarrow
\Pi_a\mathbf1_{\overline A}
\xrightarrow{\ \Pi_a(\iota_{\mathcal G}^{(1)}E_L)\ }
\Pi_aV_{\mathcal G,M}^{(1)}(A)[n],
}
\]
where the first arrow is the adjunction unit.  At every
symmetry-lifted one-variation $v$,
\[
E_{\mathcal G,L}^{(1)}(v)
\simeq
\Delta_v^{(1)}L.
\]
\end{proposition}

\begin{proof}
Apply the commutative square defining $\lambda_{\mathcal G}^1$ to the ambient full
variational derivative.  The lower-left composite is the global first Euler class
of the variational chapter, whose local adjoint is $E_L$.  The lower horizontal
map is induced by the restriction of right-Kan cones
$\iota_{\mathcal G}^{(1)}$.  Taking adjoints identifies the local composite with
$\iota_{\mathcal G}^{(1)}E_L$.  The evaluation statement follows from the
first-cubical-variation theorem of the variational chapter and the definition of
the symmetry-lifted right-Kan cone.
\end{proof}

\begin{remark}[Pairing notation]
\label{rem:symmetry-noether-pairing-notation}
On the represented tangent sector one may denote evaluation on a represented
symmetry direction $\xi$ by
\[
E_{\mathcal G,L}^{(1)}(\xi)
=
\langle E_L,\rho_{\mathcal G}^{\mathrm{rep}}(\xi)\rangle.
\]
This is an evaluative notation on the represented first-order coefficient.  It
does not assert a global tensor decomposition of $V_M^{(1)}(A)$ for arbitrary
$M$.
\end{remark}

\section{Horizontal quotients in every symmetry degree}
\label{sec:symmetry-noether-horizontal-quotients-in-every-symmetry-degree}

First totalize the horizontal direction before symmetry normalization.  As the symmetry
degree varies, define the cosimplicial stable object
\[
\mathcal D_{\mathcal G}^{\bullet}
:=
\operatorname{Tot}_H
C_{\mathcal G}^{\bullet,\bullet}(A;M),
\]
and for $q\ge1$ its stage-$q$ horizontal quotient
\[
\mathcal D_{\mathcal G}^{\bullet,<q}
:=
\operatorname{Tot}_{H,\le q-1}
C_{\mathcal G}^{\bullet,\bullet}(A;M).
\]
Thus the value in symmetry degree $p$ is
\[
\mathcal D_{\mathcal G}^{p}
=
\operatorname{Tot}_H C_{\mathcal G}^{p,\bullet}(A;M),
\qquad
\mathcal D_{\mathcal G}^{p,<q}
=
\operatorname{Tot}_{H,\le q-1}C_{\mathcal G}^{p,\bullet}(A;M).
\]

The symmetry-normalized horizontal quotient is a different object.  Define
\[
\boxed{
\mathcal D_{\mathcal G,\mathrm{nor}}^{p,<q}
:=
N_{\mathcal G}^p
\mathcal D_{\mathcal G}^{\bullet,<q}.
}
\]

\begin{lemma}[Symmetry normalization commutes with horizontal quotient]
\label{lem:symmetry-noether-symmetry-normalization-commutes-with-horizontal-quotient}
There are canonical equivalences
\[
N_{\mathcal G}^p\mathcal D_{\mathcal G}^{\bullet}
\simeq
\operatorname{Sec}_{\mathcal G}^{p}(A;M)
\]
and
\[
\boxed{
\mathcal D_{\mathcal G,\mathrm{nor}}^{p,<q}
\simeq
\operatorname{Tot}_{H,\le q-1}
B_{\mathcal G}^{p,\bullet}(A;M).
}
\]
Consequently there is a canonical fibre sequence
\[
\boxed{
F_H^q\operatorname{Sec}_{\mathcal G}^{p}(A;M)
\longrightarrow
\operatorname{Sec}_{\mathcal G}^{p}(A;M)
\longrightarrow
\mathcal D_{\mathcal G,\mathrm{nor}}^{p,<q}.
}
\]
\end{lemma}

\begin{proof}
Symmetry normalization is a finite matching-fibre construction in the symmetry
cosimplicial direction.  Horizontal totalization and horizontal partial
totalization are limits in the stable functor category.  Fubini for limits
therefore gives
\[
N_{\mathcal G}^p\operatorname{Tot}_H C_{\mathcal G}^{\bullet,\bullet}
\simeq
\operatorname{Tot}_H N_{\mathcal G}^pC_{\mathcal G}^{\bullet,\bullet}
\]
and the analogous equivalence with $\operatorname{Tot}_{H,\le q-1}$.
The first target is $\operatorname{Sec}_{\mathcal G}^p$ and the second is the
displayed normalized quotient.  The fibre sequence is then exactly the defining
horizontal-tail fibre sequence.
\end{proof}

In degree zero, normalization is the identity, so horizontal descent gives
\[
\mathcal D_{\mathcal G}^{0}
\simeq
\operatorname{Sec}^0(A;M),
\qquad
\mathcal D_{\mathcal G,\mathrm{nor}}^{0,<q}
\simeq
\mathcal D_{\mathcal G}^{0,<q}.
\]
Therefore $L$ determines a section of $\mathcal D_{\mathcal G}^{0}[n]$; let
\[
\overline L^{\,q}
:
\mathbf1_Q
\longrightarrow
\mathcal D_{\mathcal G}^{0,<q}[n]
\]
be its image in the stage-$q$ horizontal quotient.

\section{Complete symmetry realization objects}
\label{sec:symmetry-noether-complete-symmetry-realization-objects}

\begin{definition}[Partial symmetry-coherence symmetry object]
\label{def:symmetry-noether-partial-symmetry-coherence-symmetry-object}
Apply the relative zeroth-space functor to restriction from the $r$th symmetry
partial totalization to cosimplicial degree zero:
\[
\Omega_Q^\infty
\left(
\operatorname{Tot}_{\mathcal G,\le r}
\mathcal D_{\mathcal G}^{\bullet,<q}[n]
\right)
\longrightarrow
\Omega_Q^\infty
\left(
\mathcal D_{\mathcal G}^{0,<q}[n]
\right).
\]
Regard $\overline L^{\,q}$ as the corresponding section of the target over $Q$.
Define the internal parameterized symmetry object
\[
\operatorname{Sym}_{\mathcal G}^{q,\le r}(L)
:=
Q
\times_{\Omega_Q^\infty(\mathcal D_{\mathcal G}^{0,<q}[n])}
\Omega_Q^\infty
\left(
\operatorname{Tot}_{\mathcal G,\le r}
\mathcal D_{\mathcal G}^{\bullet,<q}[n]
\right)
\]
in $(\mathsf{XSp})/Q$.
\end{definition}

For $r=0$ this object is canonically equivalent to $Q$ over the chosen
Lagrangian section.  Its global $Q$-points give the corresponding space of
partial symmetry realizations.

\begin{definition}[Complete symmetry realization]
\label{def:symmetry-noether-complete-symmetry-realization}
Define the internal complete symmetry realization object by
\[
\operatorname{Sym}_{\mathcal G}^{q}(L)
:=
\lim_{r}
\operatorname{Sym}_{\mathcal G}^{q,\le r}(L)
\]
in $(\mathsf{XSp})/Q$.  Thus completeness refers to the entire symmetry simplicial
coherence at a fixed horizontal stage $q$.
\end{definition}

\begin{proposition}[Degree-one meaning]
\label{prop:symmetry-noether-degree-one-meaning}
For every test $T\to Q$, a $T$-point of
$\operatorname{Sym}_{\mathcal G}^{q,\le1}(L)$ is equivalently a section
\[
B_{\mathcal G,L}^{(q)}
:
\mathbf 1
\longrightarrow
F_H^q
\operatorname{Sec}_{\mathcal G}^1(A;M)[n]
\]
together with a specified coherent equivalence
\[
\iota_H^q B_{\mathcal G,L}^{(q)}
\simeq
\delta_{\mathcal G}^{\infty}L.
\]
Here
$\operatorname{Sec}_{\mathcal G}^1$ denotes the symmetry-normalized degree-one
horizontal totalization and $\iota_H^q$ is the horizontal-tail inclusion.
A point of the complete object additionally contains all higher symmetry simplicial
compatibilities of this boundary witness.
\end{proposition}

\begin{proof}
The first partial totalization of the symmetry-cosimplicial stable object
$\mathcal D_{\mathcal G}^{\bullet,<q}$ is the homotopy equalizer of the two
cofaces out of symmetry degree zero.  Hence lifting $\overline L^{\,q}$ through that
equalizer is precisely a specified nullhomotopy of its normalized symmetry difference
in
\[
N_{\mathcal G}^1\mathcal D_{\mathcal G}^{\bullet,<q}[n]
=
\mathcal D_{\mathcal G,\mathrm{nor}}^{1,<q}[n].
\]
The preceding normalized-quotient lemma supplies the fibre sequence
\[
F_H^q\operatorname{Sec}_{\mathcal G}^1
\longrightarrow
\operatorname{Sec}_{\mathcal G}^1
\longrightarrow
\mathcal D_{\mathcal G,\mathrm{nor}}^{1,<q}.
\]
The normalized symmetry difference of $L$ before passage to the quotient is
$\delta_{\mathcal G}^{\infty}L$.  Fibre-pasting therefore identifies a
nullhomotopy of its quotient class with a lift
$B_{\mathcal G,L}^{(q)}$ through the horizontal tail, together with
\[
\iota_H^qB_{\mathcal G,L}^{(q)}
\simeq
\delta_{\mathcal G}^{\infty}L.
\]
Higher symmetry partial totalizations impose the higher cosimplicial compatibility
cells.
\end{proof}

\begin{remark}[No inverse limit in the horizontal stage]
\label{rem:symmetry-noether-no-inverse-limit-in-the-horizontal-stage}
The index $q$ records the chosen horizontal realization depth.  A compatible
family lying in every decreasing horizontal tail would be a stronger condition
and is not declared to be the definition of symmetry.
\end{remark}

\section{Universal families of symmetry-compatible Lagrangians}
\label{sec:symmetry-noether-universal-families-of-symmetry-compatible-lagrangians}

Let
\[
\operatorname{Lag}_{A/Q}^n(M)\longrightarrow Q
\]
be the universal Lagrangian moduli object and let $L_{\mathrm{univ}}$ be its
tautological section.

\begin{construction}[Universal symmetry family]
\label{constr:symmetry-noether-universal-symmetry-family}
After base change to $\operatorname{Lag}_{A/Q}^n(M)$, apply the preceding
construction to $L_{\mathrm{univ}}$ and define
\[
\operatorname{SymLag}_{\mathcal G}^{q,\le r}(M)
:=
\operatorname{Sym}_{\mathcal G}^{q,\le r}(L_{\mathrm{univ}}),
\]
and
\[
\operatorname{SymLag}_{\mathcal G}^{q}(M)
:=
\lim_r
\operatorname{SymLag}_{\mathcal G}^{q,\le r}(M).
\]
\end{construction}

\begin{theorem}[Individual symmetry objects are pullbacks]
\label{thm:symmetry-noether-individual-symmetry-objects-are-pullbacks}
For every Lagrangian section
\[
L:Q\longrightarrow\operatorname{Lag}_{A/Q}^n(M)
\]
there are canonical equivalences
\[
L^*
\operatorname{SymLag}_{\mathcal G}^{q,\le r}(M)
\simeq
\operatorname{Sym}_{\mathcal G}^{q,\le r}(L)
\]
and
\[
L^*
\operatorname{SymLag}_{\mathcal G}^{q}(M)
\simeq
\operatorname{Sym}_{\mathcal G}^{q}(L).
\]
\end{theorem}

\begin{proof}
The symmetry--horizontal cochain object, its horizontal partial totalizations, and
symmetry partial totalizations are formed by pullback and limits from the universal
coefficient system.  Pullback preserves the finite homotopy fibres defining the
partial symmetry objects and the inverse limit in symmetry-coherence depth.
\end{proof}

\section{Symmetry restriction of the relative Lepage family}
\label{sec:symmetry-noether-symmetry-restriction-of-the-relative-lepage-family}

Fix $q\ge1$.  The relative Lepage object carries a tautological coherent full Euler
representative
\[
E_L^{(q)}
\]
and a tautological horizontal boundary realization
\[
\Theta_L^{(q)}
:
\mathbf 1_{\operatorname{Lep}^q(L)}
\longrightarrow
F_H^q\operatorname{Sec}^1(A;M)[n]
\]
such that
\[
(\pi_L^q)^*\delta_{\mathrm{var}}^\infty L
\simeq
E_L^{(q)}
+
\iota_H^q\Theta_L^{(q)}.
\]

Apply the symmetry restriction morphism of complete bifiltered objects.

\begin{construction}[Symmetry-restricted Lepage factorization]
\label{constr:symmetry-noether-symmetry-restricted-lepage-factorization}
Over $\operatorname{Lep}^q(L)$ define
\[
\mathscr E_{\mathcal G,L}^{[q]}
:=
\operatorname{res}_{\mathcal G}(E_L^{(q)})
\]
and
\[
\Theta_{\mathcal G,L}^{(q)}
:=
\operatorname{res}_{\mathcal G}(\Theta_L^{(q)}).
\]
Then there is a canonical coherent equivalence
\[
\boxed{
\delta_{\mathcal G}^{\infty}L
\simeq
\mathscr E_{\mathcal G,L}^{[q]}
+
\iota_H^q
\Theta_{\mathcal G,L}^{(q)}.
}
\]
\end{construction}

\begin{proposition}[One-stage image of the symmetry Euler representative]
\label{prop:symmetry-noether-one-stage-image-of-the-symmetry-euler-representative}
Applying the global symmetry one-stage restriction to
$\mathscr E_{\mathcal G,L}^{[q]}$ gives
\[
\lambda_{\mathcal G}^{1}
\mathscr E_{\mathcal G,L}^{[q]}
\simeq
e_{\mathcal G,L}^{(1)}
\]
after the structural pullbacks to the common parameter object.  Equivalently,
the local mate of this equality is
\[
\lambda_{\mathcal G,\mathrm{loc}}^{1}
\mathscr E_{\mathcal G,L}^{[q]}
\simeq
E_{\mathcal G,L}^{(1)}.
\]
\end{proposition}

\begin{proof}
The ambient Lepage identity is an equality of sections in the complete horizontal
tower.  Symmetry restriction is exact, so it preserves addition, fibres, the
horizontal-tail inclusion, and the specified homotopy.  The defining Helmholtz
property of the relative Lepage object says that the one-stage image of
$E_L^{(q)}$ is the first Euler class.  The commutative one-stage restriction
square and its adjoint local mate carry that statement to the two displayed
equalities.
\end{proof}

\section{The universal brave-new Noether current}
\label{sec:symmetry-noether-the-universal-brave-new-noether-current}

Define the common parameter object
\[
\mathcal N_{\mathcal G}^{q}(L)
:=
\operatorname{Sym}_{\mathcal G}^{q}(L)
\times_Q
\operatorname{Lep}^q(L).
\]
The complete symmetry point supplies
$B_{\mathcal G,L}^{(q)}$, while the Lepage point supplies
$\Theta_{\mathcal G,L}^{(q)}$.

\begin{definition}[Universal brave-new Noether current]
\label{def:symmetry-noether-universal-brave-new-noether-current}
Over $\mathcal N_{\mathcal G}^{q}(L)$ define
\[
\boxed{
J_{\mathcal G,L}^{(q)}
:=
\Theta_{\mathcal G,L}^{(q)}
-
B_{\mathcal G,L}^{(q)}.
}
\]
It is a section of the symmetry degree-one horizontal-tail current coefficient.
\end{definition}

The sign is forced by the two factorization identities and is not inserted to
imitate a classical convention.

\section{The Brave-New First Noether Theorem}
\label{sec:symmetry-noether-the-brave-new-first-noether-theorem}

\begin{theorem}[Intrinsic off-shell First Noether homotopy]
\label{thm:symmetry-noether-intrinsic-off-shell-first-noether-homotopy}
Over $\mathcal N_{\mathcal G}^{q}(L)$ there is a canonical coherent homotopy
\[
\boxed{
\iota_H^q J_{\mathcal G,L}^{(q)}
\simeq
-\mathscr E_{\mathcal G,L}^{[q]}.
}
\]
The identity is natural in the Lagrangian, the coefficient, fixed-object
formal-groupoid morphisms with the contravariant coefficient variance described
below, and arbitrary raw spectral base change.
\end{theorem}

\begin{proof}
The symmetry realization gives
\[
\iota_H^qB_{\mathcal G,L}^{(q)}
\simeq
\delta_{\mathcal G}^{\infty}L.
\]
The symmetry-restricted Lepage factorization gives
\[
\iota_H^q\Theta_{\mathcal G,L}^{(q)}
\simeq
\delta_{\mathcal G}^{\infty}L-\mathscr E_{\mathcal G,L}^{[q]}.
\]
Subtracting the first coherent equality from the second in the stable coefficient
category gives
\[
\iota_H^q
\bigl(\Theta_{\mathcal G,L}^{(q)}-B_{\mathcal G,L}^{(q)}\bigr)
\simeq
-\mathscr E_{\mathcal G,L}^{[q]}.
\]
Every term is obtained by pullback, stable limit, or subtraction from functorial
objects, so the homotopy retains all higher parameter-space coherence.
\end{proof}

\begin{corollary}[Global one-stage Euler identity]
\label{cor:symmetry-noether-global-one-stage-euler-identity}
Applying the constructed symmetry one-stage restriction gives the correctly typed
global identity
\[
\boxed{
\lambda_{\mathcal G}^1
\iota_H^qJ_{\mathcal G,L}^{(q)}
\simeq
-e_{\mathcal G,L}^{(1)}
}
\]
as a section of
$\Pi_aV_{\mathcal G,M}^{(1)}(A)[n]$, after the common parameter pullback.
\end{corollary}

\begin{proof}
Apply $\lambda_{\mathcal G}^1$ to the First Noether homotopy and use the
one-stage image proposition for $\mathscr E_{\mathcal G,L}^{[q]}$.
\end{proof}

\begin{corollary}[Local one-stage Euler evaluation]
\label{cor:symmetry-noether-local-one-stage-euler-evaluation}
Taking the local mate gives
\[
\boxed{
\lambda_{\mathcal G,\mathrm{loc}}^1
a^*\iota_H^qJ_{\mathcal G,L}^{(q)}
\simeq
-E_{\mathcal G,L}^{(1)}.
}
\]
At a represented affine point and represented infinitesimal symmetry direction
$\xi$,
\[
-E_{\mathcal G,L}^{(1)}(\xi)
\simeq
-\langle E_L,\rho_{\mathcal G}^{\mathrm{rep}}(\xi)\rangle.
\]
\end{corollary}

\begin{proof}
Adjoint the preceding global identity through $a^*\dashv\Pi_a$.  The represented
evaluation is the first-order symmetry-lift classification together with the
represented tangent shadow.
\end{proof}

\begin{corollary}[Canonical one-stage on-shell conservation]
\label{cor:symmetry-noether-canonical-one-stage-on-shell-conservation}
Put
\[
\mathcal C_{\mathcal G,L}^{q}
:=
\operatorname{Crit}(L)
\times_Q
\mathcal N_{\mathcal G}^{q}(L).
\]
After pullback to this common critical--Noether--Lepage parameter object, the local
one-stage image of the Noether current carries the canonical coherent
nullhomotopy
\[
\boxed{
\lambda_{\mathcal G,\mathrm{loc}}^1
 a^*\iota_H^q J_{\mathcal G,L}^{(q)}
\simeq0.
}
\]
Here every coefficient and section is understood after its structural pullback to
$\mathcal C_{\mathcal G,L}^{q}$.
\end{corollary}

\begin{proof}
The one-stage critical locus is the geometric zero locus $Z_{\overline A}(E_L)$ of the local first Euler
section and therefore carries
\[
j_L^*E_L\simeq0.
\]
After pullback to the common parameter object, functoriality of
$\iota_{\mathcal G}^{(1)}$ gives
\[
E_{\mathcal G,L}^{(1)}\simeq0.
\]
Apply the local one-stage First Noether identity on the same pullback base.
\end{proof}

\begin{proposition}[Full conservation is represented by compatible Euler nullification]
\label{prop:symmetry-noether-full-conservation-is-represented-by-euler-nullification}
The one-stage critical nullhomotopy does not have to be promoted by an assumption.
The derived space of promotions to a nullhomotopy of the complete coherent Euler
representative which recover the already selected one-stage critical nullhomotopy
is the compatible full Euler-nullification object constructed below.  Over that
object the First Noether homotopy canonically lifts the current to the complete
horizontal transgression fibre.
\end{proposition}

\begin{proof}
After pullback to the common critical--Noether--Lepage parameter object, the
one-stage image proposition identifies the image of the complete coherent Euler
representative with the symmetry Euler section.  The critical factor supplies its
already selected nullhomotopy.  The compatible-nullification construction
$Z_P(\sigma;u,\eta)$ therefore represents exactly the full Euler nullhomotopies
whose one-stage images equal that selected critical nullhomotopy.  The explicit
construction is given below.  Applying the stable homotopy-fibre universal
property to the First Noether homotopy then produces the complete horizontal
current transgression.
\end{proof}

\begin{remark}[The First Noether theorem is intrinsically filtration-level]
\label{rem:symmetry-noether-first-noether-is-filtration-level}
The identity
\[
\iota_H^qJ_{\mathcal G,L}^{(q)}
\simeq
-\mathscr E_{\mathcal G,L}^{[q]}
\]
is a statement in the complete horizontal filtration and is valid without an
ordinary differential realization.  No one-step integration-by-parts operator,
formal adjoint, or classical superpotential is needed to formulate or prove it.
Those are additional structures belonging to later gauge-theoretic comparison
chapters.
\end{remark}
\section{Universal Noether-current moduli}
\label{sec:symmetry-noether-universal-noether-current-moduli}

Let
\[
\operatorname{Curr}_{\mathcal G}^{q}(A;M;n)
:=
\Omega_Q^\infty
\left(
F_H^q
\operatorname{Sec}_{\mathcal G}^1(A;M)[n]
\right)
\]
be the stage-$q$ symmetry-current moduli object.

\begin{construction}[Universal Noether map]
\label{constr:symmetry-noether-universal-noether-map}
The universal current gives
\[
\operatorname{Noeth}_{\mathcal G}^{(1),q}:
\operatorname{SymLag}_{\mathcal G}^{q}(M)
\times_{\operatorname{Lag}_{A/Q}^n(M)}
\operatorname{Lep}_{\mathrm{univ}}^q
\longrightarrow
\operatorname{Curr}_{\mathcal G}^{q}(A;M;n).
\]
The corresponding critical factorization is represented by the conserved-current
object constructed immediately below.
\end{construction}


\begin{construction}[Universal one-stage conserved-current object]
\label{constr:symmetry-noether-universal-one-stage-conserved-current-object}
Compose the horizontal-tail inclusion with the constructed global one-stage
restriction:
\[
b_{\mathcal G}^{q}
:=
\lambda_{\mathcal G}^{1}\iota_H^q:
F_H^q\operatorname{Sec}_{\mathcal G}^{1}(A;M)[n]
\longrightarrow
\Pi_aV_{\mathcal G,M}^{(1)}(A)[n].
\]
Its adjoint field-evaluated morphism is
\[
\boxed{
b_{\mathcal G,\mathrm{loc}}^{q}:
a^*F_H^q\operatorname{Sec}_{\mathcal G}^{1}(A;M)[n]
\longrightarrow
V_{\mathcal G,M}^{(1)}(A)[n]
}
\]
defined by
\[
a^*F_H^q\operatorname{Sec}_{\mathcal G}^{1}(A;M)[n]
\xrightarrow{\ a^*b_{\mathcal G}^{q}\ }
a^*\Pi_aV_{\mathcal G,M}^{(1)}(A)[n]
\longrightarrow
V_{\mathcal G,M}^{(1)}(A)[n],
\]
where the final arrow is the counit of $a^*\dashv\Pi_a$.  Equivalently,
\[
b_{\mathcal G,\mathrm{loc}}^{q}
\simeq
\lambda_{\mathcal G,\mathrm{loc}}^{1}\,a^*\iota_H^q.
\]

Define the field-evaluated current object by
\[
\boxed{
\operatorname{Curr}_{\mathcal G,\mathrm{fld}}^{q}(A;M;n)
:=
\Omega_{\overline A}^\infty
\left(
a^*F_H^q\operatorname{Sec}_{\mathcal G}^{1}(A;M)[n]
\right).
}
\]
Pullback compatibility of relative zeroth space gives a canonical equivalence
\[
\operatorname{Curr}_{\mathcal G,\mathrm{fld}}^{q}(A;M;n)
\simeq
\overline A\times_Q\operatorname{Curr}_{\mathcal G}^{q}(A;M;n).
\]
After relative zeroth space, $b_{\mathcal G,\mathrm{loc}}^{q}$ gives
\[
\overline b_{\mathcal G,\mathrm{loc}}^{q}:
\operatorname{Curr}_{\mathcal G,\mathrm{fld}}^{q}(A;M;n)
\longrightarrow
\Omega_{\overline A}^\infty
\left(V_{\mathcal G,M}^{(1)}(A)[n]\right).
\]
Define the field-evaluated one-stage conserved-current object by the zero locus
\[
\boxed{
\operatorname{ConsCurr}_{\mathcal G,\mathrm{fld}}^{q}(A;M;n)
:=
\overline A
\times_{\Omega_{\overline A}^\infty(V_{\mathcal G,M}^{(1)}(A)[n])}
\operatorname{Curr}_{\mathcal G,\mathrm{fld}}^{q}(A;M;n),
}
\]
where $\overline A$ maps to the target by the zero section.  Thus a test point is a
field point, a stage-$q$ current evaluated at that field, and a specified coherent
nullhomotopy of its local one-stage image.

For the genuinely universal critical family, put
\[
\mathscr L:=\operatorname{Lag}_{A/Q}^n(M)
\]
and retain from the variational chapter the universal one-stage critical family
\[
\operatorname{Crit}_{A/Q}^{(1)}(M;n)
\longrightarrow
\overline A\times_Q\mathscr L.
\]
Put
\[
\mathscr N_{\mathcal G}^{q}
:=
\operatorname{SymLag}_{\mathcal G}^{q}(M)
\times_{\mathscr L}
\operatorname{Lep}_{\mathrm{univ}}^q
\]
and
\[
\boxed{
\mathscr C_{\mathcal G}^{q}
:=
\operatorname{Crit}_{A/Q}^{(1)}(M;n)
\times_{\mathscr L}
\mathscr N_{\mathcal G}^{q}.
}
\]
The projection through the universal critical family gives
$\mathscr C_{\mathcal G}^{q}\to\overline A$, so the field-evaluated current and its
conservation zero locus are defined on this parameter object.  Pull the universal
Noether current to $\mathscr C_{\mathcal G}^{q}$ and evaluate it along this field
projection; write
\[
J_{\mathcal G,\mathrm{univ}}^{(q)}
\]
for the resulting field-evaluated current.  Write $E_{\mathrm{univ}}$ for the
universal local one-stage Euler section on
$\overline A\times_Q\mathscr L$ and, after the structural pullbacks to the common
parameter object, put
\[
E_{\mathcal G,\mathrm{univ}}^{(1)}
:=
\iota_{\mathcal G}^{(1)}E_{\mathrm{univ}}.
\]
\end{construction}

\begin{proposition}[Universal and individual one-stage conserved-current factorizations]
\label{prop:symmetry-noether-universal-and-individual-one-stage-conserved-current-factorizations}
The universal Noether current and the universal one-stage critical nullhomotopy
determine a canonical morphism
\[
\boxed{
\operatorname{Noeth}_{\mathcal G,\mathrm{crit},\mathrm{univ}}^{(1),q}:
\mathscr C_{\mathcal G}^{q}
\longrightarrow
\operatorname{ConsCurr}_{\mathcal G,\mathrm{fld}}^{q}(A;M;n).
}
\]
For every Lagrangian section
\[
L:Q\longrightarrow\mathscr L,
\]
its pullback is the canonical fixed-L factorization
\[
\boxed{
\operatorname{Noeth}_{\mathcal G,L,\mathrm{crit}}^{(1),q}:
\mathcal C_{\mathcal G,L}^{q}
=
\operatorname{Crit}(L)\times_Q\mathcal N_{\mathcal G}^{q}(L)
\longrightarrow
\operatorname{ConsCurr}_{\mathcal G,\mathrm{fld}}^{q}(A;M;n).
}
\]
Hence one-stage on-shell conservation is represented by an actual internal moduli
object over the field geometry, and the individual factorization is obtained by
base change from the universal critical family.
\end{proposition}

\begin{proof}
Apply the local one-stage First Noether identity to the tautological Lagrangian on
$\mathscr L$.  After pullback to $\mathscr C_{\mathcal G}^{q}$ it identifies the
field-evaluated image of the universal current with the negative universal symmetry
Euler section:
\[
b_{\mathcal G,\mathrm{loc}}^{q}
\left(J_{\mathcal G,\mathrm{univ}}^{(q)}\right)
\simeq
-E_{\mathcal G,\mathrm{univ}}^{(1)}.
\]
The universal one-stage critical family carries the tautological coherent
nullhomotopy of the universal local Euler section.  Under the canonical equivalence
\[
E_{\mathcal G,\mathrm{univ}}^{(1)}
\simeq
\iota_{\mathcal G}^{(1)}E_{\mathrm{univ}},
\]
functoriality carries that critical nullhomotopy to a nullhomotopy of
$E_{\mathcal G,\mathrm{univ}}^{(1)}$.  Composing with the displayed local First
Noether homotopy gives a specified nullhomotopy of the image of the universal
current under $b_{\mathcal G,\mathrm{loc}}^{q}$.  The parameterized zero-locus
universal property therefore gives
$\operatorname{Noeth}_{\mathcal G,\mathrm{crit},\mathrm{univ}}^{(1),q}$.

The variational chapter identifies the pullback of
$\operatorname{Crit}_{A/Q}^{(1)}(M;n)$ along $L$ with $\operatorname{Crit}(L)$,
while Theorem~\ref{thm:symmetry-noether-individual-symmetry-objects-are-pullbacks}
and the corresponding Lepage pullback identify the pullback of
$\mathscr N_{\mathcal G}^{q}$ with $\mathcal N_{\mathcal G}^{q}(L)$.  The universal
local First Noether homotopy and the universal critical nullhomotopy pull back to
the already constructed fixed-L homotopies.  Thus the universal factorization pulls
back to the displayed map from $\mathcal C_{\mathcal G,L}^{q}$, with exactly the
canonical one-stage conservation nullhomotopy of
Corollary~\ref{cor:symmetry-noether-canonical-one-stage-on-shell-conservation}.
\end{proof}

\begin{proposition}[Individual current families are pullbacks]
\label{prop:symmetry-noether-individual-current-families-are-pullbacks}
Every individual Noether-current family, every path between symmetry or Lepage
realizations, and every higher homotopy between such families is obtained from the
universal Noether map by pullback along the corresponding universal points.
\end{proposition}

\section{The complete native symmetry-coherence obstruction tower}
\label{sec:symmetry-noether-the-complete-native-symmetry-coherence-obstruction-tower}

Fix the horizontal stage $q$ and abbreviate
\[
\mathcal Q_{\mathcal G}^{\bullet,q}
:=
\mathcal D_{\mathcal G}^{\bullet,<q}[n].
\]
This is a cosimplicial stable object in the symmetry direction.

For $r\ge1$, stable Dold--Kan gives a canonical fibre sequence
\[
N_{\mathcal G}^r
\mathcal Q_{\mathcal G}^{\bullet,q}[-r]
\longrightarrow
\operatorname{Tot}_{\mathcal G,\le r}
\mathcal Q_{\mathcal G}^{\bullet,q}
\longrightarrow
\operatorname{Tot}_{\mathcal G,\le r-1}
\mathcal Q_{\mathcal G}^{\bullet,q}.
\]
Equivalently there is a connecting morphism
\[
\partial_{\mathcal G}^{r,q}:
\operatorname{Tot}_{\mathcal G,\le r-1}
\mathcal Q_{\mathcal G}^{\bullet,q}
\longrightarrow
N_{\mathcal G}^r
\mathcal Q_{\mathcal G}^{\bullet,q}[1-r].
\]

\begin{definition}[Native stage-$r$ Noether obstruction coefficient]
\label{def:symmetry-noether-native-stage-r-noether-obstruction-coefficient}
Define the stable obstruction coefficient
\[
\boxed{
\mathcal O_{\mathcal G}^{r,q}(A;M,n)
:=
N_{\mathcal G}^r
\mathcal D_{\mathcal G}^{\bullet,<q}(A;M)[n+1-r]
=
N_{\mathcal G}^r
\mathcal Q_{\mathcal G}^{\bullet,q}[1-r].
}
\]
The coefficient is determined by the symmetry--horizontal cochain system; the
Lagrangian enters only through the obstruction section below.
\end{definition}

\begin{construction}[Universal stage-$r$ Noether obstruction section]
\label{constr:symmetry-noether-universal-stage-r-noether-obstruction-section}
Pull $\partial_{\mathcal G}^{r,q}$ and the stable obstruction coefficient back to
$\operatorname{Sym}_{\mathcal G}^{q,\le r-1}(L)$ and evaluate the connecting
morphism on the universal partial symmetry section.  After relative zeroth space this gives
\[
\boxed{
\mathfrak N_{\mathcal G,L}^{r,q}:
\operatorname{Sym}_{\mathcal G}^{q,\le r-1}(L)
\longrightarrow
\Omega^\infty
\mathcal O_{\mathcal G}^{r,q}(A;M,n).
}
\]
It is the stage-$r$ native Noether obstruction section.
\end{construction}

\begin{definition}[Noether-nullhomotopy object]
\label{def:symmetry-noether-noether-nullhomotopy-object}
Using the parameterized stable zero-locus construction, define
\[
\boxed{
\operatorname{NullNId}_{\mathcal G}^{r,q}(L)
:=
Z_{\operatorname{Sym}_{\mathcal G}^{q,\le r-1}(L)}
\left(\mathfrak N_{\mathcal G,L}^{r,q}\right).
}
\]
Its fibre over a partial symmetry is exactly the full space of specified coherent
nullhomotopies of the corresponding obstruction section.
\end{definition}

\section{The native symmetry-coherence Noether obstruction theorem}
\label{sec:symmetry-noether-the-native-symmetry-coherence-noether-obstruction-theorem}

\begin{theorem}[Stagewise native symmetry-coherence lifting theorem]
\label{thm:symmetry-noether-stagewise-native-symmetry-coherence-lifting-theorem}
For every $r\ge1$, there is a canonical equivalence over
$\operatorname{Sym}_{\mathcal G}^{q,\le r-1}(L)$
\[
\boxed{
\operatorname{Sym}_{\mathcal G}^{q,\le r}(L)
\simeq
\operatorname{NullNId}_{\mathcal G}^{r,q}(L).
}
\]
Thus the fibre over every partial $(r-1)$-stage symmetry realization is the full
space of specified nullhomotopies of
$\mathfrak N_{\mathcal G,L}^{r,q}$.
Consequently
\[
\boxed{
\operatorname{Sym}_{\mathcal G}^{q}(L)
\simeq
\lim_r
\operatorname{Sym}_{\mathcal G}^{q,\le r}(L)
}
\]
is exactly the internal parameterized object of coherent successive nullifications
of the complete native symmetry-coherence obstruction tower.
\end{theorem}

\begin{proof}
In a stable $\infty$-category, the displayed Dold--Kan fibre sequence identifies
extension from the $(r-1)$st partial totalization to the $r$th partial
totalization with a lift through the fibre
\[
N_{\mathcal G}^r
\mathcal Q_{\mathcal G}^{\bullet,q}[-r].
\]
Equivalently, by the associated connecting morphism, such a lift is a specified
nullhomotopy of
$\mathfrak N_{\mathcal G,L}^{r,q}$.
Mapping out of the tensor unit preserves this homotopy-fibre description.
The partial symmetry object is itself the fibre over the fixed degree-zero
Lagrangian, so fibre-pasting gives the displayed equivalence.  Taking the inverse
limit over $r$ is the definition of complete symmetry totalization and therefore
retains all higher simplicial coherence.
\end{proof}

\begin{corollary}[First Noether identity as the first obstruction]
\label{cor:symmetry-noether-first-noether-identity-as-the-first-obstruction}
For $r=1$,
\[
\mathfrak N_{\mathcal G,L}^{1,q}
\]
is the image of the normalized symmetry variation
$\delta_{\mathcal G}^{\infty}L$ in the symmetry-normalized horizontal quotient
\[
\mathcal D_{\mathcal G,\mathrm{nor}}^{1,<q}[n]
\simeq
N_{\mathcal G}^1\mathcal D_{\mathcal G}^{\bullet,<q}[n].
\]
Through the symmetry-restricted Lepage identity it is equally the normalized
horizontal-quotient class of $\mathscr E_{\mathcal G,L}^{[q]}$.
Thus a stage-$q$ infinitesimal symmetry supplies a specified nullhomotopy of
the symmetry Euler class modulo the stage-$q$ horizontal tail.
\end{corollary}

\begin{proof}
For $r=1$ the connecting morphism of the first partial totalization is the
normalized degree-one symmetry boundary.  The Lepage factorization differs from
$\delta_{\mathcal G}^{\infty}L$ by a term in
$F_H^q$, which maps to zero in the horizontal quotient.
\end{proof}

\begin{construction}[Native ambiguity coefficients of the symmetry tower]
\label{constr:symmetry-noether-native-ambiguity-coefficients}
Over $\operatorname{Sym}_{\mathcal G}^{q,\le r-1}(L)$ pull back the stage-$r$
obstruction coefficient and define
\[
\boxed{
\mathcal A_{\mathcal G}^{r,q}
:=
\Omega^\infty
\mathcal O_{\mathcal G}^{r,q}(A;M,n)[-1].
}
\]
For a fixed partial symmetry realization $x$, the stage-$r$ lift fibre is the path
space
\[
\operatorname{Path}
\left(
\mathfrak N_{\mathcal G,L}^{r,q}(x),0
\right)
\]
in the infinite-loop space of the obstruction spectrum.  Concatenation with loops
at zero gives a canonical action of
$\mathcal A_{\mathcal G}^{r,q}|_x$ on this lift fibre.
\end{construction}

\begin{theorem}[Native obstruction, ambiguity torsor, and higher automorphisms]
\label{thm:symmetry-noether-native-obstruction-ambiguity-and-higher-automorphisms}
For every $r\ge1$ and every partial $(r-1)$-stage symmetry realization $x$:
\begin{enumerate}
\item the obstruction to extension is the constructed point
      $\mathfrak N_{\mathcal G,L}^{r,q}(x)$;
\item if the stage-$r$ lift fibre is inhabited, it is a torsor for
      $\mathcal A_{\mathcal G}^{r,q}|_x$;
\item for a selected stage-$r$ lift, its $k$th higher automorphism space is canonically
      the corresponding value of
      \[
      \Omega^\infty
      \mathcal O_{\mathcal G}^{r,q}(A;M,n)[-k-1],
      \qquad k\ge1.
      \]
\end{enumerate}
Hence the native symmetry-totalization tower itself constructs its obstruction,
ambiguity, and higher-reducibility information.  No externally chosen reducibility
complex is required.  This hierarchy is not identified with the differentiated
morphism-latching syzygy hierarchy below.
\end{theorem}

\begin{proof}
The stagewise lifting theorem identifies the lift fibre with a path space from the
obstruction point to zero.  The loop space at zero acts on such a path space by
concatenation, and whenever a path exists this action is simply transitive.  Thus the
ambiguity torsor is governed by
\[
\Omega\Omega^\infty\mathcal O
\simeq
\Omega^\infty\mathcal O[-1].
\]
For a selected lift, however, its $k$th higher automorphism space is the $k$fold loop
space of the \emph{lift fibre}.  Since the lift fibre is itself a path space,
\[
\Omega^k\operatorname{Path}(\mathfrak N(x),0)
\simeq
\Omega^{k+1}\Omega^\infty\mathcal O
\simeq
\Omega^\infty\mathcal O[-k-1].
\]
This gives the higher automorphism statement with all coherences inherited from the
iterated loop-space action.
\end{proof}



\begin{remark}[This is not yet the differentiated Second Noether theorem]
\label{rem:symmetry-noether-this-is-not-yet-the-differentiated-second-noether-theorem}
The index $r$ measures extension of cosimplicial symmetry coherence.  It is distinct
from the Artinian Taylor weight and from the one-marked tangent weight.  The
differentiated Second Noether theorem below arises after the formalized Noether
section has been reflected, differentiated to a partition-Lie algebroid splitting,
and compared with the zero splitting.
\end{remark}

\section{Formalized Noether comparison}
\label{sec:symmetry-noether-formalized-noether-comparison}

The native symmetry-coherence obstruction theorem is constructed directly from the raw formal
symmetry relation.  We now form the quotient-formalized comparison without
identifying the two branches.

Let
\[
R_{\mathcal G,\bullet}^{\mathrm{fm}}
\]
be the Artinian-formalized relation of the Atiyah--Spencer chapter.  There is a
canonical simplicial map
\[
\lambda_{\mathcal G}^R:
\mathcal G_\bullet
\longrightarrow
R_{\mathcal G,\bullet}^{\mathrm{fm}},
\]
and therefore a contravariant restriction on stable cochains from the formalized
relation to the raw relation.  More explicitly, for every $(p,q)$ let
\[
\widetilde\lambda_{\mathcal G}^{p,q}
:=
\operatorname{id}_{H_q}\times_Q\lambda_{\mathcal G,p}^R
:
H_q\times_Q\mathcal G_p
\longrightarrow
H_q\times_QR_{\mathcal G,p}^{\mathrm{fm}}.
\]
If
\[
r_{p,q}^{\mathrm{raw}}:
H_q\times_Q\mathcal G_p\to Q,
\qquad
r_{p,q}^{\mathrm{fm}}:
H_q\times_QR_{\mathcal G,p}^{\mathrm{fm}}\to Q,
\]
then Construction~\ref{constr:symmetry-noether-stable-restriction-along-a-relation-morphism}
gives the coefficient map
\[
\begin{aligned}
\Pi_{r_{p,q}^{\mathrm{fm}}}
(r_{p,q}^{\mathrm{fm}})^*M
&\xrightarrow{\ \Pi_{r_{p,q}^{\mathrm{fm}}}\eta_{\widetilde\lambda}\ }
\Pi_{r_{p,q}^{\mathrm{fm}}}
\Pi_{\widetilde\lambda_{\mathcal G}^{p,q}}
(\widetilde\lambda_{\mathcal G}^{p,q})^*
(r_{p,q}^{\mathrm{fm}})^*M\\
&\simeq
\Pi_{r_{p,q}^{\mathrm{raw}}}
(r_{p,q}^{\mathrm{raw}})^*M.
\end{aligned}
\]
The simplicial compatibility lemma makes these maps a morphism of the complete
formalized symmetry--horizontal cochain system to the raw one.  Thus the
formalized-to-raw comparison is defined for the stable coefficient $M$ itself; it
is not obtained by first reducing to a scalar cochain comparison.

Apply the same horizontal quotient, symmetry normalization, partial totalization,
and connecting-morphism construction to the formalized relation on every
coefficient domain on which its stable Artinian coefficient has been constructed.
Write
\[
\mathcal O_{\mathcal G}^{\mathrm{fm},r,q}(A;M,n)
\]
for the resulting stable stage-$r$ formalized obstruction coefficient and
\[
\mathfrak N_{\mathcal G,L}^{\mathrm{fm},r,q}:
\operatorname{Sym}_{\mathcal G}^{\mathrm{fm},q,\le r-1}(L)
\longrightarrow
\Omega^\infty
\mathcal O_{\mathcal G}^{\mathrm{fm},r,q}(A;M,n)
\]
for its universal obstruction section.

In the same way define formalized partial and complete symmetry realization
objects
\[
\operatorname{Sym}_{\mathcal G}^{\mathrm{fm},q,\le r}(L),
\qquad
\operatorname{Sym}_{\mathcal G}^{\mathrm{fm},q}(L)
:=
\lim_r
\operatorname{Sym}_{\mathcal G}^{\mathrm{fm},q,\le r}(L).
\]
Using the stable zero-locus construction, put
\[
\operatorname{NullNId}_{\mathcal G}^{\mathrm{fm},r,q}(L)
:=
Z_{\operatorname{Sym}_{\mathcal G}^{\mathrm{fm},q,\le r-1}(L)}
\left(\mathfrak N_{\mathcal G,L}^{\mathrm{fm},r,q}\right).
\]
Stable Dold--Kan gives formalized stagewise lifting equivalences
\[
\operatorname{Sym}_{\mathcal G}^{\mathrm{fm},q,\le r}(L)
\simeq
\operatorname{NullNId}_{\mathcal G}^{\mathrm{fm},r,q}(L).
\]
These are formed entirely on the quotient-formalized coefficient branch.

\begin{construction}[Formalized-to-raw Noether comparison]
\label{constr:symmetry-noether-formalized-to-raw-noether-comparison}
Restriction along $\lambda_{\mathcal G}^R$ induces natural morphisms
\[
\operatorname{Sym}_{\mathcal G}^{\mathrm{fm},q,\le r}(L)
\longrightarrow
\operatorname{Sym}_{\mathcal G}^{q,\le r}(L)
\]
and, after passage to complete totalization,
\[
\operatorname{Sym}_{\mathcal G}^{\mathrm{fm},q}(L)
\longrightarrow
\operatorname{Sym}_{\mathcal G}^{q}(L).
\]
On the stable obstruction coefficients it induces
\[
\boxed{
\Lambda_{\mathcal G}^{\mathcal O,r,q}:
\mathcal O_{\mathcal G}^{\mathrm{fm},r,q}(A;M,n)
\longrightarrow
\mathcal O_{\mathcal G}^{r,q}(A;M,n).
}
\]
After pulling the two stable obstruction coefficients to the indicated symmetry
parameter objects, naturality of the connecting morphisms gives the commutative square of
parameterized obstruction sections
\begin{tikzcd}[column sep=huge,row sep=large]
\operatorname{Sym}_{\mathcal G}^{\mathrm{fm},q,\le r-1}(L)
  \arrow[r,"\mathfrak N_{\mathcal G,L}^{\mathrm{fm},r,q}"]
  \arrow[d] &
\Omega^\infty\mathcal O_{\mathcal G}^{\mathrm{fm},r,q}
  \arrow[d,"\Omega^\infty\Lambda_{\mathcal G}^{\mathcal O,r,q}"] \\
\operatorname{Sym}_{\mathcal G}^{q,\le r-1}(L)
  \arrow[r,"\mathfrak N_{\mathcal G,L}^{r,q}"'] &
\Omega^\infty\mathcal O_{\mathcal G}^{r,q}.
\end{tikzcd}
Whenever both stable coefficients are formed in one stable coefficient category,
define the formalization discrepancy by
\[
\boxed{
C_{\mathcal G,L}^{\mathrm{Noeth},r,q}
:=
\operatorname{cofib}
\left(
\Lambda_{\mathcal G}^{\mathcal O,r,q}
\right).
}
\]
The discrepancy is therefore the cofiber of a stable coefficient comparison, not
the cofiber of two parameterized obstruction sections.
\end{construction}

\begin{proposition}[Variance and compatibility]
\label{prop:symmetry-noether-variance-and-compatibility}
The comparison is contravariant on cochains and covariant on the underlying
relation geometry.  It commutes with the horizontal quotient maps, symmetry
normalization, all partial-totalization transition maps, the stable obstruction
coefficients, and the connecting morphisms defining the obstruction sections.
It consequently induces the displayed maps of nullification objects.
\end{proposition}

\begin{proof}
The relation map is simplicial and therefore commutes with every face and
degeneracy.  Stable sections are contravariant in relation maps.  Normalization,
horizontal partial totalization, symmetry partial totalization, fibres, and connecting
morphisms are functorial exact constructions, so the stable coefficient comparison
propagates through the entire tower.  Applying relative zeroth space and the
parameterized zero-locus construction gives the corresponding square of
obstruction sections and the induced maps of their nullification objects.
\end{proof}

\begin{theorem}[Represented quotient collapse]
\label{thm:symmetry-noether-represented-quotient-collapse}
If the effective formal quotient $Q_{\mathcal G}$ is represented on the theorem
domain of the Atiyah--Spencer represented-quotient comparison, then
\[
Q_{\mathcal G}\longrightarrow Q_{\mathcal G}^{\mathrm{fm}}
\]
is an equivalence.  Consequently the raw and formalized relation nerves, the
oppositely varianced jet comparison, scalar restriction, Cartan-cobar comparison,
and every stable Noether obstruction-coefficient comparison
\[
\boxed{
\Lambda_{\mathcal G}^{\mathcal O,r,q}
}
\]
are equivalences in every degree for which the corresponding coefficient has been
constructed.  The discrepancy
$C_{\mathcal G,L}^{\mathrm{Noeth},r,q}$ therefore vanishes, and the formalized and
raw obstruction sections correspond under the preceding naturality square.
\end{theorem}

\begin{proof}
The represented quotient theorem identifies quotient formalization with the
identity on represented effective formal quotients.  The two relation nerves are
therefore degreewise equivalent.  Horizontal partial totalization, symmetry
normalization, Dold--Kan layer formation, and the connecting morphisms preserve
these equivalences.  Hence the stable obstruction coefficients are equivalent and
the obstruction sections are identified by naturality.
\end{proof}


\section{Schlessinger-reflected formalized Noether coefficients}
\label{sec:symmetry-noether-schlessinger-reflected-formalized-noether-coefficients}

The formalized first obstruction varies over actual Artinian points of a formal
leaf.  We first assemble it as a stable Artinian coefficient, then apply the stable
Schlessinger reflector, and only afterwards differentiate.  The section itself is
kept at the relative infinite-loop level; no assertion that the reflector preserves
a chosen fibrewise sphere or tensor unit is required.

Fix a coherent affine-\'etale formal leaf
\[
F
:=
L_{\mathcal G,U/V}^{\mathrm{fm},\mathrm{sp}}
\in
\operatorname{FMP}^{\mathrm{sp}}_{B/R},
\qquad
U=\operatorname{Spec}B,\quad
V=\operatorname{Spec}R,
\]
and put
\[
\mathcal P_F
:=
\operatorname{Fun}
\left(
\mathsf{Art}^{\mathrm{sp}}_{R/B},\mathcal S
\right)/F.
\]

For an Artinian arrow $A\to B$ and a point $x\in F(A)$, the canonical Artinian
\v{C}ech test map into $R_{\mathcal G,\bullet}^{\mathrm{fm}}$ restricts the
formalized symmetry--horizontal coefficient to that represented test.  Let
\[
\mathcal E_{\mathcal G,L}^{1,q}(A,x)
\]
denote the resulting stable normalized symmetry-degree-one stage-$q$ horizontal
quotient coefficient.

\begin{construction}[Universal formalized first-obstruction coefficient]
\label{constr:symmetry-noether-universal-formalized-first-obstruction-coefficient}
Unstraightening
\[
x\longmapsto
\mathcal E_{\mathcal G,L}^{1,q}(A,x)
\]
over $F(A)$ and then varying $A\to B$ gives
\[
\mathcal O_{\mathcal G,L}^{\mathrm{Art},1,q}
\in
\operatorname{Sp}(\mathcal P_F).
\]
The formalized first Noether obstruction gives a section of its relative
infinite-loop object
\[
\boxed{
\mathbf n_{\mathcal G,L}^{1,q}:
F
\longrightarrow
\Omega_F^\infty
\mathcal O_{\mathcal G,L}^{\mathrm{Art},1,q}[n].
}
\]
Its value at $(A,x)$ is the formalized first Noether obstruction at that Artinian
point.
\end{construction}

\begin{proof}
The formalized relation, stable coefficient pullback, horizontal partial
totalization, symmetry normalization, and first connecting morphism are functorial in
an Artinian point.  Their values therefore assemble to a spectrum object over
$F$.  The formalized symmetry boundary is natural before any of these operations and
hence defines the displayed relative section.
\end{proof}

\begin{construction}[Stable Schlessinger reflection of the coefficient and section]
\label{constr:symmetry-noether-stable-schlessinger-reflection-of-the-coefficient-and-section}
Apply the exact stable Schlessinger reflector
\[
L_F^{\mathrm{Sch},\mathrm{st}}:
\operatorname{Sp}(\mathcal P_F)
\longrightarrow
\operatorname{Rep}^{\mathrm{fm}}(F)
\]
and put
\[
\widehat{\mathcal O}_{\mathcal G,L}^{1,q}
:=
L_F^{\mathrm{Sch},\mathrm{st}}
\left(
\mathcal O_{\mathcal G,L}^{\mathrm{Art},1,q}
\right).
\]
Let
\[
\eta_{\mathcal O}:
\mathcal O_{\mathcal G,L}^{\mathrm{Art},1,q}
\longrightarrow
\widehat{\mathcal O}_{\mathcal G,L}^{1,q}
\]
be the localization unit.  Define the reflected Noether section by
\[
\boxed{
\widehat{\mathbf n}_{\mathcal G,L}^{1,q}
:=
\Omega_F^\infty(\eta_{\mathcal O}[n])
\circ
\mathbf n_{\mathcal G,L}^{1,q}
:
F\longrightarrow
\Omega_F^\infty
\widehat{\mathcal O}_{\mathcal G,L}^{1,q}[n].
}
\]
No source object is reflected and no unit-preservation theorem is used.
\end{construction}

\begin{construction}[Relative unstable Schlessinger reflector over a formal leaf]
\label{constr:symmetry-noether-relative-unstable-schlessinger-reflector-over-a-formal-leaf}
The differentiated Second Noether theorem requires the symmetry parameter itself,
not only its stable obstruction coefficient, to lie in the formal-moduli category.
We therefore form the unstable relative reflector explicitly.

For every Artinian Schlessinger generator
\[
w:K\longrightarrow L
\]
of Chapter 14 and every morphism $L\to F$, regard $w$ as a morphism in the slice
$\mathcal P_F$.  Let $W_F^{\mathrm{Sch}}$ be the resulting set of relative
generators.  Accessible localization gives
\[
\boxed{
L_F^{\mathrm{Sch}}:
\mathcal P_F
\longrightarrow
\bigl(\operatorname{FMP}^{\mathrm{sp}}_{B/R}\bigr)/F .
}
\]
Its local objects are precisely the maps $Z\to F$ whose underlying Artinian
presheaf is a spectral formal-moduli problem.
\end{construction}

\begin{proof}
The category $\mathcal P_F$ is a slice of a presentable functor category and is
therefore presentable.  The fixed-universe family $W_F^{\mathrm{Sch}}$ is a set, so
the accessible localization exists.

It remains to identify its local objects.  Fix an ordinary Schlessinger generator
\[
w:K\longrightarrow L
\]
and a map $\ell:L\to F$.  For a map $z:Z\to F$, consider the commutative square
\[
\begin{tikzcd}[column sep=large,row sep=large]
\operatorname{Map}(L,Z)
  \arrow[r]
  \arrow[d] &
\operatorname{Map}(K,Z)
  \arrow[d] \\
\operatorname{Map}(L,F)
  \arrow[r] &
\operatorname{Map}(K,F).
\end{tikzcd}
\]
Because $F$ is already a spectral formal-moduli problem, the bottom horizontal map
is an equivalence.  The mapping space in the slice from $(L,\ell)$ to $Z$ is the
homotopy fibre of the left vertical map over $\ell$; the corresponding mapping
space from $(K,\ell w)$ is the fibre of the right vertical map over $\ell w$.

Suppose first that the underlying Artinian presheaf $Z$ is Schlessinger-local.
Then the top horizontal map in the displayed square is an equivalence.  Taking the homotopy
fibres over the corresponding points of the equivalent bottom mapping spaces shows
that $Z\to F$ is local for the relative generator associated to $(w,\ell)$.
Hence every ordinary Schlessinger-local object is relative-local.

Conversely suppose that $Z\to F$ is local for every relative generator.  Let
$k\in\operatorname{Map}(K,F)$.  Since the bottom map in the displayed square is an equivalence,
the space of lifts $\ell\in\operatorname{Map}(L,F)$ with $\ell w\simeq k$ is
contractible.  Relative locality identifies, for every such lift, the corresponding
homotopy fibres of the two vertical maps.  Thus the top horizontal map in the displayed square
is a map over an equivalence of bases which is an equivalence on every homotopy
fibre.  It is therefore itself an equivalence.  This holds for every Schlessinger
generator $w$, so the underlying Artinian presheaf of $Z$ is a spectral
formal-moduli problem.

Thus the local objects are precisely
\[
\bigl(\operatorname{FMP}^{\mathrm{sp}}_{B/R}\bigr)/F,
\]
which proves the stated target of the localization.
\end{proof}

\begin{proposition}[Spectral-\'etale exchange for the relative unstable Schlessinger reflector]
\label{prop:symmetry-noether-spectral-etale-exchange-for-relative-unstable-schlessinger-reflector}
Let
\[
f:(R,B)\longrightarrow(R',B')
\]
be a coherent one- or two-variable spectral-\'etale coefficient transition in the
sense of Chapter 14, let
\[
F\in\operatorname{FMP}^{\mathrm{sp}}_{B/R},
\qquad
F'\in\operatorname{FMP}^{\mathrm{sp}}_{B'/R'}
\]
be corresponding transported formal leaves, and let
\[
f_{\mathcal P,!}:\mathcal P_F\longrightarrow\mathcal P_{F'}
\]
denote the induced ambient Artinian-slice transport.  Explicitly,
$f_{\mathcal P,!}$ is left Kan transport along the coherent Artinian transition,
with its structural map to $F'$ obtained from the canonical formal-moduli transport
map.  Let
\[
f_{\mathrm{FMP},!}:
\bigl(\operatorname{FMP}^{\mathrm{sp}}_{B/R}\bigr)/F
\longrightarrow
\bigl(\operatorname{FMP}^{\mathrm{sp}}_{B'/R'}\bigr)/F'
\]
be the induced transport on the Schlessinger-local slices.  Then there is a
canonical exchange equivalence
\[
\boxed{
L_{F'}^{\mathrm{Sch}}\,f_{\mathcal P,!}
\simeq
f_{\mathrm{FMP},!}\,L_F^{\mathrm{Sch}}.
}
\]
These equivalences satisfy identity, composition, and higher-pasting coherence.

In particular, whenever the Artinian complete-symmetry presheaf and all of its
formalized input data are transported along $f$, the reflected complete-symmetry
parameter is transported by a canonical equivalence
\[
f_{\mathrm{FMP},!}
\left(
\widehat{\operatorname{Sym}}_{\mathcal G,L}^{\mathrm{Sch},q}
\right)
\simeq
\widehat{\operatorname{Sym}}_{\mathcal G',L'}^{\mathrm{Sch},q},
\]
where primes denote the transported formalized symmetry and Lagrangian data.
\end{proposition}

\begin{proof}
The coherent Artinian transport of Chapter 14 preserves the identity Artinian
object and every structured square-zero Schlessinger pullback.  Hence its
right-adjoint restriction functor preserves ordinary Schlessinger-local presheaves.

Let $R_{\mathcal P}$ denote that right adjoint on ambient presheaves.  The right
adjoint to the induced slice transport sends a morphism $Y\to F'$ to
\[
R_{\mathcal P}Y
\times_{R_{\mathcal P}F'}
F
\longrightarrow
F,
\]
where $F\to R_{\mathcal P}F'$ is adjoint to the canonical transport map
$f_{\mathcal P,!}F\to F'$.  If $Y$ is Schlessinger-local, then
$R_{\mathcal P}Y$ and $R_{\mathcal P}F'$ are Schlessinger-local; the
Schlessinger-local subcategory is closed under limits, so the displayed pullback is
again local.  Thus the right adjoint to $f_{\mathcal P,!}$ on slices preserves the
local objects identified above.

The universal property of accessible localization therefore induces the left
adjoint $f_{\mathrm{FMP},!}$ on local slices and identifies localization after
ambient transport with local transport after localization.  This is the displayed
exchange equivalence.

One can see the same comparison directly on the chosen generating families.
Artinian transport carries the identity generator and every structured
square-zero generator to the corresponding generators after coefficient change;
a map from the target of such a generator to $F$ is transported to the
corresponding map to $F'$.  Thus the relative generating family
$W_F^{\mathrm{Sch}}$ is carried to the relative generating family
$W_{F'}^{\mathrm{Sch}}$.  Contractible uniqueness of the localization
factorizations, together with the identity, composition, and higher coherence of
Artinian transport, supplies the asserted higher-pasting coherence.

The final equivalence is the exchange formula applied to the transported
Artinian complete-symmetry presheaf.
\end{proof}

\begin{construction}[Artinian complete-symmetry presheaf over the formal leaf]
\label{constr:symmetry-noether-artinian-complete-symmetry-presheaf-over-the-formal-leaf}
For an Artinian arrow $A\to B$ and a point $x\in F(A)$, restrict the complete
quotient-formalized symmetry problem to the canonical Artinian \v{C}ech test
determined by $x$.  Denote its space of complete stage-$q$ symmetry realizations by
\[
\mathcal S_{\mathcal G,L}^{\mathrm{Art},q}(A,x).
\]
Restriction of Artinian tests, the formalized relation maps, horizontal quotient,
symmetry partial totalizations, and their inverse limit make these spaces functorial in
$(A,x)$.  Unstraightening therefore gives an object
\[
\mathcal S_{\mathcal G,L}^{\mathrm{Art},q}
\longrightarrow F
\]
of $\mathcal P_F$.

Define the Schlessinger-reflected complete-symmetry parameter by
\[
\boxed{
\widehat{\operatorname{Sym}}_{\mathcal G,L}^{\mathrm{Sch},q}
:=
L_F^{\mathrm{Sch}}
\left(
\mathcal S_{\mathcal G,L}^{\mathrm{Art},q}
\right)
\longrightarrow F .
}
\]
This object is a spectral formal-moduli problem over $F$.  It is not identified
with the native object
$\operatorname{Sym}_{\mathcal G}^{\mathrm{fm},q}(L)$; the former is obtained by
Artinian restriction and Schlessinger reflection of the latter's formalized
symmetry problem on the fixed leaf.
\end{construction}

\begin{proposition}[Compatibility with the formalized first obstruction]
\label{prop:symmetry-noether-compatibility-with-the-formalized-first-obstruction}
The universal complete symmetry over
$\mathcal S_{\mathcal G,L}^{\mathrm{Art},q}$ carries a specified nullhomotopy of
the Artinian first Noether obstruction.  After application of the stable
localization unit
\[
\mathcal O_{\mathcal G,L}^{\mathrm{Art},1,q}
\longrightarrow
\widehat{\mathcal O}_{\mathcal G,L}^{1,q},
\]
this gives a specified nullhomotopy of the pullback of
$\widehat{\mathbf n}_{\mathcal G,L}^{1,q}$ to
$\mathcal S_{\mathcal G,L}^{\mathrm{Art},q}$.
\end{proposition}

\begin{proof}
At every Artinian point, a complete formalized symmetry is by definition a lift
through the full symmetry totalization.  Its first partial-totalization component is
therefore a specified nullhomotopy of the formalized first obstruction.  These
nullhomotopies are natural in Artinian restriction.  Applying
$\Omega_F^\infty$ to the stable Schlessinger localization unit transports them to
the reflected coefficient, preserving all higher homotopies.
\end{proof}

\begin{definition}[Differentiated Noether representation]
\label{def:symmetry-noether-differentiated-noether-representation}
Put
\[
\mathfrak a_{\mathcal G}
:=
\operatorname{Diff}^{\pi,E_\infty}_{B/R}(F)
=
\mathbb A_{\mathcal G}^{\pi,E_\infty}|_{U/V}
\]
and
\[
\boxed{
\mathbb O_{\mathcal G,L}^{1,q}
:=
D_F^{\mathrm{lin}}
\left(
\widehat{\mathcal O}_{\mathcal G,L}^{1,q}
\right)
\in
\operatorname{Rep}^{\pi,E_\infty}(\mathfrak a_{\mathcal G}).
}
\]
\end{definition}

\begin{construction}[Retractive algebroid of the reflected Noether coefficient]
\label{constr:symmetry-noether-retractive-algebroid-of-the-reflected-noether-coefficient}
Form the relative infinite-loop object
\[
Y_{\mathcal G,L}^{1,q}
:=
\Omega_F^\infty
\widehat{\mathcal O}_{\mathcal G,L}^{1,q}[n]
\longrightarrow
F.
\]
It has the zero section
\[
z:F\longrightarrow Y_{\mathcal G,L}^{1,q}
\]
and the reflected Noether section
\[
\widehat{\mathbf n}_{\mathcal G,L}^{1,q}:
F\longrightarrow Y_{\mathcal G,L}^{1,q}.
\]
Differentiate the projection:
\[
\mathfrak b_{\mathcal G,L}^{1,q}
:=
\operatorname{Diff}^{\pi,E_\infty}_{B/R}
\left(
Y_{\mathcal G,L}^{1,q}
\right)
\longrightarrow
\mathfrak a_{\mathcal G}.
\]
Differentiating the two sections gives two splittings
\[
z_{\mathfrak a},\ s_{\mathfrak n}:
\mathfrak a_{\mathcal G}
\rightrightarrows
\mathfrak b_{\mathcal G,L}^{1,q}
\]
of this morphism of spectral partition-Lie algebroids.
\end{construction}

\begin{construction}[Reflected symmetry zero locus and universal section homotopy]
\label{constr:symmetry-noether-reflected-symmetry-zero-locus-and-universal-section-homotopy}
The stable Schlessinger theorem identifies every shifted relative infinite-loop
object of a local spectrum object with a Schlessinger-local object in the slice.
Hence
\[
Y_{\mathcal G,L}^{1,q}\longrightarrow F
\]
is an object of
$\bigl(\operatorname{FMP}^{\mathrm{sp}}_{B/R}\bigr)/F$.

Define the zero locus of the reflected Noether section in that slice by
\[
Z_F\!\left(\widehat{\mathbf n}_{\mathcal G,L}^{1,q}\right)
:=
F
\times_{Y_{\mathcal G,L}^{1,q}}
F,
\]
where the two maps are
$\widehat{\mathbf n}_{\mathcal G,L}^{1,q}$ and the zero section $z$.

The Artinian symmetry nullhomotopy of the preceding proposition gives a morphism
\[
\mathcal S_{\mathcal G,L}^{\mathrm{Art},q}
\longrightarrow
Z_F\!\left(\widehat{\mathbf n}_{\mathcal G,L}^{1,q}\right)
\]
over $F$.  Since the target is already Schlessinger-local, the universal property
of $L_F^{\mathrm{Sch}}$ gives a contractibly unique factorization
\[
\boxed{
\widehat{\operatorname{Sym}}_{\mathcal G,L}^{\mathrm{Sch},q}
\longrightarrow
Z_F\!\left(\widehat{\mathbf n}_{\mathcal G,L}^{1,q}\right).
}
\]
Equivalently, over
$\widehat{\operatorname{Sym}}_{\mathcal G,L}^{\mathrm{Sch},q}$ there is a
canonical specified homotopy
\[
\boxed{
\widehat{\mathbf n}_{\mathcal G,L}^{1,q}\simeq z.
}
\]
\end{construction}

\begin{proof}
The zero locus is a finite limit of spectral formal-moduli problems over $F$, so it
is again a spectral formal-moduli problem over $F$.  The preceding proposition
supplies the morphism from the unreflected Artinian symmetry presheaf.  Localization
adjunction then gives the displayed factorization and its contractible uniqueness.
The mapping-space interpretation of the homotopy pullback is exactly the stated
section homotopy.
\end{proof}

\begin{theorem}[Differentiated retractive Noether algebroid]
\label{thm:symmetry-noether-differentiated-retractive-noether-algebroid}
The stable fibre of
\[
\mathfrak b_{\mathcal G,L}^{1,q}
\longrightarrow
\mathfrak a_{\mathcal G}
\]
under the fibre-at-the-base equivalence is canonically
\[
U\mathbb O_{\mathcal G,L}^{1,q}[n].
\]
The two splittings therefore determine a canonical vertical difference
\[
\boxed{
d_0\mathbf n_{\mathcal G,L}^{1,q}:
U\mathfrak a_{\mathcal G}
\longrightarrow
U\mathbb O_{\mathcal G,L}^{1,q}[n].
}
\]
For every structured square-zero test, evaluation of this map is the reduced first
difference of the reflected Noether section.  Consequently the construction agrees
with the square-zero Yoneda derivative and is independent of an Artinian
presentation.
\end{theorem}

\begin{proof}
The tangent differentiation equivalence is obtained by stabilizing the slice over
$F$.  The fibre-at-$F$ equivalence identifies the tangent fibre of the relative
infinite-loop object of a representation with its underlying representation
coefficient.  Hence the fibre of the differentiated projection is
$U\mathbb O_{\mathcal G,L}^{1,q}[n]$.

The difference between the two splittings is vertical because their composites
with $\mathfrak b_{\mathcal G,L}^{1,q}\to\mathfrak a_{\mathcal G}$ are both the
identity.  This gives the displayed stable map.

Now evaluate on a structured square-zero extension $B_\alpha\to B$.  Under the
affine formal-moduli/partition-Lie equivalence, this test is represented by the
corresponding free anchored partition-Lie cell.  Full faithfulness identifies the
difference of the two differentiated splittings on that cell with
\[
\widehat{\mathbf n}_{\mathcal G,L}^{1,q}(B_\alpha)
-
\widehat{\mathbf n}_{\mathcal G,L}^{1,q}(B).
\]
This is exactly the reduced first difference.  Free square-zero cells detect the
map, giving the intrinsic characterization.
\end{proof}

\begin{proposition}[Spectral-\'etale descent]
\label{prop:symmetry-noether-spectral-etale-descent}
The reflected complete-symmetry parameter
$\widehat{\operatorname{Sym}}_{\mathcal G,L}^{\mathrm{Sch},q}$, the coefficients
$\widehat{\mathcal O}_{\mathcal G,L}^{1,q}$ and
$\mathbb O_{\mathcal G,L}^{1,q}$, the retractive algebroids
$\mathfrak b_{\mathcal G,L}^{1,q}$, both splittings, and
$d_0\mathbf n_{\mathcal G,L}^{1,q}$ are compatible with the coherent
spectral-\'etale transition equivalences of the formalized coefficient system and
therefore descend on the global theorem domain.
\end{proposition}

\begin{proof}
The unstable reflected symmetry parameter is transported by
Proposition~\ref{prop:symmetry-noether-spectral-etale-exchange-for-relative-unstable-schlessinger-reflector}.
Stable Artinian sections and the stable Schlessinger reflector carry their
independently constructed spectral-\'etale exchange equivalences.  Relative
infinite-loop objects, spectral formal-moduli differentiation, and the
fibre-at-the-base equivalence carry the two-variable spectral-\'etale
Beck--Chevalley equivalences constructed in the preceding chapters.  The zero and
Noether sections and their difference are natural for those equivalences.
\end{proof}

\begin{convention}[Global descended notation]
\label{conv:symmetry-noether-global-descended-notation}
On the global locally coherent theorem domain, use the same symbols for the
descended representation and retractive algebroid.  In particular write
\[
\mathbb O_{\mathcal G,L}^{1,q}
\in
\operatorname{Rep}^{\pi,E_\infty}
\left(
\mathbb A_{\mathcal G}^{\pi,E_\infty}
\right)
\]
and
\[
\boxed{
d_0\mathbf n_{\mathcal G,L}^{1,q}:
U\mathbb A_{\mathcal G}^{\pi,E_\infty}
\longrightarrow
U\mathbb O_{\mathcal G,L}^{1,q}[n].
}
\]
The retractive algebroid and its two splittings are likewise understood as the
descended cocartesian family on the global coefficient system.  All subsequent
global syzygy and Second Noether statements use this descended notation; their
affine restrictions are the constructions above.
\end{convention}

\section{Directional Taylor classes and higher differentiated Noether coherence}
\label{sec:symmetry-noether-directional-taylor-classes-and-higher-differentiated-noether-coherence}

The directional multi-filtered square-zero geometry of Chapter 15 and the
cotangent-comonadic primitive operation system remain different constructions.
The exact labelled Taylor class is defined first from the differentiated
partition-Lie operation object and the actual tangent variables of the controller.
The directional Artinian branch then reaches that same exact operation only through
the separately constructed all-positive-to-exact scalar diagonal, comonadic
Hurewicz comparison, primitive cotangent totalization, and continuous duality.
In particular, the Boolean cross-effect of the values of the reflected Noether
section is not identified with a higher partition-Lie corolla.

Let $I$ be finite and nonempty.  Choose connective almost-perfect derivations
\[
\alpha_i:L_{B/R}[-1]\longrightarrow M_i,
\qquad i\in I,
\]
and put
\[
P_i:=M_i^\vee,
\qquad
\bar\alpha_i:=\alpha_i^\vee:
P_i\longrightarrow T^{pc}_{B/R}[1].
\]
Let $A_I\to B$ be the corresponding multi-filtered square-zero test of Chapter 15,
and for $J\subseteq I$ let $A_J$ be obtained by switching off the directions
outside $J$.

\begin{definition}[Actual labelled tangent variables]
\label{def:symmetry-noether-actual-labelled-tangent-variables}
An actual labelled tangent variable for $(M_i,\alpha_i)$ is a morphism of anchored
pro-coherent objects
\[
\boxed{
\xi_i:
\left(
P_i\xrightarrow{\bar\alpha_i}T^{pc}_{B/R}[1]
\right)
\longrightarrow
\left(
U\mathfrak a_{\mathcal G}
\xrightarrow{\rho_{\mathcal G}^{\pi,E_\infty}}
T^{pc}_{B/R}[1]
\right).
}
\]
A labelled tangent tuple is a family
$((M_i,\alpha_i,\xi_i))_{i\in I}$.
The derivation $(M_i,\alpha_i)$ specifies the anchored square-zero probe; the map
$\xi_i$ specifies its actual value in the controller.  No tangent value is attached
to the derivation alone.
\end{definition}

\begin{construction}[Compatible formal-leaf points of a labelled tangent tuple]
\label{constr:symmetry-noether-compatible-formal-leaf-points-of-a-labelled-tangent-tuple}
For $J\subseteq I$, let
\[
A_J^{0}:=\operatorname{ev}_{\mathbf0}(A_J)\longrightarrow B
\]
be the underlying ordinary structured square-zero test obtained by evaluation at
multiweight zero.  The finite multi-filtered cotangent--square-zero adjunction
identifies $A_J^{0}$ with the ordinary square-zero test classified by
\[
\alpha_J:
L_{B/R}[-1]
\longrightarrow
\bigoplus_{j\in J}M_j.
\]
Finite continuous duality gives
\[
P_J:=\left(\bigoplus_{j\in J}M_j\right)^\vee
\simeq
\bigoplus_{j\in J}P_j,
\]
and the tuple $(\xi_j)_{j\in J}$ induces an anchored morphism
\[
\xi_J:
\left(
P_J\xrightarrow{\bar\alpha_J}T^{pc}_{B/R}[1]
\right)
\longrightarrow
\left(
U\mathfrak a_{\mathcal G}
\xrightarrow{\rho_{\mathcal G}^{\pi,E_\infty}}
T^{pc}_{B/R}[1]
\right).
\]
The complete-almost-finite extension of the affine formal-moduli/partition-Lie
recognition theorem gives
\[
F^{\wedge\mathrm{aft}}(A_J^{0})
\simeq
\operatorname{Map}_{\operatorname{Anch}^{pc}_{B/R}}
\left(
P_J\xrightarrow{\bar\alpha_J}T^{pc}_{B/R}[1],
U\mathfrak a_{\mathcal G}
\xrightarrow{\rho_{\mathcal G}^{\pi,E_\infty}}
T^{pc}_{B/R}[1]
\right).
\]
Define
\[
\boxed{
x_J\in F^{\wedge\mathrm{aft}}(A_J^{0})
}
\]
to be the point corresponding to $\xi_J$.  When $A_J^{0}\to B$ is Artinian this is
the ordinary point of $F(A_J^{0})$.  For $J\subseteq K$, restriction along the
coordinate projection carries $x_K$ to $x_J$.  Thus the individual tangent
variables canonically determine the full directional cube of formal-leaf points;
no additional coherent $I$-point is retained as input.
\end{construction}

\begin{proof}
The affine formal-moduli/partition-Lie equivalence sends a structured square-zero
extension classified by $\alpha_J$ to the free anchored partition-Lie cell on
$P_J\xrightarrow{\bar\alpha_J}T^{pc}_{B/R}[1]$.  Its complete-almost-finite
extension therefore identifies evaluation of $F$ on $A_J^{0}$ with the displayed
anchored mapping space.  Colimits in the anchored slice are created in the
underlying pro-coherent category, so the anchored direct-sum probe is the coproduct
of the labelled probes and the maps $\xi_j$ assemble uniquely to $\xi_J$.
Coordinate projection $A_K^{0}\to A_J^{0}$ is Koszul-dual to inclusion of the
corresponding free anchored summand; precomposition therefore carries $\xi_K$ to
$\xi_J$.  The identity, composition, and all higher compatibilities are those of
the square-zero/free-cell equivalence.
\end{proof}

\begin{construction}[Exact labelled Taylor class of the Noether splitting]
\label{constr:symmetry-noether-exact-labelled-taylor-class-of-the-noether-splitting}
For a labelled tangent tuple $((M_i,\alpha_i,\xi_i))_{i\in I}$, let
\[
\mathcal C_I(\alpha_\bullet)
:=
\operatorname{Op}^{\pi,E_\infty}_{B/R}
\left(
I;
\left(
P_i\xrightarrow{\bar\alpha_i}T^{pc}_{B/R}[1]
\right)_{i\in I}
\right)
\]
be the exact labelled $I$-corolla operation object of Chapter 15.  The
partition-Lie algebroid structure of $\mathfrak a_{\mathcal G}$ and the actual input
maps $\xi_i$ give the canonical corolla evaluation
\[
\mu^{I}_{\mathfrak a,\xi}:
\mathcal C_I(\alpha_\bullet)
\longrightarrow
U\mathfrak a_{\mathcal G}.
\]
Define the exact labelled Noether Taylor morphism by
\[
\boxed{
\operatorname{Tay}_{\mathcal G,L}^{I,q}
\left((M_i,\alpha_i,\xi_i)_{i\in I}\right)
:=
 d_0\mathbf n_{\mathcal G,L}^{1,q}
 \circ
 \mu^{I}_{\mathfrak a,\xi}:
\mathcal C_I(\alpha_\bullet)
\longrightarrow
U\mathbb O_{\mathcal G,L}^{1,q}[n].
}
\]
Equivalently, evaluate the two actual splittings
\[
s_{\mathfrak n},z_{\mathfrak a}:
\mathfrak a_{\mathcal G}
\rightrightarrows
\mathfrak b_{\mathcal G,L}^{1,q}
\]
on the same exact labelled source corolla.  Their difference has zero projection to
$U\mathfrak a_{\mathcal G}$ and therefore factors uniquely through the stable
vertical fibre $U\mathbb O_{\mathcal G,L}^{1,q}[n]$; under the fibre-at-the-base
identification that vertical factor is exactly the displayed Taylor morphism.

The construction is natural in the labelled derivations and actual tangent
variables and equivariant for bijections of $I$.  For $|I|=1$, the exact unary
corolla is the identity operation, so the Taylor morphism is precisely evaluation
of $d_0\mathbf n_{\mathcal G,L}^{1,q}$ on the chosen tangent variable.
\end{construction}

\begin{construction}[Directional Artinian exact-cell extraction]
\label{constr:symmetry-noether-directional-artinian-cross-effect-and-exact-multilinearization}
Retain the multi-filtered derivation
\[
D_I:
L_{B/R}[-1]
\longrightarrow
\bigoplus_{i\in I}M_i(e_i)
\]
and the directional square-zero cube $J\mapsto A_J$ of Chapter 15.  Its
source-fixed directional \v{C}ech primitive is the Boolean total fibre obtained by
switching the square-zero directions themselves on and off.  The directional
cross-effect theorem identifies that primitive with the all-positive filtration
stage.  The constructed directional meeting map then passes from the all-positive
degree-zero primitive to the exact scalar diagonal, the comonadic scalar-diagonal
Hurewicz map passes to the primitive cotangent cobar, and totalization gives the
exact multiweight co-operation object
\[
\operatorname{coOp}^{\mathrm{Sp}}_{B/R}
\left(I;(M_i,\alpha_i)_{i\in I}\right).
\]
Continuous pro-coherent duality identifies its dual with
$\mathcal C_I(\alpha_\bullet)$.  Evaluating the resulting exact operation on the
actual tangent variables $\xi_i$ and postcomposing with
$d_0\mathbf n_{\mathcal G,L}^{1,q}$ defines
\[
\boxed{
\operatorname{Tay}^{I,q}_{\mathcal G,L,\mathrm{dir}}
\left((M_i,\alpha_i,\xi_i)_{i\in I}\right):
\mathcal C_I(\alpha_\bullet)
\longrightarrow
U\mathbb O_{\mathcal G,L}^{1,q}[n].
}
\]
When the underlying tests $A_J^{0}$ lie in the spectral Artinian subcategory, the
compatible points $x_J$ constructed above allow the reflected Noether section to be
evaluated on the same directional square-zero tests.  Those section values form a
separate raw directional cube.  Its Boolean total fibre is retained as a raw
section-value cross-effect, but no equivalence of that cross-effect with
$\mathcal C_I(\alpha_\bullet)$ is asserted.  Outside the Artinian range no such
section-value cube is used in the Taylor construction.
\end{construction}

\begin{theorem}[Directional recognition of the exact Taylor class]
\label{thm:symmetry-noether-directional-recognition-of-the-exact-taylor-class}
For every labelled tangent tuple there is a canonical equivalence of morphisms
\[
\boxed{
\operatorname{Tay}^{I,q}_{\mathcal G,L,\mathrm{dir}}
\left((M_i,\alpha_i,\xi_i)_{i\in I}\right)
\simeq
\operatorname{Tay}_{\mathcal G,L}^{I,q}
\left((M_i,\alpha_i,\xi_i)_{i\in I}\right).
}
\]
Thus the directional Artinian branch recovers the exact labelled Noether Taylor
morphism only after the constructed exact scalar-diagonal and cotangent-operation
extraction.  It is not the raw Boolean cross-effect of the values of
$\widehat{\mathbf n}_{\mathcal G,L}^{1,q}$.
\end{theorem}

\begin{proof}
The multi-filtered diagonal primitive theorem identifies the exact $1_I$ cotangent
primitive of $A_I$ with
$\operatorname{coOp}^{\mathrm{Sp}}_{B/R}(I;(M_i,\alpha_i))$, compatibly with
symmetric relabelling and substitution.  Continuous duality is exactly the
construction of the labelled operation object
$\mathcal C_I(\alpha_\bullet)$.  Supplying the actual tangent variables $\xi_i$
then gives the controller corolla evaluation
$\mu^I_{\mathfrak a,\xi}$ on both sides of that identification.  Postcomposition
with the already constructed stable map
$d_0\mathbf n_{\mathcal G,L}^{1,q}$ therefore gives the same morphism to
$U\mathbb O_{\mathcal G,L}^{1,q}[n]$.  This is the displayed equivalence, with all
higher coherence inherited from the multi-filtered comonad and its dual
partition-Lie substitution system.
\end{proof}

\begin{construction}[Taylor--operation comparison on the scalar Artinian overlap]
\label{constr:symmetry-noether-taylor-operation-comparison-on-the-scalar-artinian-overlap}
On the scalar Artinian coefficient range of the Atiyah--Spencer operation
calculus, the directional branch used above factors through the explicit comparison
ladder.  Apply the directional meeting map from the all-positive degree-zero
\v{C}ech primitive to the exact scalar diagonal and then the comonadic Hurewicz
morphism
\[
\operatorname{ScCotBar}_{B/R}(I;D_I)
\longrightarrow
\operatorname{CotBar}_{B/R}(I;D_I).
\]
After totalization, the primitive cotangent comparison identifies the target with
the exact multiweight co-operation object, and continuous duality gives the
corresponding partition-Lie operation object.  Evaluation on
$(\xi_i)_{i\in I}$ followed by
$d_0\mathbf n_{\mathcal G,L}^{1,q}$ is the Taylor morphism constructed above.
This is the canonical scalar directional comparison used here; no direct
identification with the raw cube of reflected Noether-section values is made.
\end{construction}

\begin{proposition}[Non-conflation of directional and cotangent operation systems]
\label{prop:symmetry-noether-non-conflation-of-directional-and-cotangent-operation-systems}
The preceding comparison does not identify the support-marked geometric
\v{C}ech system, the directional Artinian square-zero cube, the raw cube of
reflected Noether-section values, or the cotangent-comonadic cosimplicial cobar.
The directional and cotangent cosimplicial directions meet at the constructed exact
scalar diagonal, and the exact labelled partition-Lie corolla is reached only after
the stated Hurewicz, totalization, and continuous-duality steps.
\end{proposition}

\subsection{The morphism corolla--latching hierarchy}
\label{subsec:symmetry-noether-the-morphism-corolla-latching-hierarchy}

Put
\[
h_{\mathfrak n}:=Us_{\mathfrak n}:
U\mathfrak a_{\mathcal G}
\longrightarrow
U\mathfrak b_{\mathcal G,L}^{1,q}.
\]
This is the underlying anchored morphism of an actual spectral partition-Lie
algebroid morphism.

\begin{construction}[Noether morphism latching objects]
\label{constr:symmetry-noether-noether-morphism-latching-objects}
Apply the parallel morphism obstruction construction of the Atiyah--Spencer
chapter to $h_{\mathfrak n}$.  For every unmarked weight $r\ge1$, write
\[
L_r^{\mathrm{Noeth}},
\qquad
C_r^{\mathrm{Noeth}},
\qquad
Q_r^{\mathrm{Noeth}}
:=
\operatorname{cofib}
\left(
L_r^{\mathrm{Noeth}}
\longrightarrow
C_r^{\mathrm{Noeth}}
\right),
\]
and
\[
R_r^{\mathrm{Noeth}}
:=
Q_r^{\mathrm{Noeth}}[-1].
\]
Here $C_r^{\mathrm{Noeth}}$ is the marked morphism corolla of total unmarked weight
$r$, while $L_r^{\mathrm{Noeth}}$ is the colimit of all proper reduced
morphism trees obtained from lower-weight coefficient actions, the fixed source
and target algebroid structures, and the already determined lower-weight
components of $h_{\mathfrak n}$.
\end{construction}

\begin{theorem}[Higher differentiated Noether coherence]
\label{thm:symmetry-noether-higher-differentiated-noether-coherence}
For every $r\ge1$, the already constructed lower-weight data determine a canonical
morphism
\[
g_r^{\mathrm{Noeth}}:
L_r^{\mathrm{Noeth}}
\longrightarrow
U\mathfrak b_{\mathcal G,L}^{1,q}.
\]
The extension obstruction is the connecting morphism
\[
\operatorname{ob}^{\mathrm{Noeth}}_r:
R_r^{\mathrm{Noeth}}
\longrightarrow
U\mathfrak b_{\mathcal G,L}^{1,q}.
\]
Because $s_{\mathfrak n}$ is an actual morphism of spectral partition-Lie
algebroids, its weight-$r$ corolla component supplies a canonical filler and hence
a canonical nullhomotopy
\[
\boxed{
\operatorname{ob}^{\mathrm{Noeth}}_r\simeq0.
}
\]
The complete family of these fillers and nullhomotopies is equivalent to the full
partition-Lie morphism coherence of $s_{\mathfrak n}$.
\end{theorem}

\begin{proof}
This is the morphism corolla--latching extension theorem applied to the actual
algebroid morphism $s_{\mathfrak n}$.  Every proper reduced weight-$r$ morphism
tree is evaluable from lower-weight morphism components together with the fixed
source and target algebra structures.  These evaluations assemble to
$g_r^{\mathrm{Noeth}}$.  The cofiber sequence
\[
L_r^{\mathrm{Noeth}}
\longrightarrow
C_r^{\mathrm{Noeth}}
\longrightarrow
Q_r^{\mathrm{Noeth}}
\]
identifies an extension with a specified nullhomotopy of the connecting map from
$R_r^{\mathrm{Noeth}}$.  The actual weight-$r$ component of
$s_{\mathfrak n}$ is such an extension.  Compatibility in all weights is exactly
the monad-morphism coherence already carried by the differentiated splitting.
\end{proof}

\begin{proposition}[Relation with the Taylor classes]
\label{prop:symmetry-noether-relation-with-the-taylor-classes}
For every finite nonempty label set $I$ and every labelled tangent tuple
$((M_i,\alpha_i,\xi_i))_{i\in I}$, evaluation of the indecomposable labelled source
corolla under the vertical difference
$s_{\mathfrak n}-z_{\mathfrak a}$ is precisely
\[
\operatorname{Tay}_{\mathcal G,L}^{I,q}
\left((M_i,\alpha_i,\xi_i)_{i\in I}\right).
\]
The full morphism-latching relation at that weight is not the vanishing of this
corolla alone: it compares that corolla with all proper decompositions built from
lower Taylor components and the partition-Lie operations of the source and target.
\end{proposition}

\begin{proof}
The first assertion is the vertical-factor characterization in
Construction~\ref{constr:symmetry-noether-exact-labelled-taylor-class-of-the-noether-splitting}.
The second is the defining latching relation for a morphism of algebras over the
partition-Lie monad.
\end{proof}

\section{Zero-action parameters, actual isotropy, and first Noether syzygies}
\label{sec:symmetry-noether-zero-action-parameters-actual-isotropy-and-first-noether-syzygies}

\begin{definition}[First Noether-syzygy coefficient]
\label{def:symmetry-noether-first-noether-syzygy-coefficient}
Define
\[
\boxed{
\operatorname{Syz}_{\mathcal G,L}^{1,q}
:=
\operatorname{fib}
\left(
d_0\mathbf n_{\mathcal G,L}^{1,q}:
U\mathbb A_{\mathcal G}^{\pi,E_\infty}
\longrightarrow
U\mathbb O_{\mathcal G,L}^{1,q}[n]
\right).
}
\]
This is the first differentiated syzygy coefficient.  Higher differentiated
coherence is measured by the morphism latching objects
$R_r^{\mathrm{Noeth}}$ and their obstruction maps, not by declaring the fibres of
the higher directional Taylor classes to be higher syzygies.
\end{definition}

\begin{lemma}[Zero-action cells canonically nullify the first Noether derivative]
\label{lem:symmetry-noether-the-first-noether-derivative-factors-through-the-field-action}
Let
\[
i_F:
UF_{\mathcal G}^{\pi,E_\infty}
\longrightarrow
U\mathbb A_{\mathcal G}^{\pi,E_\infty}
\]
be the underlying map from the homotopy anchor fibre.  Work on a coherent affine
spectral chart $U=\operatorname{Spec}B$.  For every connective coherent
$B$-module $M$ and every $k<0$, put
\[
P(M,k):=M^\vee[k].
\]
For every morphism
\[
u:P(M,k)\longrightarrow UF_{\mathcal G}^{\pi,E_\infty},
\]
the composite
\[
P(M,k)
\xrightarrow{\ u\ }
UF_{\mathcal G}^{\pi,E_\infty}
\xrightarrow{\ i_F\ }
U\mathbb A_{\mathcal G}^{\pi,E_\infty}
\xrightarrow{\ d_0\mathbf n_{\mathcal G,L}^{1,q}\ }
U\mathbb O_{\mathcal G,L}^{1,q}[n]
\]
carries a canonical nullhomotopy.  These nullhomotopies are natural in maps of
such cells and in all higher simplices.  After passage to the stabilized tangent
slice they are compatible with shifts, finite fibres and cofibres, and retracts.
Consequently they extend coherently to the idempotent-complete finite stable
closure of the cells $P(M,k)$.  The entire construction is compatible with
coherent spectral-\'etale transport.
\end{lemma}

\begin{proof}
By the zero-anchor free--forgetful adjunction, a map
\[
u:P(M,k)\longrightarrow UF_{\mathcal G}^{\pi,E_\infty}
\]
determines a morphism from the free zero-anchor spectral partition-Lie algebroid
on
\[
P(M,k)\xrightarrow{0}T_{B/R}^{pc}[1]
\]
to $F_{\mathcal G}^{\pi,E_\infty}$ and hence to
$\mathbb A_{\mathcal G}^{\pi,E_\infty}$.  For $k<0$, the negative-cell
calculation of the differentiation chapter identifies this free cell with the
Koszul-dual image of a split structured square-zero Artinian extension.  More
explicitly, with
\[
M':=M[-k-1],
\]
which is again connective and coherent, one has
\[
\operatorname{Free}_{B/R}^{\pi,E_\infty}
\left(M^\vee[k]\xrightarrow0T_{B/R}^{pc}[1]\right)
\simeq
D_{B/R}^{\mathrm{sp}}
\left(B\oplus M'[1]\longrightarrow B\right).
\]
Thus the chosen map is a structured square-zero infinitesimal symmetry parameter.
Because it factors through the homotopy anchor fibre, it carries the tautological
nullhomotopy of its induced first-order field action.

The reflected first Noether section is obtained from the formalized normalized
source--target difference of the Lagrangian by horizontal quotient, stable
Schlessinger reflection, and relative infinite-loop formation.  Evaluation of
$d_0\mathbf n_{\mathcal G,L}^{1,q}$ on the preceding square-zero tangent cell is
therefore the reduced first difference of that reflected Noether section.  The
tautological zero-action homotopy identifies the two induced first-order field
maps.  Pulling the normalized Lagrangian difference through this homotopy gives a
specified nullhomotopy of the reduced first difference, hence of the displayed
composite.

This construction occurs before forgetting the spectrum-object tangent structure.
Stable Schlessinger localization and tangent differentiation are exact on that
stable structure, so the nullhomotopies commute with suspension and desuspension,
finite fibre/cofiber constructions, and retracts.  They therefore extend over the
finite stable idempotent closure of the negative square-zero cells.  Naturality of
the square-zero/free-cell equivalence gives naturality in cell maps and all higher
simplices, and the two-variable spectral-\'etale Beck--Chevalley equivalences
transport the construction between coherent charts.
\end{proof}

\begin{lemma}[Pro-coherent generator detection of coherent nullhomotopies]
\label{lem:symmetry-noether-pro-coherent-generator-detection-of-coherent-nullhomotopies}
Let $B$ be coherent and define
\[
\boxed{
\mathcal P_B
:=
\operatorname{Thick}^{\natural}
\left\{
M^\vee[k]
\ \middle|\
M\in\operatorname{Coh}(B),\ M\text{connective},\ k<0
\right\}
\subset QC_B^\vee.
}
\]
Then $\mathcal P_B$ is a small stable idempotent-complete subcategory of compact
generators and restricted spectral Yoneda
\[
y_{\mathcal P_B}:
QC_B^\vee
\longrightarrow
\operatorname{Fun}(\mathcal P_B^{\mathrm{op}},\operatorname{Sp})
\]
is fully faithful.

Let $f:X\to Y$ be a morphism in $QC_B^\vee$.  Suppose that for every generating
negative square-zero cell $P(M,k)$ and every map $u:P(M,k)\to X$ there is a
specified nullhomotopy $fu\simeq0$, natural in maps of cells and all higher
simplices, and that this family is compatible with shifts, finite stable
fibre/cofiber constructions, and retracts.  Then these data determine a canonical
coherent nullhomotopy
\[
\boxed{f\simeq0.}
\]
The space of nullhomotopies inducing the prescribed cellwise family is
contractible.
\end{lemma}

\begin{proof}
The pro-coherent category is
\[
QC_B^\vee\simeq
\operatorname{Ind}\bigl(\operatorname{Coh}(B)^{\mathrm{op}}\bigr),
\]
and the continuous duals of coherent modules are compact generators.  The
subcategory $\mathcal P_B$ contains that compact generating family.  Indeed, for
$N\in\operatorname{Coh}(B)$ choose $r$ such that $N[r]$ is connective.  Then
\[
(N[r])^\vee[-1]
\simeq
N^\vee[-r-1]
\]
is one of the displayed negative cells, and stability of $\mathcal P_B$ implies
$N^\vee\in\mathcal P_B$.  Hence $\mathcal P_B$ is a compact generating stable
subcategory, and restricted spectral Yoneda on it is fully faithful.

The two functors
\[
P\longmapsto\operatorname{map}_{QC_B^\vee}(P,X),
\qquad
P\longmapsto\operatorname{map}_{QC_B^\vee}(P,Y)
\]
are exact spectral functors on $\mathcal P_B^{\mathrm{op}}$.  By the assumed
stable compatibility, the cellwise nullhomotopies extend over the finite stable
and idempotent completion to a homotopy of natural transformations
\[
y_{\mathcal P_B}(f)\simeq0.
\]
Full faithfulness identifies the path space between $f$ and the zero morphism
with the corresponding path space between their restricted-Yoneda images.  The
specified natural homotopy therefore lifts through a contractible space to the
required coherent nullhomotopy in $QC_B^\vee$.
\end{proof}

\begin{theorem}[Field-action factorization of the first Noether derivative]
\label{thm:symmetry-noether-field-action-factorization-of-the-first-noether-derivative}
There is a canonical descended pro-coherent morphism
\[
\boxed{
\partial_{\mathrm{fld}}\mathbf n_{\mathcal G,L}^{1,q}:
T^{pc}_{\overline A/Q}[1]
\longrightarrow
U\mathbb O_{\mathcal G,L}^{1,q}[n]
}
\]
and a specified coherent factorization
\[
\boxed{
d_0\mathbf n_{\mathcal G,L}^{1,q}
\simeq
\partial_{\mathrm{fld}}\mathbf n_{\mathcal G,L}^{1,q}
\circ
\rho_{\mathcal G}^{\pi,E_\infty}.
}
\]
The construction is intrinsic to the differentiated formalized Noether section and
does not produce a finite-order differential operator on an ordinary parameter
bundle.
\end{theorem}

\begin{proof}
Work first on one coherent affine-\'etale chart and put
\[
f:=
d_0\mathbf n_{\mathcal G,L}^{1,q}\,i_F:
UF_{\mathcal G}^{\pi,E_\infty}
\longrightarrow
U\mathbb O_{\mathcal G,L}^{1,q}[n].
\]
The zero-action-cell lemma supplies the prescribed nullhomotopy of $fu$ on every
negative square-zero generator and, after stabilization, on its finite stable
idempotent closure.  The preceding generator-detection lemma therefore gives the
canonical coherent path
\[
\boxed{
d_0\mathbf n_{\mathcal G,L}^{1,q}\,i_F\simeq0.
}
\]

By construction of the homotopy anchor fibre there is a fibre sequence in the
underlying stable pro-coherent category
\[
UF_{\mathcal G}^{\pi,E_\infty}
\xrightarrow{\ i_F\ }
U\mathbb A_{\mathcal G}^{\pi,E_\infty}
\xrightarrow{\ \rho_{\mathcal G}^{\pi,E_\infty}\ }
T^{pc}_{\overline A/Q}[1].
\]
Stability makes this also a cofiber sequence.  Applying
\[
\operatorname{Map}
\left(-,U\mathbb O_{\mathcal G,L}^{1,q}[n]\right)
\]
gives a fibre sequence of mapping spaces in which
\[
\operatorname{Map}
\left(T^{pc}_{\overline A/Q}[1],U\mathbb O_{\mathcal G,L}^{1,q}[n]\right)
\]
is the homotopy fibre over zero of restriction along $i_F$.  The pair consisting
of $d_0\mathbf n_{\mathcal G,L}^{1,q}$ and the displayed nullhomotopy is therefore
precisely a point of that fibre.  It determines, through a contractible space,
\[
\partial_{\mathrm{fld}}\mathbf n_{\mathcal G,L}^{1,q}:
T^{pc}_{\overline A/Q}[1]
\longrightarrow
U\mathbb O_{\mathcal G,L}^{1,q}[n]
\]
and the stated coherent factorization through
$\rho_{\mathcal G}^{\pi,E_\infty}$.

The cellwise nullhomotopies, the stable generator extension, and the anchor
fibre/cofiber sequence are all transported by the coherent spectral-\'etale
transition equivalences.  Hence the chartwise factorizations descend to the global
locally coherent theorem domain.
\end{proof}

\begin{theorem}[Anchor-fibre syzygy map]
\label{thm:symmetry-noether-anchor-fibre-syzygy-map}
The canonical nullhomotopy
\[
d_0\mathbf n_{\mathcal G,L}^{1,q}\,i_F\simeq0
\]
determines a canonical morphism
\[
\boxed{
UF_{\mathcal G}^{\pi,E_\infty}
\longrightarrow
\operatorname{Syz}_{\mathcal G,L}^{1,q}.
}
\]
Under the factorization theorem this is the lift induced by the tautological
nullhomotopy of
$\rho_{\mathcal G}^{\pi,E_\infty}i_F$.
\end{theorem}

\begin{proof}
The first assertion is the universal property of the homotopy fibre defining
$\operatorname{Syz}_{\mathcal G,L}^{1,q}$.  For the second, substitute
\[
d_0\mathbf n_{\mathcal G,L}^{1,q}
\simeq
\partial_{\mathrm{fld}}\mathbf n_{\mathcal G,L}^{1,q}
\rho_{\mathcal G}^{\pi,E_\infty}
\]
and compose the tautological anchor-fibre nullhomotopy with
$\partial_{\mathrm{fld}}\mathbf n_{\mathcal G,L}^{1,q}$.
\end{proof}

\begin{construction}[Actual-isotropy-to-Noether comparison]
\label{constr:symmetry-noether-actual-isotropy-to-noether-comparison}
Compose the independently constructed actual-isotropy comparison
\[
i_{\mathcal G}^{\pi,E_\infty}
\xrightarrow{\gamma_{\mathcal G}^{\pi,E_\infty}}
F_{\mathcal G}^{\pi,E_\infty}
\]
with the preceding syzygy map:
\[
\boxed{
Ui_{\mathcal G}^{\pi,E_\infty}
\longrightarrow
\operatorname{Syz}_{\mathcal G,L}^{1,q}.
}
\]
This is the actual-isotropy contribution to the first differentiated Noether
syzygy.
\end{construction}

\begin{proposition}[Isotropy-exact Noether syzygy comparison]
\label{prop:symmetry-noether-isotropy-exact-noether-syzygy-comparison}
After pullback to $\operatorname{IsoEq}_{\mathcal G}$, the zero-action syzygy map
\[
UF_{\mathcal G}^{\pi,E_\infty}
\longrightarrow
\operatorname{Syz}_{\mathcal G,L}^{1,q}
\]
and the actual-isotropy syzygy map
\[
Ui_{\mathcal G}^{\pi,E_\infty}
\longrightarrow
\operatorname{Syz}_{\mathcal G,L}^{1,q}
\]
agree under the universal inverse of
$U\gamma_{\mathcal G}^{\pi,E_\infty}$.  Away from this realization their failure
to be identified remains measured by the constructed discrepancy
$C_{\mathcal G}^{\mathrm{iso}}$.
\end{proposition}

\begin{proof}
The actual-isotropy syzygy map was defined as the composite of
$U\gamma_{\mathcal G}^{\pi,E_\infty}$ with the zero-action syzygy map.  Pull back
the universal coherent inverse supplied by the walking-equivalence point of
$\operatorname{IsoEq}_{\mathcal G}$ and cancel
that equivalence.  The final statement is the definition of the discrepancy
cofiber.
\end{proof}

\section{The differentiated Brave-New Second Noether theorem}
\label{sec:symmetry-noether-the-differentiated-brave-new-second-noether-theorem}

The native symmetry-coherence obstruction tower records extension of a Lagrangian
through the raw cosimplicial symmetry totalization.  A different theorem appears after
Artinian formalization and differentiation.  The parameter of that theorem must
itself be a formal-moduli object.  Accordingly the complete formalized symmetry is
first restricted to Artinian points of a formal leaf and then reflected by the
relative unstable Schlessinger localization constructed above.

Fix a coherent affine-\'etale formal leaf $F$ and abbreviate
\[
\widehat{\operatorname{Sym}}^{q}
:=
\widehat{\operatorname{Sym}}_{\mathcal G,L}^{\mathrm{Sch},q}.
\]
Write
\[
\mathfrak s_{\mathcal G,L}^{q}
:=
\operatorname{Diff}^{\pi,E_\infty}_{B/R}
\left(
\widehat{\operatorname{Sym}}^{q}
\right)
\]
and let
\[
\vartheta_{\mathfrak s}:
\mathfrak s_{\mathcal G,L}^{q}
\longrightarrow
\mathfrak a_{\mathcal G}
\]
be the morphism obtained by differentiating
$\widehat{\operatorname{Sym}}^{q}\to F$.

\begin{construction}[Base-changed retractive Noether algebroid]
\label{constr:symmetry-noether-base-changed-retractive-noether-algebroid}
Pull the relative infinite-loop object back to the reflected symmetry parameter:
\[
Y_{\mathfrak s}
:=
Y_{\mathcal G,L}^{1,q}
\times_F
\widehat{\operatorname{Sym}}^{q}.
\]
The affine formal-moduli/partition-Lie equivalence preserves finite limits.
Therefore
\[
\operatorname{Diff}^{\pi,E_\infty}_{B/R}(Y_{\mathfrak s})
\simeq
\mathfrak b_{\mathcal G,L}^{1,q}
\times_{\mathfrak a_{\mathcal G}}
\mathfrak s_{\mathcal G,L}^{q}.
\]
Denote this pullback algebroid by
\[
\mathfrak b_{\mathfrak s}.
\]
The pulled-back zero and Noether sections differentiate to splittings
\[
z_{\mathfrak s},\ s_{\mathfrak n,\mathfrak s}:
\mathfrak s_{\mathcal G,L}^{q}
\rightrightarrows
\mathfrak b_{\mathfrak s}.
\]
\end{construction}

\begin{theorem}[Differentiated Brave-New Second Noether theorem]
\label{thm:symmetry-noether-differentiated-brave-new-second-noether-theorem}
Over the reflected complete-symmetry formal-moduli parameter
$\widehat{\operatorname{Sym}}_{\mathcal G,L}^{\mathrm{Sch},q}$ there is a
canonical homotopy of spectral partition-Lie algebroid splittings
\[
\boxed{
s_{\mathfrak n,\mathfrak s}
\simeq
z_{\mathfrak s}.
}
\]
Consequently:
\begin{enumerate}
\item the pullback of the first differentiated Noether morphism is canonically
      nullhomotopic,
      \[
      \boxed{
      d_0\mathbf n_{\mathcal G,L}^{1,q}
      \circ
      U\vartheta_{\mathfrak s}
      \simeq
      0
      }
      \]
      as a morphism
      $U\mathfrak s_{\mathcal G,L}^{q}
      \to U\mathbb O_{\mathcal G,L}^{1,q}[n]$;
\item the nullhomotopy determines a canonical morphism
      \[
      U\mathfrak s_{\mathcal G,L}^{q}
      \longrightarrow
      \operatorname{Syz}_{\mathcal G,L}^{1,q}
      \]
      over $U\mathfrak a_{\mathcal G}$;
\item for every finite nonempty $I$ and every labelled tangent tuple in
      $\mathfrak s_{\mathcal G,L}^{q}$, the exact labelled Taylor morphism obtained
      after transport along
      $\vartheta_{\mathfrak s}:\mathfrak s_{\mathcal G,L}^{q}\to\mathfrak a_{\mathcal G}$
      is canonically nullhomotopic;
\item the complete morphism corolla--latching hierarchy of the pulled-back
      Noether splitting is transported coherently to the zero-splitting hierarchy.
\end{enumerate}
\end{theorem}

\begin{proof}
The preceding section constructs, in the formal-moduli slice over $F$, a specified
homotopy
\[
\widehat{\mathbf n}_{\mathcal G,L}^{1,q}\simeq z
\]
after pullback to
$\widehat{\operatorname{Sym}}_{\mathcal G,L}^{\mathrm{Sch},q}$.
The spectral formal-moduli/partition-Lie equivalence is an equivalence of
$\infty$-categories and therefore preserves finite limits, maps, paths, and all
higher simplices.  Differentiating the pullback of
$Y_{\mathcal G,L}^{1,q}\to F$ gives
$\mathfrak b_{\mathfrak s}\to\mathfrak s$, and differentiating the displayed
section homotopy gives
$s_{\mathfrak n,\mathfrak s}\simeq z_{\mathfrak s}$.

Taking the vertical stable fibre of the two splittings gives the first
nullhomotopy.  The universal property of
$\operatorname{Syz}_{\mathcal G,L}^{1,q}$ then gives the displayed syzygy map.
For a labelled tangent tuple in $\mathfrak s_{\mathcal G,L}^{q}$, evaluate the
homotopy of actual partition-Lie algebroid splittings on the corresponding exact
labelled source corolla.  Functoriality of the multi-input operation system carries
the homotopy to a homotopy between the two corolla evaluations.  Their vertical
difference is, by the Taylor construction, the pullback of
$\operatorname{Tay}_{\mathcal G,L}^{I,q}$; hence that Taylor morphism is
canonically nullhomotopic.  Finally the morphism corolla--latching construction is
functorial in an actual morphism of partition-Lie algebroids and in homotopies
between such morphisms, so the entire pulled-back filler system is transported to
the zero system.
\end{proof}

\begin{construction}[Complete-symmetry section realization and universal evaluation]
\label{constr:symmetry-noether-complete-symmetry-section-realization-and-universal-evaluation}
Let
\[
\mathcal F_F
:=
\bigl(\operatorname{FMP}^{\mathrm{sp}}_{B/R}\bigr)/F
\]
and define the section-realization space
\[
\boxed{
K_{\mathcal G,L}^{q}
:=
\operatorname{SymSect}_{\mathcal G,L}^{\mathrm{Sch},q}
:=
\operatorname{Map}_{\mathcal F_F}
\left(
F,
\widehat{\operatorname{Sym}}_{\mathcal G,L}^{\mathrm{Sch},q}
\right).
}
\]
Its possible emptiness is the obstruction to a strong symmetry realization on the
entire formal leaf.

The Schlessinger-local formal-moduli category is presentable, hence its slice
$\mathcal F_F$ is canonically tensored over spaces.  Regard
$F\xrightarrow{\mathrm{id}}F$ as the terminal object of the slice and write
$K_{\mathcal G,L}^{q}\odot F$ for its tensor by the section space.  The tensor--mapping
adjunction sends the identity of $K_{\mathcal G,L}^{q}$ to the canonical evaluation
morphism
\[
\boxed{
\operatorname{ev}_{\mathrm{Sym}}:
K_{\mathcal G,L}^{q}\odot F
\longrightarrow
\widehat{\operatorname{Sym}}_{\mathcal G,L}^{\mathrm{Sch},q}
}
\]
in $\mathcal F_F$.  This is the universal family of complete reflected symmetries;
its restriction to a point
$\sigma\in K_{\mathcal G,L}^{q}$ is the corresponding section
\[
\sigma:F\longrightarrow
\widehat{\operatorname{Sym}}_{\mathcal G,L}^{\mathrm{Sch},q}.
\]
\end{construction}

\begin{theorem}[Universal strong differentiated Second Noether form]
\label{thm:symmetry-noether-universal-strong-differentiated-second-noether-form}
Let
\[
K:=K_{\mathcal G,L}^{q}.
\]
Differentiation of the universal evaluation family gives a morphism in the slice
over $\mathfrak a_{\mathcal G}$
\[
\boxed{
\operatorname{ev}_{\mathrm{Diff}}:
\mathfrak a_K
\longrightarrow
\mathfrak s_{\mathcal G,L}^{q},
}
\]
where
\[
\boxed{
\mathfrak a_K
:=
\operatorname{Diff}^{\pi,E_\infty}_{B/R}(K\odot F)
\simeq
K\odot\mathfrak a_{\mathcal G}
}
\]
and
\[
\operatorname{pr}_K:
\mathfrak a_K
\longrightarrow
\mathfrak a_{\mathcal G}
\]
is its structural morphism in the slice.  The evaluation satisfies
\[
\vartheta_{\mathfrak s}\operatorname{ev}_{\mathrm{Diff}}
\simeq
\operatorname{pr}_K.
\]

Define the explicitly pulled-back retractive Noether algebroid
\[
\boxed{
\mathfrak b_K
:=
\mathfrak a_K
\times_{\mathfrak a_{\mathcal G}}
\mathfrak b_{\mathcal G,L}^{1,q}
}
\]
and let
\[
s_{\mathfrak n,K},\ z_K:
\mathfrak a_K
\rightrightarrows
\mathfrak b_K
\]
be the base changes of the Noether and zero splittings.  In the tangent
representation category put
\[
\boxed{
\mathbb O_K
:=
\operatorname{pr}_K^*\mathbb O_{\mathcal G,L}^{1,q},
}
\]
where $\operatorname{pr}_K^*$ is the functorial pullback in the partition-Lie
tangent-representation slice.  The vertical stable fibre of
$\mathfrak b_K\to\mathfrak a_K$ is canonically
$U\mathbb O_K[n]$.

Pullback of the relative differentiated Second Noether homotopy along
$\operatorname{ev}_{\mathrm{Diff}}$ gives the correctly typed homotopy of actual
spectral partition-Lie algebroid splittings
\[
\boxed{
s_{\mathfrak n,K}\simeq z_K.
}
\]
Taking their vertical difference gives a canonical morphism
\[
(d_0\mathbf n_{\mathcal G,L}^{1,q})_K:
U\mathfrak a_K
\longrightarrow
U\mathbb O_K[n]
\]
and a coherent nullhomotopy
\[
\boxed{
(d_0\mathbf n_{\mathcal G,L}^{1,q})_K\simeq0.
}
\]
These homotopies carry the parameterized global syzygy lift, the coherent
nullification of every exact labelled Taylor morphism after restriction to
$\mathfrak a_K$, and the compatible identification of the complete Noether
morphism-latching hierarchy with the zero hierarchy.

For every point
\[
\sigma\in
\operatorname{SymSect}_{\mathcal G,L}^{\mathrm{Sch},q}=K,
\]
tensoring $\sigma:*\to K$ with the terminal object of the slice determines a
section
\[
i_\sigma:
\mathfrak a_{\mathcal G}
\longrightarrow
\mathfrak a_K,
\qquad
\operatorname{pr}_K i_\sigma\simeq\operatorname{id}_{\mathfrak a_{\mathcal G}}.
\]
Pullback along $i_\sigma$ recovers the strong form on the whole controller; in
particular it gives
\[
\boxed{
d_0\mathbf n_{\mathcal G,L}^{1,q}\simeq0
}
\]
for that selected complete reflected symmetry section, together with the
corresponding syzygy, Taylor, and complete morphism-latching nullifications.
\end{theorem}

\begin{proof}
The category $\mathcal F_F$ is a slice of an accessible localization of a
presentable functor category and is therefore presentable and tensored over
spaces.  The evaluation map is the adjoint of
$\operatorname{id}_{K}$ under
\[
\operatorname{Map}_{\mathcal F_F}
(K\odot F,\widehat{\operatorname{Sym}})
\simeq
\operatorname{Map}_{\mathcal S}
\left(
K,
\operatorname{Map}_{\mathcal F_F}(F,\widehat{\operatorname{Sym}})
\right).
\]
The affine formal-moduli/partition-Lie differentiation functor is an equivalence of
presentable $\infty$-categories.  It therefore preserves the tensor by $K$ and
carries the universal evaluation to
$\operatorname{ev}_{\mathrm{Diff}}:\mathfrak a_K\to\mathfrak s$.  Because the
original evaluation is a morphism over $F$, its differentiated composite with
$\vartheta_{\mathfrak s}$ is the structural map
$\operatorname{pr}_K:\mathfrak a_K\to\mathfrak a_{\mathcal G}$.

The relative Second Noether theorem gives
\[
s_{\mathfrak n,\mathfrak s}\simeq z_{\mathfrak s}
\]
in the pullback
\[
\mathfrak b_{\mathfrak s}
=
\mathfrak s\times_{\mathfrak a_{\mathcal G}}
\mathfrak b_{\mathcal G,L}^{1,q}.
\]
Pulling this homotopy back along
$\operatorname{ev}_{\mathrm{Diff}}$ gives a homotopy in
\[
\mathfrak a_K
\times_{\mathfrak s}
\mathfrak b_{\mathfrak s}
\simeq
\mathfrak a_K
\times_{\mathfrak a_{\mathcal G}}
\mathfrak b_{\mathcal G,L}^{1,q}
=
\mathfrak b_K.
\]
Under this canonical identification the two pulled-back splittings are exactly
$s_{\mathfrak n,K}$ and $z_K$.  This proves
$s_{\mathfrak n,K}\simeq z_K$.

Tangent-representation functoriality identifies the vertical stable fibre of
$\mathfrak b_K\to\mathfrak a_K$ with
$U\operatorname{pr}_K^*\mathbb O_{\mathcal G,L}^{1,q}[n]=U\mathbb O_K[n]$.
Taking the vertical difference of the two homotopic splittings therefore gives the
map $(d_0\mathbf n)_K$ together with its canonical nullhomotopy.  The syzygy,
exact labelled Taylor, and morphism corolla--latching statements follow by the
same functoriality used in the relative Second Noether theorem, now applied to the
actual map $\mathfrak a_K\to\mathfrak a_{\mathcal G}$ and the actual retractive
algebroid $\mathfrak b_K$.

Finally a point $\sigma:*\to K$ induces the displayed section
$i_\sigma:\mathfrak a_{\mathcal G}\to\mathfrak a_K$ of $\operatorname{pr}_K$.
Pulling the homotopy of splittings and its vertical difference along $i_\sigma$
identifies $\mathbb O_K$ with the original representation and recovers
$s_{\mathfrak n}\simeq z_{\mathfrak a}$, hence
$d_0\mathbf n_{\mathcal G,L}^{1,q}\simeq0$, on the whole controller.  All higher
coherences are preserved by the same pullback.
\end{proof}

\begin{remark}[Why this theorem carries the Second Noether name]
\label{rem:symmetry-noether-why-this-theorem-carries-the-second-noether-name}
The theorem concerns the differentiated dependence of the Euler/Noether
obstruction on symmetry parameters and their higher coherent relations.  It is
therefore separated from the earlier symmetry-totalization lifting theorem, whose
index records cosimplicial coherence rather than differentiated symmetry dependence.
The universal theorem is naturally relative to the differentiated symmetry
parameter.  The stronger nullhomotopy on the whole controller is itself universal
over the constructed section-realization space
$\operatorname{SymSect}_{\mathcal G,L}^{\mathrm{Sch},q}$ rather than being stated
under a section-existence assumption.
\end{remark}

\section{Integral Noether--Spencer cocycles and the tangent representation bar}
\label{sec:symmetry-noether-integral-noether-spencer-cocycles-and-the-tangent-representation-bar}

The integral tangent representation theory of Chapter 15 has two constructions
which must remain distinct.  A representation has a simplicial one-marked monadic
bar whose stable normalization is a \emph{chain} object.  A section of an
infinite-loop representation, after differentiation, gives instead a splitting of
a retractive partition-Lie algebroid.  The Noether cocycle is the latter splitting;
the former supplies its intrinsic coefficient-action resolution.

Put
\[
\mathfrak a:=\mathfrak a_{\mathcal G},
\qquad
M_{\mathcal G,L}^{q}:=
U\mathbb O_{\mathcal G,L}^{1,q}[n].
\]

\begin{definition}[Integral Noether--Spencer cocycle space]
\label{def:symmetry-noether-integral-noether-spencer-cocycle-space}
Let
\[
\mathfrak b_{\mathcal G,L}^{1,q}\longrightarrow\mathfrak a
\]
be the differentiated retractive algebroid on a coherent affine-\'etale chart.
Define
\[
\boxed{
\operatorname{SpCocyc}_{\mathcal G,L}^{q}
:=
\operatorname{Map}_{\operatorname{LieAlgd}^{\pi,E_\infty}/\mathfrak a}
\left(
\mathfrak a,
\mathfrak b_{\mathcal G,L}^{1,q}
\right).
}
\]
It has the distinguished zero point $z_{\mathfrak a}$ and distinguished Noether
point $s_{\mathfrak n}$.  Under coherent spectral-\'etale transport these mapping
spaces form the \'etale mapping sheaf constructed above; the same notation denotes
its descended global object of choices.
\end{definition}

\begin{construction}[Universal finite-weight splitting tower in the slice]
\label{constr:symmetry-noether-universal-finite-weight-splitting-tower}
Write
\[
p:
\mathfrak b_{\mathcal G,L}^{1,q}
\longrightarrow
\mathfrak a
\]
for the differentiated retractive algebroid projection.  Define the weight-zero
morphism-action space for every underlying anchored map $h$ by
\[
\boxed{
\operatorname{MorAct}_{\le0}(h):=*.
}
\]
Accordingly set
\[
\boxed{
\operatorname{Mor}_{\le0}(\mathfrak a,\mathfrak b_{\mathcal G,L}^{1,q})
:=
\operatorname{Map}_{\operatorname{Anch}^{pc}_{B/R}}
(U\mathfrak a,U\mathfrak b_{\mathcal G,L}^{1,q}),
}
\]
and define $\operatorname{Mor}_{\le0}(\mathfrak a,\mathfrak a)$ analogously.  Thus
weight zero remembers exactly the underlying anchored morphism and introduces no
additional relation.

For $r\ge1$, let
\[
\operatorname{Mor}_{\le r}(\mathfrak a,\mathfrak b_{\mathcal G,L}^{1,q})
\]
be the total space, over the space of underlying anchored maps
\[
h:
U\mathfrak a
\longrightarrow
U\mathfrak b_{\mathcal G,L}^{1,q},
\]
of the Chapter 15 spaces $\operatorname{MorAct}_{\le r}(h)$ of coherent
algebroid-morphism relations through unmarked weight $r$.  Equivalently,
\[
\operatorname{Mor}_{\le r}(\mathfrak a,\mathfrak b_{\mathcal G,L}^{1,q})
:=
\int_{
h\in
\operatorname{Map}_{\operatorname{Anch}^{pc}_{B/R}}
(U\mathfrak a,U\mathfrak b_{\mathcal G,L}^{1,q})
}
\operatorname{MorAct}_{\le r}(h),
\qquad r\ge1.
\]
Define $\operatorname{Mor}_{\le r}(\mathfrak a,\mathfrak a)$ in the same way.

Because $p$ is an actual spectral partition-Lie algebroid morphism,
postcomposition induces
\[
p_{\le r,*}:
\operatorname{Mor}_{\le r}(\mathfrak a,\mathfrak b_{\mathcal G,L}^{1,q})
\longrightarrow
\operatorname{Mor}_{\le r}(\mathfrak a,\mathfrak a).
\]
The identity algebroid morphism determines compatible points
\[
\operatorname{id}_{\mathfrak a}^{\le r}
\in
\operatorname{Mor}_{\le r}(\mathfrak a,\mathfrak a).
\]
Define
\[
\boxed{
\operatorname{SpCocyc}_{\mathcal G,L}^{q,\le r}
:=
\operatorname{hofib}_{\operatorname{id}_{\mathfrak a}^{\le r}}
\left(
p_{\le r,*}
\right).
}
\]
Thus the finite stage is formed in the slice over $\mathfrak a$.  In particular,
a point of $\operatorname{SpCocyc}_{\mathcal G,L}^{q,\le0}$ is an underlying
anchored map
\[
h:
U\mathfrak a
\longrightarrow
U\mathfrak b_{\mathcal G,L}^{1,q}
\]
together with the prescribed coherent identity
\[
Up\circ h\simeq\operatorname{id}_{U\mathfrak a}.
\]

More explicitly, suppose
\[
x_{r-1}\in
\operatorname{SpCocyc}_{\mathcal G,L}^{q,\le r-1}
\]
is a partial splitting with underlying anchored map $h_{r-1}$.  The morphism
corolla--latching theorem gives the weight-$r$ extension space
\[
\operatorname{Ext}^{\mathfrak b}_r(h_{r-1})
:=
\operatorname{hofib}_{g_r^{\mathfrak b}}
\left(
\underline{\operatorname{Map}}
(C_r^{\mathrm{mor}}(h_{r-1}),U\mathfrak b_{\mathcal G,L}^{1,q})
\longrightarrow
\underline{\operatorname{Map}}
(L_r^{\mathrm{mor}}(h_{r-1}),U\mathfrak b_{\mathcal G,L}^{1,q})
\right),
\]
equivalently the space of specified nullhomotopies of the corresponding
connecting obstruction together with the corolla filler.  Postcomposition with
$p$ gives
\[
\operatorname{Ext}^{\mathfrak b}_r(h_{r-1})
\longrightarrow
\operatorname{Ext}^{\mathfrak a}_r(ph_{r-1}).
\]
The slice datum in $x_{r-1}$ identifies $ph_{r-1}$ with the identity through
weight $r-1$, while the actual identity morphism supplies its prescribed
weight-$r$ filler
\[
e_r^{\operatorname{id}}
\in
\operatorname{Ext}^{\mathfrak a}_r(\operatorname{id}_{\mathfrak a}^{\le r-1}).
\]
Under the induced identification of target extension spaces, the fibre of
\[
\operatorname{SpCocyc}_{\mathcal G,L}^{q,\le r}
\longrightarrow
\operatorname{SpCocyc}_{\mathcal G,L}^{q,\le r-1}
\]
over $x_{r-1}$ is canonically
\[
\boxed{
\operatorname{hofib}_{e_r^{\operatorname{id}}}
\left(
\operatorname{Ext}^{\mathfrak b}_r(h_{r-1})
\longrightarrow
\operatorname{Ext}^{\mathfrak a}_r(ph_{r-1})
\right).
}
\]
Hence a weight-$r$ filler is retained only when its image under $p$ is the
already prescribed weight-$r$ identity filler, including the required higher
coherence.
\end{construction}

\begin{theorem}[Finite-weight reconstruction of integral Noether--Spencer cocycles]
\label{thm:symmetry-noether-finite-weight-reconstruction-of-integral-noether-spencer-cocycles}
There is a canonical equivalence
\[
\boxed{
\operatorname{SpCocyc}_{\mathcal G,L}^{q}
\simeq
\lim_r
\operatorname{SpCocyc}_{\mathcal G,L}^{q,\le r}.
}
\]
Pullback of the universal tower along the Noether point $s_{\mathfrak n}$ recovers
\[
L_r^{\mathrm{Noeth}},\qquad
C_r^{\mathrm{Noeth}},\qquad
R_r^{\mathrm{Noeth}},\qquad
\operatorname{ob}^{\mathrm{Noeth}}_r,
\]
while pullback along the zero point gives the zero-splitting hierarchy.
\end{theorem}

\begin{proof}
Let
\[
T_{\pi}:=\operatorname{Lie}^{\pi,E_\infty}_{B/R}
\]
be the full spectral partition-Lie monad.  Chapter 15 equips its underlying
operation endofunctor with the exhaustive anchor-weight filtration
\[
W_{\le0}T_{\pi}\longrightarrow
W_{\le1}T_{\pi}\longrightarrow\cdots,
\qquad
\operatorname*{colim}_{r}W_{\le r}T_{\pi}\simeq T_{\pi}.
\]
The finite pieces are pieces of one filtered monad; they are not used here as
independent monads.

Fix an underlying anchored map
\[
h:U\mathfrak a\longrightarrow U\mathfrak b_{\mathcal G,L}^{1,q}.
\]
Every unit, substitution, and algebra-morphism relation is represented by a
finite operation tree and therefore has finite anchor weight.  The parallel
morphism corolla--latching theorem consequently identifies a full coherent
algebroid-morphism structure on $h$ with a compatible point of all of the
finite-weight spaces $\operatorname{MorAct}_{\le r}(h)$.  Hence
\[
\operatorname{Map}_{\operatorname{LieAlgd}^{\pi,E_\infty}}
\left(
\mathfrak a,\mathfrak b_{\mathcal G,L}^{1,q}
\right)
\simeq
\lim_r
\operatorname{Mor}_{\le r}
\left(
\mathfrak a,\mathfrak b_{\mathcal G,L}^{1,q}
\right).
\]
The same argument gives
\[
\operatorname{Map}_{\operatorname{LieAlgd}^{\pi,E_\infty}}
(\mathfrak a,\mathfrak a)
\simeq
\lim_r
\operatorname{Mor}_{\le r}(\mathfrak a,\mathfrak a),
\]
and these equivalences are compatible with postcomposition by
$p:\mathfrak b_{\mathcal G,L}^{1,q}\to\mathfrak a$ and with the identity
points.

By definition of a mapping space in the slice,
\[
\operatorname{SpCocyc}_{\mathcal G,L}^{q}
\simeq
\operatorname{hofib}_{\operatorname{id}_{\mathfrak a}}
\left(
\operatorname{Map}_{\operatorname{LieAlgd}^{\pi,E_\infty}}
(\mathfrak a,\mathfrak b_{\mathcal G,L}^{1,q})
\longrightarrow
\operatorname{Map}_{\operatorname{LieAlgd}^{\pi,E_\infty}}
(\mathfrak a,\mathfrak a)
\right).
\]
Homotopy fibres commute with inverse limits in spaces.  Substituting the two
finite-weight reconstruction equivalences therefore gives
\[
\operatorname{SpCocyc}_{\mathcal G,L}^{q}
\simeq
\lim_r
\operatorname{hofib}_{\operatorname{id}_{\mathfrak a}^{\le r}}
\left(
\operatorname{Mor}_{\le r}
(\mathfrak a,\mathfrak b_{\mathcal G,L}^{1,q})
\longrightarrow
\operatorname{Mor}_{\le r}(\mathfrak a,\mathfrak a)
\right),
\]
which is exactly
\[
\operatorname{SpCocyc}_{\mathcal G,L}^{q}
\simeq
\lim_r\operatorname{SpCocyc}_{\mathcal G,L}^{q,\le r}.
\]
The explicit relative corolla-fibre description in the preceding construction
shows stagewise that the projection to $\mathfrak a$ is the prescribed identity
coherence.  Naturality of the morphism corolla--latching construction therefore
gives the pullback statements at the Noether and zero points.
\end{proof}

\begin{construction}[Tangent representation Spencer chain object]
\label{constr:symmetry-noether-tangent-representation-spencer-chain-object}
The differentiated coefficient
$\mathbb O_{\mathcal G,L}^{1,q}$ is a representation of $\mathfrak a$ and hence
an algebra over the tangent enveloping monad
\[
T:=U^\pi_{\mathfrak a}.
\]
Form the simplicial monadic bar
\[
B^{\mathrm{tan}}_m
\left(
\mathfrak a;\mathbb O_{\mathcal G,L}^{1,q}
\right)
:=
T^{m+1}
\left(
U\mathbb O_{\mathcal G,L}^{1,q}
\right),
\qquad m\ge0.
\]
The nonterminal faces use monad multiplication, the terminal face uses the
representation action, and the degeneracies use the monad unit.  Stable
simplicial normalization gives the chain object
\[
\boxed{
B^{\mathrm{Sp}}_{\mathfrak a}
\left(
U\mathbb O_{\mathcal G,L}^{1,q}
\right)
\in
\operatorname{Ch}_{\ge0}(QC_B^\vee).
}
\]
Its canonical normalized bar augmentation is
\[
\tau^{\mathrm{Sp}}_{\mathfrak a,\mathbb O}:
B^{\mathrm{Sp}}_{\mathfrak a}
\left(
U\mathbb O_{\mathcal G,L}^{1,q}
\right)
\longrightarrow
U\mathbb O_{\mathcal G,L}^{1,q}[0].
\]
\end{construction}

\begin{proposition}[Weight filtration of the representation bar]
\label{prop:symmetry-noether-weight-filtration-of-the-representation-bar}
The additive one-marked weight filtration of the tangent enveloping monad induces
an exhaustive increasing filtration
\[
W_{\le0}B^{\mathrm{Sp}}_{\mathfrak a}
\longrightarrow
W_{\le1}B^{\mathrm{Sp}}_{\mathfrak a}
\longrightarrow\cdots
\longrightarrow
B^{\mathrm{Sp}}_{\mathfrak a}.
\]
The weight-zero piece is the unary identity coefficient, the first positive
graded piece is the universal anchored first-order action on
$U\mathbb O_{\mathcal G,L}^{1,q}$, and the normalized twisting identity contains
the curvature, Bianchi, and higher representation relations of Chapter 15.
\end{proposition}


\subsection{Infinitesimal connection and curvature realizations}
\label{subsec:symmetry-noether-infinitesimal-connection-and-curvature-realizations}

\begin{construction}[Canonical Noether coefficient connection]
\label{constr:symmetry-noether-canonical-noether-coefficient-connection}
Let $T=U^\pi_{\mathfrak a}$ be the tangent enveloping monad and retain its constructed
one-marked weight filtration.  The full representation action
\[
T\left(U\mathbb O_{\mathcal G,L}^{1,q}\right)
\longrightarrow
U\mathbb O_{\mathcal G,L}^{1,q}
\]
restricts to the first-weight stage and gives
\[
\boxed{
\nabla_{\mathcal G,L}^{\mathrm{Noeth}}:
F_{\le1}T\left(U\mathbb O_{\mathcal G,L}^{1,q}\right)
\longrightarrow
U\mathbb O_{\mathcal G,L}^{1,q},
}
\]
whose restriction along
$F_{\le0}T\simeq\operatorname{id}$ is the identity.  Thus
$\nabla_{\mathcal G,L}^{\mathrm{Noeth}}$ is a canonically constructed coefficient
connection; it is not an input datum.
\end{construction}

\begin{theorem}[Flatness and Bianchi hierarchy of the Noether coefficient connection]
\label{thm:symmetry-noether-flatness-and-bianchi-hierarchy-of-noether-coefficient-connection}
Every corolla-latching obstruction to extending
$\nabla_{\mathcal G,L}^{\mathrm{Noeth}}$ through the next marked weight carries the
canonical filler obtained by restricting the already existing full representation
action.  In particular its weight-two curvature has a canonical nullhomotopy
\[
\boxed{F_{\nabla^{\mathrm{Noeth}}}\simeq0,}
\]
the next monad-defect syzygy supplies the canonical Bianchi coherence, and the same
construction produces the complete higher flatness hierarchy.
\end{theorem}

\begin{proof}
The coefficient connection is the weight-one truncation of an actual algebra over
the complete tangent monad.  The Chapter 15 corolla-latching theorem identifies
extension from weight $r-1$ to weight $r$ with a filler of the corresponding
corolla--latching map, equivalently a nullhomotopy of its obstruction.  The
weight-$r$ component of the given full representation action is exactly such a
filler for every $r$.  At weight two this is the curvature nullification, and the
next universal monad-defect syzygy is the Bianchi homotopy.  Iteration gives the
higher hierarchy.
\end{proof}

\begin{construction}[Terminal identity-anchor algebroid and intrinsic pro-coherent anchor splittings]
\label{constr:symmetry-noether-anchor-splitting-connection-realization-and-curvature}
On a coherent affine-\'etale chart
\[
U=\operatorname{Spec}B,\qquad V=\operatorname{Spec}R,
\]
retain the terminal identity-anchor spectral partition-Lie algebroid of Chapter 14,
\[
\boxed{
\mathbb T_{B/R}^{\pi,E_\infty}
:=
\left(
T_{B/R}^{pc}[1]
\xrightarrow{\ \operatorname{id}\ }
T_{B/R}^{pc}[1]
\right).
}
\]
Its algebra structure is the unique one on the terminal anchored object.  Spectral
\'etale cotangent base change identifies these terminal anchored objects on
overlaps, and uniqueness identifies their algebra structures coherently.  Thus they
descend to the global terminal identity-anchor algebroid
\[
\mathbb T_{\overline A/Q}^{\pi,E_\infty}
\]
on the locally coherent differentiated coefficient system.

The differentiated anchor of
$\mathbb A_{\mathcal G}^{\pi,E_\infty}$ is the underlying anchored map of the
unique algebroid morphism
\[
\widetilde\rho_{\mathcal G}^{\pi,E_\infty}:
\mathbb A_{\mathcal G}^{\pi,E_\infty}
\longrightarrow
\mathbb T_{\overline A/Q}^{\pi,E_\infty},
\qquad
U\widetilde\rho_{\mathcal G}^{\pi,E_\infty}
=
\rho_{\mathcal G}^{\pi,E_\infty}.
\]
On each coherent affine-\'etale chart define
\[
\boxed{
\operatorname{Conn}_{\mathcal G}^{\mathrm{anch},pc}(U/V)
:=
\operatorname{Map}_{\operatorname{Anch}^{pc}_{B/R}}
\left(
U\mathbb T_{B/R}^{\pi,E_\infty},
U\mathbb A_{\mathcal G}^{\pi,E_\infty}|_{U/V}
\right).
}
\]
Because the mapping space is formed in the anchored overcategory, a point is
exactly an underlying anchored splitting
\[
h:
U\mathbb T_{B/R}^{\pi,E_\infty}
\longrightarrow
U\mathbb A_{\mathcal G}^{\pi,E_\infty}|_{U/V}
\]
together with the specified identity
\[
\rho_{\mathcal G}^{\pi,E_\infty}h
\simeq
\operatorname{id}_{T_{B/R}^{pc}[1]}.
\]
The coherent spectral-\'etale mapping-sheaf construction makes these chartwise
spaces descend.  Denote the resulting realization object by
\[
\boxed{
\operatorname{Conn}_{\mathcal G}^{\mathrm{anch},pc}.
}
\]
No represented tangent complex is used in this definition.
\end{construction}

\begin{theorem}[Intrinsic pro-coherent anchor curvature and Bianchi hierarchy]
\label{thm:symmetry-noether-intrinsic-procoherent-anchor-curvature-and-bianchi-hierarchy}
Let
\[
h\in
\operatorname{Conn}_{\mathcal G}^{\mathrm{anch},pc}(U/V)
\]
be a pro-coherent anchor splitting on a coherent affine-\'etale chart.  Apply the
Chapter 15 morphism corolla--latching tower to the underlying anchored map
\[
h:
U\mathbb T_{B/R}^{\pi,E_\infty}
\longrightarrow
U\mathbb A_{\mathcal G}^{\pi,E_\infty}.
\]
Its weight-two morphism obstruction
\[
\operatorname{ob}^{\mathrm{mor}}_2(h):
R_2^{\mathrm{mor}}(h)
\longrightarrow
U\mathbb A_{\mathcal G}^{\pi,E_\infty}
\]
has canonically null composite with the differentiated anchor.  Hence it factors
through the homotopy anchor fibre:
\[
\boxed{
F_h^{\mathrm{anch},pc}:
R_2^{\mathrm{mor}}(h)
\longrightarrow
UF_{\mathcal G}^{\pi,E_\infty}.
}
\]
The next weight of the universal morphism-defect syzygy supplies a canonical
Bianchi nullhomotopy for this curvature, and the higher morphism coherences supply
the corresponding higher Bianchi hierarchy.  The construction is natural in $h$,
in fixed-object morphisms of formal symmetry groupoids, and under coherent
spectral-\'etale transport.
\end{theorem}

\begin{proof}
The target
$\mathbb T_{B/R}^{\pi,E_\infty}$ is an actual spectral partition-Lie algebroid, and
the restriction of
\[
\widetilde\rho_{\mathcal G}^{\pi,E_\infty}:
\mathbb A_{\mathcal G}^{\pi,E_\infty}
\longrightarrow
\mathbb T_{B/R}^{\pi,E_\infty}
\]
is an actual algebroid morphism.  Naturality of the morphism corolla--latching
obstruction tower under postcomposition therefore carries
$\operatorname{ob}^{\mathrm{mor}}_2(h)$ to the weight-two obstruction of the
composite
\[
\rho_{\mathcal G}h.
\]
The defining anchored-splitting homotopy identifies this composite with
$\operatorname{id}_{\mathbb T_{B/R}^{\pi,E_\infty}}$.  The identity is already a
full algebroid morphism, so its weight-two obstruction carries the canonical zero
nullhomotopy.  Consequently
\[
\rho_{\mathcal G}^{\pi,E_\infty}
\operatorname{ob}^{\mathrm{mor}}_2(h)
\simeq0.
\]

By definition
\[
UF_{\mathcal G}^{\pi,E_\infty}
\simeq
\operatorname{fib}
\left(
U\mathbb A_{\mathcal G}^{\pi,E_\infty}
\xrightarrow{\rho_{\mathcal G}^{\pi,E_\infty}}
T_{B/R}^{pc}[1]
\right).
\]
The homotopy-fibre universal property therefore gives the displayed factorization
through a contractible space; this factorization is
$F_h^{\mathrm{anch},pc}$.

The universal morphism-defect syzygy of Chapter 15 is functorial under
postcomposition and its next weight gives the Bianchi nullhomotopy for the
weight-two defect.  Applying the same fibre factorization to that coherent
nullhomotopy yields the asserted Bianchi coherence in the anchor fibre.  The
higher statement is the corresponding restriction of the higher morphism
coherences.  Spectral-\'etale naturality follows because the terminal
identity-anchor algebroid, the differentiated anchor, the morphism-latching tower,
and the anchor-fibre construction are all transported by the coherent
spectral-\'etale coefficient system.
\end{proof}

\begin{construction}[Represented anchor-splitting realization]
\label{constr:symmetry-noether-represented-anchor-splitting-realization}
On the represented coherent first-order sector retain the internal-dual anchor
\[
\rho_{\mathcal G}^{\mathrm{rep}}:
a_{\mathcal G}
\longrightarrow
T_{\overline A/Q}^{\mathrm{rep}}.
\]
Define
\[
\boxed{
\operatorname{Conn}^{\mathrm{anch},rep}_{\mathcal G}
:=
\operatorname{hofib}_{\operatorname{id}}
\left(
\underline{\operatorname{Map}}
\left(T_{\overline A/Q}^{\mathrm{rep}},a_{\mathcal G}\right)
\longrightarrow
\underline{\operatorname{Map}}
\left(T_{\overline A/Q}^{\mathrm{rep}},T_{\overline A/Q}^{\mathrm{rep}}\right)
\right),
}
\]
where the displayed map is postcomposition with
$\rho_{\mathcal G}^{\mathrm{rep}}$.  A point is therefore an ordinary represented
anchor splitting
\[
h^{\mathrm{rep}}:
T_{\overline A/Q}^{\mathrm{rep}}
\longrightarrow
a_{\mathcal G}
\]
with its specified homotopy
$\rho_{\mathcal G}^{\mathrm{rep}}h^{\mathrm{rep}}\simeq\operatorname{id}$.
This construction exists on the represented first-order domain independently of
whether it can be compared with the partition-Lie controller.
\end{construction}

\begin{proposition}[Represented finite-locally-free comparison of anchor splittings and curvature]
\label{prop:symmetry-noether-represented-finite-locally-free-comparison-of-anchor-splittings-and-curvature}
Assume in addition that the represented locally coherent effective-quotient
comparison with
$\mathbb A_{\mathcal G}^{\pi,E_\infty}$ is available and that
\[
L_{\overline A/Q},
\qquad
a_{\mathcal G}^{\vee}
\]
are finite locally free on the represented chart.  Write
\[
D(-):=(-)^{\vee,\mathrm{cont}}
\]
for continuous pro-coherent duality.  On the represented locally coherent overlap,
the differentiated first-order comparison gives the explicit shifted
identifications
\[
\boxed{
U\mathbb A_{\mathcal G}^{\pi,E_\infty}
\simeq
D(a_{\mathcal G}^{\vee})[1],
\qquad
U\mathbb T_{\overline A/Q}^{\pi,E_\infty}
=
T_{\overline A/Q}^{pc}[1]
\simeq
D(L_{\overline A/Q})[1].
}
\]
Under these identifications the differentiated anchor is the simultaneous
$[1]$-suspension of the continuous dual of the cotangent transitivity morphism
\[
L_{\overline A/Q}
\longrightarrow
a_{\mathcal G}^{\vee},
\]
while the represented anchor
\[
\rho_{\mathcal G}^{\mathrm{rep}}:
a_{\mathcal G}\longrightarrow T_{\overline A/Q}^{\mathrm{rep}}
\]
is its unshifted internal dual.

Finite-locally-free duality identifies the two splitting mapping problems after
simultaneous suspension/desuspension.  Consequently, after restricting the
represented realization to the common coherent affine-\'etale comparison site,
there is a canonical equivalence
\[
\boxed{
\left.
\operatorname{Conn}^{\mathrm{anch},rep}_{\mathcal G}
\right|_{\operatorname{AffEt}^{\mathrm{coh}}}
\simeq
\operatorname{Conn}^{\mathrm{anch},pc}_{\mathcal G}.
}
\]
On this overlap we abbreviate the displayed restriction by the unqualified
comparison
$\operatorname{Conn}^{\mathrm{anch},rep}_{\mathcal G}
\simeq
\operatorname{Conn}^{\mathrm{anch},pc}_{\mathcal G}$.

Under this equivalence a represented splitting $h^{\mathrm{rep}}$ determines the
intrinsic pro-coherent splitting $h$, and the intrinsic curvature
\[
F_h^{\mathrm{anch},pc}
\]
and its Bianchi hierarchy are exactly the partition-Lie curvature and Bianchi
hierarchy attached to the represented splitting through that comparison.  On this
overlap write
\[
F_{h^{\mathrm{rep}}}^{\mathrm{anch}}
\]
for the same curvature.  Outside the finite-locally-free comparison locus no
identification between
$\operatorname{Conn}^{\mathrm{anch},rep}_{\mathcal G}$ and
$\operatorname{Conn}^{\mathrm{anch},pc}_{\mathcal G}$ is asserted.
\end{proposition}

\begin{proof}
Put
\[
L_0:=L_{\overline A/Q},
\qquad
L_1:=a_{\mathcal G}^{\vee}.
\]
On the represented locally coherent quotient domain the underlying shifted
spectral partition-Lie anchor is
\[
D(L_1)[1]
\longrightarrow
D(L_0)[1],
\]
the continuous dual, followed by the common suspension, of the cotangent
transitivity map $L_0\to L_1$.  The represented tangent anchor is
\[
L_1^\vee
\longrightarrow
L_0^\vee.
\]
Because $L_0$ and $L_1$ are finite locally free, continuous duality agrees with
internal duality on them, biduality holds, and simultaneous suspension is an
autoequivalence.  Hence there are canonical equivalences of mapping spaces
\[
\operatorname{Map}_{QC_B^\vee}
\left(D(L_0)[1],D(L_1)[1]\right)
\simeq
\operatorname{Map}_{QC_B^\vee}
\left(D(L_0),D(L_1)\right)
\simeq
\operatorname{Map}_{QCoh(B)}
\left(L_0^\vee,L_1^\vee\right).
\]
Under these equivalences, postcomposition by the shifted pro-coherent anchor
corresponds to postcomposition by $\rho_{\mathcal G}^{\mathrm{rep}}$, and the
identity of $D(L_0)[1]$ corresponds to the identity of $L_0^\vee$.  Taking the
homotopy fibre over that identity therefore identifies the two splitting spaces.
The pro-coherent side is a coherent spectral-\'etale mapping sheaf, so the literal
global comparison is the displayed equivalence after restriction of the
represented realization to the same coherent affine-\'etale site.

The morphism corolla--latching construction is functorial in the underlying
anchored map.  Under the preceding equivalence the represented splitting and the
intrinsic pro-coherent splitting are the same anchored map after the explicit
common $[1]$ suspension/desuspension.  Their weight-two morphism defect, its
canonical factorization through the homotopy anchor fibre, the next universal
morphism-defect syzygy, and all higher Bianchi coherences therefore identify.  This
proves the curvature and Bianchi comparison.
\end{proof}

\begin{remark}[The Noether section is not ``applied to the bar'']
\label{rem:symmetry-noether-the-noether-section-is-not-applied-to-the-bar}
The chain object $B^{\mathrm{Sp}}_{\mathfrak a}$ is constructed from the
\emph{representation action}.  The Noether section itself is encoded by the
separate splitting $s_{\mathfrak n}$, hence by a point of
$\operatorname{SpCocyc}_{\mathcal G,L}^{q}$.  No cochain object is obtained by
applying the simplicial bar to the section.
\end{remark}

\begin{theorem}[Rational dg cocycle interpretation]
\label{thm:symmetry-noether-rational-dg-cocycle-interpretation}
On the rational dg representation theorem domain, rectification carries
\[
\mathfrak b_{\mathcal G,L}^{1,q}\longrightarrow\mathfrak a
\]
to a square-zero semidirect dg Lie-algebroid extension
\[
\mathfrak a_{\mathrm{dg}}\ltimes M_{\mathrm{dg}}
\longrightarrow
\mathfrak a_{\mathrm{dg}}.
\]
A splitting has the form
\[
x\longmapsto (x,c(x)),
\]
and it is a dg Lie-algebroid morphism precisely when $c$ is the corresponding
Chevalley--Eilenberg/Spencer cocycle.  In particular, on homogeneous local
sections,
\[
c([x,y])
\simeq
\nabla_xc(y)
-
(-1)^{|x||y|}
\nabla_yc(x)
\]
together with the differential compatibility and its higher coherent version.
Thus the integral splitting object rectifies to the ordinary dg
Noether--Spencer cocycle space on this theorem domain.
\end{theorem}

\begin{proof}
Rational rectification identifies the spectral partition-Lie representation
category with the dg Lie-algebroid representation category and carries a
retractive object to the corresponding semidirect square-zero extension.
A section over the base is a dg Lie-algebroid morphism exactly when its vertical
component satisfies the dg derivation/cocycle equation.  Expanding the bracket of
the semidirect product gives the displayed formula.  The representation bar
rectifies independently to the usual enveloping-module bar.
\end{proof}

\section{Derived current-trivialization and Noether transgression objects}
\label{sec:symmetry-noether-derived-current-trivialization-and-noether-transgression-objects}

The one-stage critical locus canonically kills the first Euler source, but it does
not by itself identify a nullhomotopy of the complete coherent Euler representative.
The required promotion is represented rather than assumed.

Retain the common parameter object
\[
\mathcal C_{\mathcal G,L}^{q}
=
\operatorname{Crit}(L)
\times_Q
\mathcal N_{\mathcal G}^{q}(L)
\]
constructed in the one-stage conservation corollary, and pull the universal current
and the complete symmetry Euler representative to this object.

\begin{definition}[Compatible full Euler-nullification object]
\label{def:symmetry-noether-full-euler-nullification-object}
Let
\[
\pi_{\mathrm{crit}}:
\mathcal C_{\mathcal G,L}^{q}\longrightarrow\operatorname{Crit}(L),
\qquad
c_{\mathcal C}:=
j_L\pi_{\mathrm{crit}}:
\mathcal C_{\mathcal G,L}^{q}\longrightarrow\overline A
\]
be the critical projection and its composite with the critical inclusion.  After
pulling every coefficient to $\mathcal C_{\mathcal G,L}^{q}$, write
\[
\sigma_{\mathcal C}
:=
\left.\mathscr E_{\mathcal G,L}^{[q]}\right|_{\mathcal C_{\mathcal G,L}^{q}}
\]
for the complete coherent symmetry Euler representative and
\[
u_{\mathcal C}
:=
c_{\mathcal C}^*
\lambda_{\mathcal G,\mathrm{loc}}^1
\]
for the pulled-back local one-stage restriction.  The one-stage image proposition
gives a canonical equivalence
\[
\beta_{\mathcal C}:
u_{\mathcal C}\sigma_{\mathcal C}
\simeq
c_{\mathcal C}^*E_{\mathcal G,L}^{(1)}.
\]
The critical factor carries the canonical nullhomotopy
\[
h_{\mathrm{crit}}:
c_{\mathcal C}^*E_L\simeq0.
\]
Using
$E_{\mathcal G,L}^{(1)}\simeq\iota_{\mathcal G}^{(1)}E_L$, compose
$\beta_{\mathcal C}$ with the image of $h_{\mathrm{crit}}$ to obtain the specified
one-stage image nullhomotopy
\[
\eta_{\mathcal C}:
u_{\mathcal C}\sigma_{\mathcal C}\simeq0.
\]

Define
\[
\boxed{
\operatorname{EulNull}_{\mathcal G}^{q}(L)
:=
Z_{\mathcal C_{\mathcal G,L}^{q}}
\left(
\sigma_{\mathcal C};
u_{\mathcal C},
\eta_{\mathcal C}
\right).
}
\]
Equivalently,
\[
\operatorname{EulNull}_{\mathcal G}^{q}(L)
\simeq
\mathcal C_{\mathcal G,L}^{q}
\times_{Z_{\mathcal C_{\mathcal G,L}^{q}}
(u_{\mathcal C}\sigma_{\mathcal C})}
Z_{\mathcal C_{\mathcal G,L}^{q}}(\sigma_{\mathcal C}),
\]
where the first map is the section classified by $\eta_{\mathcal C}$.  Thus a
point is a specified coherent nullhomotopy
\[
H:\sigma_{\mathcal C}\simeq0
\]
whose image under $u_{\mathcal C}$ is the already selected nullhomotopy
$\eta_{\mathcal C}$, including the required higher compatibility.  The
construction is a geometric object over $\mathcal C_{\mathcal G,L}^{q}$ by the
parameterized compatible-nullification construction; its inhabitation is a
property rather than an admissibility condition.
\end{definition}

For the symmetry-horizontal calculus define the complete horizontal transgression
coefficient
\[
\operatorname{Trans}_{\mathcal G,H}^{1,q}(A;M)
:=
\operatorname{fib}
\left(
F_H^q\operatorname{Sec}_{\mathcal G}^{1}(A;M)
\longrightarrow
\operatorname{Sec}_{\mathcal G}^{1}(A;M)
\right).
\]

\begin{construction}[Native on-shell Noether transgression]
\label{constr:symmetry-noether-native-on-shell-noether-transgression}
Over $\operatorname{EulNull}_{\mathcal G}^{q}(L)$, compose the First Noether
homotopy
\[
\iota_H^qJ_{\mathcal G,L}^{(q)}
\simeq
-\mathscr E_{\mathcal G,L}^{[q]}
\]
with the tautological nullhomotopy of
$\mathscr E_{\mathcal G,L}^{[q]}$.  This nullhomotopes the image of the current in the
complete horizontal coefficient.  The homotopy-fibre universal property gives
\[
\boxed{
\mathcal U_{\mathcal G,L}^{(q)}:
\mathbf1
\longrightarrow
\operatorname{Trans}_{\mathcal G,H}^{1,q}(A;M)[n].
}
\]
\end{construction}

\begin{proposition}[Meaning of the native transgression]
\label{prop:symmetry-noether-meaning-of-the-native-transgression}
A point of $\operatorname{EulNull}_{\mathcal G}^{q}(L)$ together with
$\mathcal U_{\mathcal G,L}^{(q)}$ is precisely a complete-horizontal
transgression witness for the Noether current.  Paths and higher homotopies of
Euler nullifications map functorially to paths and higher homotopies of current
transgressions.
\end{proposition}

\begin{proof}
The transgression coefficient is the homotopy fibre of the horizontal-tail
inclusion.  A point of that fibre is exactly a point in the source together with
a nullhomotopy of its image.  Apply this to the current and the composite
nullhomotopy just constructed.
\end{proof}

\begin{proposition}[Complete transgression refines the selected one-stage conservation]
\label{prop:symmetry-noether-complete-transgression-refines-selected-one-stage-conservation}
The nullhomotopy defining
$\mathcal U_{\mathcal G,L}^{(q)}$ maps under the pulled-back local one-stage
restriction to the canonical one-stage on-shell conservation nullhomotopy of
Corollary~\ref{cor:symmetry-noether-canonical-one-stage-on-shell-conservation}.
Thus the complete-horizontal current transgression is a refinement of the
already selected one-stage conservation homotopy, not an independent
nullification.
\end{proposition}

\begin{proof}
The First Noether homotopy identifies the image of the current with
$-\sigma_{\mathcal C}$.  The tautological point of
$\operatorname{EulNull}_{\mathcal G}^{q}(L)$ supplies a nullhomotopy
$H:\sigma_{\mathcal C}\simeq0$ satisfying
\[
(u_{\mathcal C})_*H=\eta_{\mathcal C}.
\]
Hence the nullhomotopy used to define
$\mathcal U_{\mathcal G,L}^{(q)}$ has one-stage image
$-\eta_{\mathcal C}$.  By construction, $\eta_{\mathcal C}$ is the composite of
the one-stage identification
$u_{\mathcal C}\sigma_{\mathcal C}\simeq
c_{\mathcal C}^*E_{\mathcal G,L}^{(1)}$
with the critical nullhomotopy induced from
$c_{\mathcal C}^*E_L\simeq0$.  This is exactly the nullhomotopy used in the
canonical one-stage on-shell conservation corollary, with the same sign supplied
by the First Noether identity.
\end{proof}

\section{The critical formal symmetry groupoid}
\label{sec:symmetry-noether-the-critical-formal-symmetry-groupoid}

Let
\[
C_L:=\operatorname{Crit}(L),
\qquad
j_L:C_L\longrightarrow\overline A.
\]
Restriction of a formal groupoid along an arbitrary object-space morphism is
a groupoid construction, but unitwise formality is not invariant under
arbitrary object-space replacement.  We therefore separate raw restriction
from formalization.

\begin{construction}[Raw critical restriction groupoid]
\label{constr:symmetry-noether-raw-critical-restriction-groupoid}
For every $p\ge0$ define
\[
\mathcal H_{\mathcal G,L,p}
:=
\mathcal G_p
\times_{\overline A^{\times_Q(p+1)}}
C_L^{\times_Q(p+1)},
\]
where $\mathcal G_p\to\overline A^{\times_Q(p+1)}$ is the ordered tuple of
vertex maps.  Deletion and repetition of vertices give a simplicial object
\[
\mathcal H_{\mathcal G,L,\bullet}
\rightrightarrows
C_L.
\]
\end{construction}

\begin{proposition}[Raw critical restriction is a groupoid]
\label{prop:symmetry-noether-raw-critical-restriction-is-a-groupoid}
The simplicial object
$\mathcal H_{\mathcal G,L,\bullet}$ is a fixed-object internal groupoid over
$Q$, and the evident map
\[
(\mathcal H_{\mathcal G,L,\bullet},C_L)
\longrightarrow
(\mathcal G_\bullet,\overline A)
\]
is a functor of internal groupoids.
\end{proposition}

\begin{proof}
Each simplicial operator of $\mathcal G_\bullet$ commutes with its ordered
vertex maps.  Pulling the Segal diagrams of $\mathcal G_\bullet$ back along
$C_L^{\times_Q(\bullet+1)}\to
\overline A^{\times_Q(\bullet+1)}$ therefore gives the Segal diagrams of
$\mathcal H_{\mathcal G,L,\bullet}$.  Finite limits also transport the unit,
inversion, and their higher coherences.  Hence the raw restriction is a
groupoid object.  The projections to $\mathcal G_\bullet$ form the displayed
groupoid functor.
\end{proof}

\begin{construction}[Canonical critical formalization]
\label{constr:symmetry-noether-canonical-critical-formalization}
Apply the formalization coreflection of the formal-groupoid chapter to the
raw critical groupoid and define
\[
\boxed{
\mathcal K_{\mathcal G,L,\bullet}
:=
\widehat{\mathcal H_{\mathcal G,L,\bullet}}^{\,\mathrm{dR}/Q}.
}
\]
It comes with the formalization counit
\[
\kappa_{\mathcal G,L}:
\mathcal K_{\mathcal G,L,\bullet}
\longrightarrow
\mathcal H_{\mathcal G,L,\bullet}
\]
and hence with a canonical formal-groupoid morphism
\[
\mathcal K_{\mathcal G,L,\bullet}
\longrightarrow
\mathcal G_\bullet.
\]
\end{construction}

\begin{theorem}[Critical formal symmetry]
\label{thm:symmetry-noether-critical-formal-symmetry}
The groupoid
\[
\mathcal K_{\mathcal G,L,\bullet}
\rightrightarrows C_L
\]
is a formal Lie groupoid over $Q$.  Its canonical formal anchor gives a
contractibly unique simplicial morphism
\[
\phi_{\mathcal K}:
\mathcal K_{\mathcal G,L,\bullet}
\longrightarrow
R_{C_L/Q,\bullet}.
\]
The composite to $\mathcal G_\bullet$ is compatible with the two canonical
anchors.
\end{theorem}

\begin{proof}
Formalization is the right adjoint to the inclusion of formal fixed-object
groupoids and carries every fixed-object groupoid to a formal fixed-object
groupoid.  The maximal formal relation on $C_L/Q$ is terminal among formal
fixed-object groupoids, so it supplies the canonical anchor.  The composite
$\mathcal K_{\mathcal G,L,\bullet}\to\mathcal G_\bullet$ is a morphism
between formal groupoids.  Contractible uniqueness of the two anchors
identifies the two composites to the maximal ambient relation.
\end{proof}

\begin{corollary}[Recognition when raw critical restriction is already formal]
\label{cor:symmetry-noether-recognition-when-raw-critical-restriction-is-already-formal}
If the raw critical restriction
$\mathcal H_{\mathcal G,L,\bullet}$ is unitwise formal, then the counit
\[
\mathcal K_{\mathcal G,L,\bullet}
\longrightarrow
\mathcal H_{\mathcal G,L,\bullet}
\]
is an equivalence.
\end{corollary}

\begin{remark}[No critical-formality axiom]
\label{rem:symmetry-noether-no-critical-formality-axiom}
No condition asserting that the raw restriction is formal is retained as
input.  Formality is restored functorially by the canonical coreflection.
\end{remark}


\section{The symmetry-reduced critical object}
\label{sec:symmetry-noether-the-symmetry-reduced-critical-object}

\begin{definition}[Formal critical symmetry quotient]
\label{def:symmetry-noether-formal-critical-symmetry-quotient}
Define
\[
\boxed{
\operatorname{Crit}(L)//\mathcal G
:=
\left|
\mathcal K_{\mathcal G,L,\bullet}
\right|.
}
\]
Write
\[
q_{\mathcal G,L}:
C_L
\twoheadrightarrow
\operatorname{Crit}(L)//\mathcal G
\]
for the realization map.
\end{definition}

\begin{theorem}[Effective critical quotient]
\label{thm:symmetry-noether-effective-critical-quotient}
The map $q_{\mathcal G,L}$ is the effective quotient of the canonically
formalized critical symmetry groupoid.  Its \v{C}ech nerve is canonically
equivalent to $\mathcal K_{\mathcal G,L,\bullet}$.  The construction is
functorial in fixed-object morphisms of formal symmetry groupoids and commutes
with every raw spectral base change for which the formal-groupoid chapter
constructs base change.
\end{theorem}

\begin{proof}
Every groupoid object in an $\infty$-topos is effective.  Hence realization of
$\mathcal K_{\mathcal G,L,\bullet}$ is an effective epimorphism and its kernel
pair recovers the groupoid.  Formalization, finite-limit restriction, and
realization are functorial, and the formal-groupoid base-change theorem transports
all three constructions.
\end{proof}

\begin{construction}[Parameterized critical formal groupoid]
\label{constr:symmetry-noether-parameterized-critical-formal-groupoid}
Let
\[
P\longrightarrow Q
\]
be an arbitrary parameter object.  Define
\[
C_{L,P}:=C_L\times_QP
\]
and
\[
\boxed{
\mathcal K_{\mathcal G,L,P,\bullet}
:=
\mathcal K_{\mathcal G,L,\bullet}\times_QP.
}
\]
The formal-groupoid base-change theorem equips this simplicial object with its
base-changed formal Lie groupoid structure over $C_{L,P}/P$ and with the
base-changed canonical anchor.
\end{construction}

\begin{theorem}[Parameterized quotient and jet calculus]
\label{thm:symmetry-noether-parameterized-quotient-and-jet-calculus}
For every $P\to Q$ there are canonical equivalences
\[
\left|
\mathcal K_{\mathcal G,L,P,\bullet}
\right|
\simeq
\left(
\operatorname{Crit}(L)//\mathcal G
\right)\times_QP
\]
and the canonical Beck--Chevalley equivalence of comonads
\[
c_P^*J_{\mathcal K}
\simeq
J_{\mathcal K,P}c_P^*,
\]
where $c_P:C_L\times_QP\to C_L$ is the projection.  Pullback along
\[
q_{\mathcal G,L,P}:
C_L\times_QP
\twoheadrightarrow
\left(
\operatorname{Crit}(L)//\mathcal G
\right)\times_QP
\]
induces the effective descent equivalence
\[
E_{
(\operatorname{Crit}(L)//\mathcal G)\times_QP
}
\simeq
\operatorname{Coalg}_{J_{\mathcal K,P}}
\left(
E_{C_L\times_QP}
\right).
\]
\end{theorem}

\begin{proof}
Apply arbitrary base change of formal Lie groupoids to
$\mathcal K_{\mathcal G,L,\bullet}$ along $P\to Q$.  The theorem transports the
groupoid realization, quotient, disk--jet adjunction, jet comonad, and effective
descent equivalence by their canonical Beck--Chevalley maps.
\end{proof}

\begin{proposition}[No critical-formality or parameterized-descent datum]
\label{prop:symmetry-noether-no-critical-formality-or-parameterized-descent-datum}
The parameterized critical formal groupoid and its descent equivalence are
functorial outputs of base change.  No assertion that the raw critical restriction
is formal, no chosen quotient presentation, and no independent Cartan descent
datum on a parameter family is required.
\end{proposition}

\section{The affine-point linear tangent of the derived critical locus}
\label{sec:symmetry-noether-the-affine-point-linear-tangent-of-the-derived-critical-locus}

The native critical object belongs to the big spectral Nisnevich $\infty$-topos.
Chapter 12 nevertheless constructs, on the represented smooth scalar degree-zero
sector, a finite-locally-free affine-point linear coefficient.  We first use that
coefficient globally in the affine-point quasi-coherent category and only then pass
to ordinary tangent complexes on represented affine charts.

\begin{convention}[Represented smooth scalar critical domain]
\label{conv:symmetry-noether-represented-smooth-scalar-critical-domain}
Assume
\[
M=\mathcal O_Q^{\mathrm{add}},
\qquad
n=0,
\]
and assume that
\[
a:\overline A\longrightarrow Q
\]
is representably smooth in the sense of Chapter 12: after every represented affine
test $T\to Q$, the pullback
$\overline A_T:=T\times_Q\overline A\to T$ is a smooth spectral morphism.

Let
\[
\mathcal E
:=
L_{\overline A/Q}^{\mathrm{lin}}
\in
QCoh(\overline A)
\]
be the finite-locally-free affine-point linear refinement constructed there.  Its
underlying stable coefficient is
\[
U_{\overline A}(\mathcal E)
\simeq
V_{\mathcal O_Q^{\mathrm{add}}}^{(1)}(A).
\]
Write
\[
\mathcal T:=\mathcal E^\vee.
\]
No corresponding linear-space or tangent formula is asserted for a general stable
coefficient $M$.
\end{convention}

\begin{definition}[Affine-point linear-space object]
\label{def:symmetry-noether-affine-point-linear-space-object}
Define the native affine-point linear-space object
\[
\boxed{
\mathbb V_{\overline A}^{\mathrm{lin}}(\mathcal T)
:=
\Omega^\infty_{\overline A}
U_{\overline A}(\mathcal E).
}
\]
For every represented affine point
\[
c:T=\operatorname{Spec}B\longrightarrow\overline A
\]
its pullback is canonically represented by the ordinary relative linear spectral
space
\[
\mathbb V_T(c^*\mathcal T).
\]
\end{definition}

\begin{proof}
At $c$, affine-point recognition gives
\[
c^*U_{\overline A}(\mathcal E)
\simeq
U\Gamma(T,c^*\mathcal E).
\]
Since $c^*\mathcal E$ is finite locally free,
\[
\Gamma(T,c^*\mathcal E)
\simeq
\operatorname{map}_{\operatorname{QCoh}(T)}
(c^*\mathcal T,\mathcal O_T).
\]
The represented functor-of-points theorem for relative linear spaces identifies
the corresponding relative zeroth space with
$\mathbb V_T(c^*\mathcal T)$.  Naturality under affine restriction supplies the
Cartesian affine-point family defining the displayed native object.
\end{proof}

Under the preceding identification, the local Euler section determines a native
section
\[
e_L:
\overline A
\longrightarrow
\mathbb V_{\overline A}^{\mathrm{lin}}(\mathcal T).
\]

\begin{construction}[Affine-point Euler zero locus]
\label{constr:symmetry-noether-affine-point-euler-zero-locus}
Define
\[
\boxed{
C_L^{\mathrm{lin}}
:=
\overline A
\times_{\mathbb V_{\overline A}^{\mathrm{lin}}(\mathcal T)}
\overline A,
}
\]
where the second map is the zero section.
\end{construction}

\begin{theorem}[The native critical object is the affine-point Euler zero locus]
\label{thm:symmetry-noether-the-native-critical-object-is-the-affine-point-euler-zero-locus}
There is a canonical equivalence in the native big spectral Nisnevich
$\infty$-topos
\[
\boxed{
\operatorname{Crit}(L)
\simeq
C_L^{\mathrm{lin}}.
}
\]
\end{theorem}

\begin{proof}
The native one-stage critical object is the geometric zero locus of
\[
E_L:
\mathbf1_{\overline A}
\longrightarrow
V_{\mathcal O_Q^{\mathrm{add}}}^{(1)}(A).
\]
The affine-point linear refinement identifies the target with
$U_{\overline A}(\mathcal E)$ and hence identifies its relative zeroth-space object
with
$\mathbb V_{\overline A}^{\mathrm{lin}}(\mathcal T)$.
Under this identification the Euler and zero sections are $e_L$ and the displayed
zero section.  The two homotopy fibres are therefore the same finite limit.
\end{proof}

\begin{construction}[Affine-point Euler linearization]
\label{constr:symmetry-noether-affine-point-euler-linearization}
Let
\[
j_L:C_L\longrightarrow\overline A
\]
be the critical projection.  At an affine critical point
$c:T=\operatorname{Spec}B\to C_L$, the universal split square-zero derivative of
the represented Euler section gives a $B$-linear morphism
\[
de_{L,c}:
c^*\mathcal T
\longrightarrow
c^*\mathcal E.
\]
The Chapter 12 affine restriction and base-change compatibilities make these maps
Cartesian in $c$.  Hence they assemble to a morphism in $QCoh(C_L)$
\[
\boxed{
de_L^{\mathrm{lin}}:
j_L^*\mathcal T
\longrightarrow
j_L^*\mathcal E.
}
\]
Define the affine-point linear critical tangent coefficient by
\[
\boxed{
T_{C_L/Q}^{\mathrm{lin}}
:=
\operatorname{fib}
\left(
de_L^{\mathrm{lin}}
\right)
\in QCoh(C_L).
}
\]
\end{construction}

\begin{theorem}[Affine-point Jacobi identification]
\label{thm:symmetry-noether-affine-point-jacobi-identification}
On the represented smooth scalar Jacobi overlap,
\[
de_L^{\mathrm{lin}}
\simeq
\operatorname{Jac}_L^{\mathrm{lin}}
\]
in $QCoh(C_L)$.  Consequently
\[
\boxed{
T_{C_L/Q}^{\mathrm{lin}}
\simeq
\operatorname{fib}
\left(
\operatorname{Jac}_L^{\mathrm{lin}}
\right).
}
\]
Applying the exact underlying-coefficient functor gives
\[
U_{C_L}
\left(
T_{C_L/Q}^{\mathrm{lin}}
\right)
\simeq
\operatorname{fib}(\operatorname{Jac}_L)
=:J_L.
\]
\end{theorem}

\begin{proof}
At every affine critical point, Chapter 12 identifies the affine linear Jacobi map
with the first split square-zero linearization of the brave-new Euler source.
The map $de_{L,c}$ is characterized by that same universal derivative.
The identifications commute with affine restriction, so they assemble to the
displayed equality in the affine-point quasi-coherent category.  Exactness of
$U_{C_L}$ preserves the fibre.
\end{proof}

\subsection{Ordinary tangent complexes on represented affine charts}
\label{subsec:symmetry-noether-ordinary-tangent-complexes-on-represented-affine-charts}

Let
\[
t:T=\operatorname{Spec}R\longrightarrow Q
\]
be a represented affine test and put
\[
\overline A_T:=T\times_Q\overline A,
\qquad
C_{L,T}:=T\times_QC_L.
\]
By the theorem-domain hypothesis,
$\overline A_T\to T$ is a smooth spectral morphism.  The pullback
$\mathcal E_T$ is finite locally free and
\[
\mathbb V_{\overline A}^{\mathrm{lin}}(\mathcal T)
\times_QT
\simeq
\mathbb V_{\overline A_T}(\mathcal T_T).
\]
Thus $C_{L,T}$ is the ordinary derived zero locus of the represented Euler section
on the spectral scheme $\overline A_T$.

\begin{proposition}[Chartwise quasi-smooth Euler zero locus]
\label{prop:symmetry-noether-chartwise-quasi-smooth-euler-zero-locus}
The projection
\[
j_{L,T}:C_{L,T}\longrightarrow\overline A_T
\]
is a quasi-smooth closed immersion and
\[
L_{C_{L,T}/\overline A_T}
\simeq
j_{L,T}^*\mathcal E_T^\vee[1].
\]
Cotangent transitivity gives
\[
j_{L,T}^*L_{\overline A_T/T}
\longrightarrow
L_{C_{L,T}/T}
\longrightarrow
j_{L,T}^*\mathcal E_T^\vee[1].
\]
\end{proposition}

\begin{proof}
This is the Chapter 2 derived-zero-locus theorem for a section of the
finite-locally-free represented linear space
$\mathbb V_{\overline A_T}(\mathcal T_T)$, followed by cotangent transitivity.
\end{proof}

\begin{theorem}[Chartwise ordinary critical tangent]
\label{thm:symmetry-noether-chartwise-ordinary-critical-tangent}
All terms in the preceding cotangent triangle are perfect.  Exact duality gives
\[
T_{C_{L,T}/T}
\longrightarrow
j_{L,T}^*T_{\overline A_T/T}
\xrightarrow{\,de_{L,T}\,}
j_{L,T}^*\mathcal E_T
\]
and hence
\[
\boxed{
T_{C_{L,T}/T}
\simeq
\operatorname{fib}(de_{L,T})
\simeq
\operatorname{fib}(\operatorname{Jac}_{L,T}).
}
\]
Moreover this ordinary tangent complex is the pullback of
$T_{C_L/Q}^{\mathrm{lin}}$ to the chart under the affine-point recognition
equivalence.
\end{theorem}

\begin{proof}
Dualize the perfect cotangent triangle.  The dual conormal morphism is the
ordinary derivative of the represented Euler section.  By the affine-point Jacobi
theorem that derivative is the represented Jacobi morphism.  The final statement
follows because every construction is the affine restriction of
$de_L^{\mathrm{lin}}$ and its fibre.
\end{proof}

\begin{convention}[Ordinary tangent notation]
\label{conv:symmetry-noether-ordinary-tangent-notation}
The notation
\[
T_{C_L/Q}
\]
is used only after passage to a represented chart on which the ordinary tangent
complex has been constructed.  Globally in the native topos the intrinsic
linearized coefficient is
$T_{C_L/Q}^{\mathrm{lin}}$.
\end{convention}

\section{Critical symmetry directions are Jacobi zero-modes}
\label{sec:symmetry-noether-critical-symmetry-directions-are-jacobi-zero-modes}

Work on the overlap of the represented first-order groupoid domain with a
represented affine critical chart
\[
T\longrightarrow Q
\]
as in the preceding section.  Then
\[
C_{L,T}
\]
and the raw critical arrow object
\[
\mathcal H_{\mathcal G,L,1,T}
=
\mathcal H_{\mathcal G,L,1}\times_QT
\]
are represented spectral schemes.

Let
\[
a_{\mathcal K,T}^\vee
:=
u^*
L_{\mathcal H_{\mathcal G,L,1,T}/C_{L,T},s}.
\]
By the represented first-neighbourhood theorem for unitwise formalization, this
coefficient corepresents the split first-order probe of
$\mathcal K_{\mathcal G,L,\bullet}\times_QT$ even if the completed arrow itself is
not represented.  Define its internal dual by
\[
a_{\mathcal K,T}:=(a_{\mathcal K,T}^\vee)^\vee.
\]
Linearization of the canonical critical formal anchor constructs
\[
\rho_{\mathcal K,T}^{\mathrm{rep}}:
a_{\mathcal K,T}
\longrightarrow
T_{C_{L,T}/T}.
\]
No finite-local-freeness hypothesis on $a_{\mathcal K,T}^\vee$ is required for
this internal-dual morphism.  The perfectness already available for the critical
tangent target is used only in the subsequent Jacobi and bilinear identifications.

\begin{theorem}[Jacobi zero-mode factorization]
\label{thm:symmetry-noether-jacobi-zero-mode-factorization}
There is a canonical factorization
\[
\boxed{
a_{\mathcal K,T}
\xrightarrow{\rho_{\mathcal K,T}^{\mathrm{rep}}}
T_{C_{L,T}/T}
\xrightarrow[\sim]{}
\operatorname{fib}(\operatorname{Jac}_{L,T})
\longrightarrow
j_{L,T}^*T_{\overline A_T/T}.
}
\]
Equivalently,
\[
\operatorname{Jac}_{L,T}
\circ
\left(
a_{\mathcal K,T}
\to
j_{L,T}^*T_{\overline A_T/T}
\right)
\simeq0
\]
with the tautological fibre nullhomotopy.
\end{theorem}

\begin{proof}
The base-changed critical formalization is a formal symmetry groupoid on
$C_{L,T}/T$.  Its represented first-order theorem gives the tangent anchor into
$T_{C_{L,T}/T}$.  The chartwise critical-tangent theorem identifies this tangent
complex with the homotopy fibre of the represented Jacobi operator.  Composition
with the fibre inclusion is the ambient critical field variation and carries the
canonical Jacobi nullhomotopy.
\end{proof}

\begin{remark}[Why the formalized critical groupoid is used]
\label{rem:symmetry-noether-why-the-formalized-critical-groupoid-is-used}
No unconditional map
\[
j_L^*a_{\mathcal G}
\longrightarrow
\operatorname{fib}(\operatorname{Jac}_L)
\]
is asserted.  The raw restriction of $\mathcal G$ may fail to be formal and an
arbitrary ambient symmetry arrow through a critical point need not define a critical
formal arrow.  The object
$\mathcal K_{\mathcal G,L,\bullet}$ constructs the formal critical symmetry
relation first; its represented chart then supplies the tangent anchor used in the
theorem.
\end{remark}


\section{Symmetry degeneracy of the critical Hessian}
\label{sec:symmetry-noether-symmetry-degeneracy-of-the-critical-hessian}

Remain on a represented smooth scalar degree-zero critical chart
$C_{L,T}\to T$.  The critical Hessian
\[
\operatorname{Hess}_{L,T}^{\mathrm{crit}}
\]
is the symmetric bilinear morphism whose adjoint is
$\operatorname{Jac}_{L,T}$.

\begin{definition}[Homotopy radical of the critical Hessian]
\label{def:symmetry-noether-homotopy-radical-of-the-critical-hessian}
Define
\[
\operatorname{Rad}
\left(
\operatorname{Hess}_{L,T}^{\mathrm{crit}}
\right)
\]
to be the homotopy fibre over zero of the adjoint contraction
\[
j_{L,T}^*T_{\overline A_T/T}
\longrightarrow
j_{L,T}^*\mathcal E_T.
\]
\end{definition}

\begin{corollary}[Critical symmetry-radical factorization]
\label{cor:symmetry-noether-critical-symmetry-radical-factorization}
The Jacobi zero-mode theorem gives a canonical morphism
\[
\boxed{
a_{\mathcal K,T}
\longrightarrow
\operatorname{Rad}
\left(
\operatorname{Hess}_{L,T}^{\mathrm{crit}}
\right).
}
\]
\end{corollary}

\begin{proof}
The adjoint contraction of the critical Hessian is
$\operatorname{Jac}_{L,T}$ under the finite-locally-free duality of the represented
critical Hessian.  Apply the Jacobi nullhomotopy and the universal property of its
homotopy fibre.
\end{proof}

\section{The mixed symmetry--ambient-vertical--horizontal calculus}
\label{sec:symmetry-noether-the-mixed-symmetry-ambient-vertical-horizontal-calculus}

The symmetry--horizontal calculus has symmetry and horizontal directions, while the
presymplectic transgression of the variational chapter is built from the ambient
vertical relation.  To differentiate a symmetry current vertically and compare it
with presymplectic contraction, both constructions must first be placed in one
coefficient system.  We construct that common system directly from the \v{C}ech
presentation of the maximal relative infinitesimal relation.  This avoids any
variance ambiguity in an abstract simplicial cotensor.

Put
\[
V_\bullet:=R_{\overline A/Q,\bullet},
\qquad
G_A:=Q_{\overline A/Q}.
\]
Thus
\[
V_p\simeq
\underbrace{
\overline A\times_{G_A}\cdots\times_{G_A}\overline A
}_{p+1\ \mathrm{factors}}.
\]

\begin{construction}[Ambient rectangular \v{C}ech grid]
\label{constr:symmetry-noether-ambient-rectangular-relation-grid}
For $p,r\ge0$ define
\[
\boxed{
\operatorname{Rect}_V(p,r)
:=
\underbrace{
\overline A\times_{G_A}\cdots\times_{G_A}\overline A
}_{(p+1)(r+1)\ \mathrm{factors}}.
}
\]
Index the factors by
\[
(i,j)\in\{0,\ldots,p\}\times\{0,\ldots,r\}.
\]
For fixed $j$, projection to the row
\[
(0,j),\ldots,(p,j)
\]
gives a morphism to $V_p$.  Together the $r+1$ rows give
\[
\operatorname{row}_{p,r}:
\operatorname{Rect}_V(p,r)
\longrightarrow
\underbrace{
V_p\times_{G_A}\cdots\times_{G_A}V_p
}_{r+1\ \mathrm{factors}}.
\]
\end{construction}

\begin{lemma}[Rectangular functoriality]
\label{lem:symmetry-noether-rectangular-functoriality}
Deleting or repeating the $i$-coordinate and deleting or repeating the $j$-coordinate
make
\[
([p],[r])\longmapsto\operatorname{Rect}_V(p,r)
\]
a bisimplicial object.  The two operations commute with their full simplicial
coherence, and
\[
\operatorname{Rect}_V(p,0)\simeq V_p,
\qquad
\operatorname{Rect}_V(0,r)\simeq V_r.
\]
\end{lemma}

\begin{proof}
Every simplicial operator of the \v{C}ech nerve is deletion or repetition of one
of the ordered vertices.  Applying an operator in the first coordinate acts on
each row, while applying one in the second coordinate deletes or repeats entire
rows.  These operations commute because they act on different indices of the
finite grid.  The axis identities are the cases in which the grid has one row or
one column.
\end{proof}

\begin{definition}[Symmetry--vertical mixed grid]
\label{def:symmetry-noether-symmetry-vertical-mixed-grid}
The canonical anchor
\[
\phi_{\mathcal G,p}:\mathcal G_p\longrightarrow V_p
\]
and the structural map $V_p\to G_A$ give
\[
(\mathcal G_p)^{\times_{G_A}(r+1)}
\longrightarrow
(V_p)^{\times_{G_A}(r+1)}.
\]
Define
\[
\boxed{
M_{\mathcal G,V}^{p,r}
:=
\operatorname{Rect}_V(p,r)
\times_{(V_p)^{\times_{G_A}(r+1)}}
(\mathcal G_p)^{\times_{G_A}(r+1)}.
}
\]
Thus every row of the ambient infinitesimal rectangle is lifted to a symmetry
$p$-simplex, while the common $G_A$-image of the complete rectangle is retained.
\end{definition}

\begin{theorem}[Bisimplicial mixed relation]
\label{thm:symmetry-noether-bisimplicial-mixed-relation}
The assignment
\[
([p],[r])\longmapsto M_{\mathcal G,V}^{p,r}
\]
is canonically bisimplicial.  Its two axes are
\[
M_{\mathcal G,V}^{p,0}\simeq\mathcal G_p,
\qquad
M_{\mathcal G,V}^{0,r}\simeq V_r.
\]
\end{theorem}

\begin{proof}
The first simplicial direction acts simultaneously on every ambient row and on
every lifted symmetry row.  The second deletes or repeats rows.  The anchor is
simplicial, so the defining pullback is preserved.  For $r=0$ the row map has one
factor and the pullback is
\[
V_p\times_{V_p}\mathcal G_p\simeq\mathcal G_p.
\]
For $p=0$, one has $\mathcal G_0=\overline A=V_0$, and both the ambient and lifted
row products are
\[
\overline A^{\times_{G_A}(r+1)}\simeq V_r.
\]
Hence the other axis is $V_r$.
\end{proof}

\begin{definition}[Symmetry--vertical--horizontal triple relation]
\label{def:symmetry-noether-symmetry-vertical-horizontal-triple-relation}
For $p,r,q\ge0$ define
\[
R_{\mathcal G,V,H}^{p,r,q}(A)
:=
H_q\times_QM_{\mathcal G,V}^{p,r}.
\]
This is a trisimplicial relation object.  For $M\in E_Q$ put
\[
C_{\mathcal G,V,H}^{p,r,q}(A;M)
:=
\Pi_{R_{\mathcal G,V,H}^{p,r,q}\to Q}
M|_{R_{\mathcal G,V,H}^{p,r,q}}.
\]
Normalize in the three directions and write
\[
\Omega_{\mathcal G,V,H}^{p,r,q}(A;M)
:=
N_{\mathcal G}^pN_V^rN_H^q
C_{\mathcal G,V,H}^{\bullet,\bullet,\bullet}(A;M).
\]
\end{definition}

\begin{theorem}[Coherent mixed tricomplex]
\label{thm:symmetry-noether-coherent-mixed-tricomplex}
Stable Dold--Kan applied successively in the three functor categories gives
coherent structure maps
\[
d_{\mathcal G},\qquad d_V^{\mathrm{mix}},\qquad d_H
\]
whose squares carry specified nullhomotopies and whose pairwise composites carry
the canonical interchange homotopies.  After the Koszul rephasing
\[
d_{\mathcal G}=\partial_{\mathcal G},
\qquad
d_V^{\mathrm{mix}}=(-1)^p\partial_V,
\qquad
d_H=(-1)^{p+r}\partial_H,
\]
their homotopy-category shadow satisfies pairwise signed anticommutation.  All
higher cube coherences are inherited from the product-indexed coherent
tricochain object.
\end{theorem}

\begin{proof}
The three simplicial directions commute before stable coefficients are taken.
Normalization is a finite matching-fibre construction, and totalization is a
limit in a stable functor category.  Thus all orders of normalization and partial
totalization are canonically related by Fubini.  Stable Dold--Kan is natural for
exact functors, producing a coherent cochain object on the product of the three
pointed cochain indexing categories.  The displayed signs are the standard Koszul
rephasing by the number of preceding degrees.
\end{proof}

For fixed $(p,r)$ horizontally totalize and write
\[
\operatorname{Sec}_{\mathcal G,V}^{p,r}(A;M)
:=
\operatorname{Tot}_H
N_{\mathcal G}^pN_V^r
C_{\mathcal G,V,H}^{\bullet,\bullet,\bullet}(A;M).
\]
Define the horizontal tail
$F_H^q\operatorname{Sec}_{\mathcal G,V}^{p,r}$ and
\[
\operatorname{Trans}_{\mathcal G,V,H}^{p,r,q}
:=
\operatorname{fib}
\left(
F_H^q\operatorname{Sec}_{\mathcal G,V}^{p,r}
\longrightarrow
\operatorname{Sec}_{\mathcal G,V}^{p,r}
\right).
\]
The axis equivalence gives
\[
\operatorname{Sec}_{\mathcal G,V}^{p,0}
\simeq
\operatorname{Sec}_{\mathcal G}^{p}.
\]

\begin{construction}[Ordered staircase]
\label{constr:symmetry-noether-ordered-staircase}
For $(p,r)=(1,1)$ the rectangular grid has vertices
\[
(0,0),(1,0),(0,1),(1,1).
\]
Projection to the ordered staircase
\[
(0,0)\longrightarrow(1,0)\longrightarrow(1,1)
\]
gives
\[
\operatorname{Rect}_V(1,1)
\longrightarrow
\overline A\times_{G_A}\overline A\times_{G_A}\overline A
\simeq V_2.
\]
Restricting to the symmetry-row pullback defines
\[
\boxed{
\tau_{\mathcal G,V}:
M_{\mathcal G,V}^{1,1}
\longrightarrow
V_2.
}
\]
\end{construction}

\begin{lemma}[Normalization of the ordered contraction]
\label{lem:symmetry-noether-normalization-of-the-ordered-contraction}
Let $\omega$ be a normalized ambient vertical degree-two cochain.  Its pullback
along $\tau_{\mathcal G,V}$ has zero restriction under the symmetry codegeneracy
and under the mixed-vertical codegeneracy.  Hence there is a canonical
factorization
\[
N_V^2C_V^\bullet
\longrightarrow
N_{\mathcal G}^1N_V^1
C_{\mathcal G,V}^{\bullet,\bullet}.
\]
The factorization is natural in the cochain and in the formal symmetry groupoid.
\end{lemma}

\begin{proof}
If the symmetry direction is degenerate, the first two vertices of the ordered
staircase coincide, so the resulting ambient $2$-simplex is degenerate.  If the
mixed-vertical direction is degenerate, the final two vertices coincide.  A
normalized ambient degree-two cochain vanishes on either degeneracy.  Normalization
is the simultaneous fibre of the two codegeneracy maps, so the required
factorization follows from its universal property.  Naturality is termwise.
\end{proof}

\begin{construction}[Ordered mixed presymplectic contraction]
\label{constr:symmetry-noether-ordered-mixed-presymplectic-contraction}
Pull the ambient two-vertical presymplectic transgression
$\omega_L^{(q)}$ along $\tau_{\mathcal G,V}$.  The preceding lemma gives the mixed
normalization, and its existing horizontal transgression nullhomotopy is preserved
by exact pullback.  Thus there is a canonical section
\[
\boxed{
\iota_{\mathcal G}\omega_L^{(q)}
:
\mathbf1
\longrightarrow
\operatorname{Trans}_{\mathcal G,V,H}^{1,1,q}(A;M)[n].
}
\]
This is a complete relation-level construction.  No tangent bundle, interior
product, or multilinear model is used in its definition.
\end{construction}

\subsection{Exact represented $(1,1)$-linearization of the ordered contraction}
\label{subsec:symmetry-noether-exact-represented-one-one-linearization}

The complete mixed coefficient is retained before any linearization.  We now
construct a separate exact first-order image on the represented smooth scalar
sector.  This image depends only on two first-order side directions because that
dependence is built by a reduced two-variable stabilization; it is not an
identification of the unlinearized mixed coefficient with a bilinear form.

Let
\[
c:T=\operatorname{Spec}B\longrightarrow\overline A
\]
be a represented affine centre and retain the affine cotangent coefficient $N_c$
which classifies split first variations.  Let $P,Q$ be connective
$B$-modules and let
\[
\alpha:N_c\longrightarrow P,
\qquad
\beta:N_c\longrightarrow Q
\]
classify two split first variations.  Suppose the $P$-variation is equipped with a
symmetry lift
\[
\widetilde\sigma_\alpha:T[P]\longrightarrow\mathcal G_1.
\]
Write
\[
T[P,Q]
:=
\operatorname{Spec}
\left((B\oplus P)\otimes_B(B\oplus Q)\right)
\]
with projections $p_P,p_Q$ and central retraction $r_{P,Q}:T[P,Q]\to T$.

\begin{construction}[Ordered side evaluation]
\label{constr:symmetry-noether-ordered-side-evaluation}
Let $u_\alpha:T[P]\to\overline A$ and $u_\beta:T[Q]\to\overline A$ be the two
classified field variations.  On $T[P,Q]$ take the first symmetry row to be
\[
\widetilde\sigma_\alpha p_P:T[P,Q]\longrightarrow\mathcal G_1
\]
and the second symmetry row to be the groupoid unit at $u_\beta p_Q$.
Their four anchor vertices are
\[
cr_{P,Q},\qquad
u_\alpha p_P,\qquad
u_\beta p_Q,\qquad
u_\beta p_Q.
\]
The central section $T\hookrightarrow T[P,Q]$ is a finite composite of split
square-zero thickenings and is inverted by relative de Rham localization.  Hence
the four vertices have a canonical coherent common image in $G_A$.  They therefore
define a canonical rectangle point, and the two specified groupoid rows lift it to
\[
\boxed{
\operatorname{side}_{\widetilde\sigma_\alpha,\beta}:
T[P,Q]
\longrightarrow
M_{\mathcal G,V}^{1,1}.
}
\]
The ordered staircase of this point has vertices
\[
cr_{P,Q},\qquad
u_\alpha p_P,\qquad
u_\beta p_Q.
\]
\end{construction}

For $\epsilon,\delta\in\{0,1\}$ switch the $P$-direction or the $Q$-direction off
when the corresponding index is zero.  When the $P$-direction is switched off,
use the groupoid unit as the symmetry lift.  Pulling a mixed section through the
four ordered side evaluations gives a pointed square of stable sections.

\begin{construction}[Exact $(1,1)$ side linearization]
\label{constr:symmetry-noether-exact-one-one-side-linearization}
Take the Boolean total fibre of the preceding square and then apply reduced
stabilization successively in the independent $P$- and $Q$-variables.  Exactness of
stabilization makes the two orders canonically equivalent.  Denote the resulting
exact biderivative by
\[
\operatorname{Lin}_{1,1}^{\mathrm{side}}
\]
and the image of the ordered mixed contraction by
\[
\lambda_{\mathcal G,V,c}^{(1,1)}
\left(\iota_{\mathcal G}\omega_L^{(q)}\right).
\]
After passage to the stabilized direction categories the biderivative is exact
separately in the $P$- and $Q$-variables.  Its evaluated section is natural in the
two derivations and in the chosen symmetry lift, and the whole construction is
natural under represented affine restriction.  Applying the identical construction
with the maximal relation groupoid $V_\bullet$ and its identity anchor defines
\[
\lambda_{V,c}^{(1,1)}(\omega)
\]
for every ambient vertical degree-two section $\omega$.  In that case the first row
is the canonical maximal-relation edge associated to the first field variation, so
no symmetry-lift choice is involved.  The same notation $\lambda_{V^C,c}^{(1,1)}$
is used for the maximal critical relation.
\end{construction}

\begin{lemma}[Side-coefficient and evaluated-section separation]
\label{lem:symmetry-noether-side-coefficient-and-evaluated-section-separation}
Fix the represented affine centre $c:T=\operatorname{Spec}B\to\overline A$ and
connective $B$-modules $P,Q$.  Before a pair of field derivations, or a
symmetry lift of the first derivation, is chosen, the four target coefficients of
the ordered side-evaluation square are canonically determined by $c$, $P$, and $Q$.
Write
\[
\mathscr E_{\omega,c}^{q,\square}(P,Q)
\]
for that square of stable target coefficients.  A choice of
$\alpha:N_c\to P$, $\beta:N_c\to Q$, and, on the symmetry branch, a lift
$\widetilde\sigma_\alpha$, determines a section of this already constructed
coefficient square.  Changing those choices changes the evaluated section but not
the coefficient square.  The assertion remains true for the ambient maximal
relation and for the maximal critical relation.
\end{lemma}

\begin{proof}
Every vertex of the split parameter square has structural morphism to $Q$ obtained
from its split retraction to $T$ followed by the fixed centre and the structural
map to $Q$.  The stable scalar coefficient is pulled back from $Q$, so its four
vertex targets and all restriction maps depend only on the split parameter algebras
attached to $B,P,Q$.  The maps to the field object and to the symmetry relation are
used only to pull the already constructed mixed section into those targets.  Thus
$\alpha$, $\beta$, and $\widetilde\sigma_\alpha$ affect the value of the section,
not its target coefficient.  Normalization, horizontal totalization, horizontal
tails, transgression fibres, Boolean total fibres, and reduced stabilization are
functorial stable limit constructions, so this separation survives through the
exact side linearization.  The maximal and maximal-critical cases are the same
argument with the canonical relation edge in place of a chosen symmetry lift.
\end{proof}

\begin{theorem}[Represented first-order contraction recognition]
\label{thm:symmetry-noether-represented-first-order-contraction-recognition}
On the represented smooth scalar degree-zero overlap, let $\xi$ be a represented
first-order symmetry direction and let $\eta$ be an arbitrary represented vertical
first-order direction.  Then the exact side linearization satisfies the canonical
identity
\[
\boxed{
\lambda_{\mathcal G,V,c}^{(1,1)}
\left(\iota_{\mathcal G}\omega_L^{(q)}\right)(\xi,\eta)
\simeq
\lambda_{V,c}^{(1,1)}
\left(\omega_L^{(q)}\right)
\left(\rho_{\mathcal G}^{\mathrm{rep}}(\xi),\eta\right).
}
\]
Here the right side denotes the same exact two-side first-order extraction formed in
the ambient maximal relation.  No equality of raw mixed cochains with an ordinary
bilinear form is asserted.
\end{theorem}

\begin{proof}
The represented first-order symmetry theorem identifies the first-order image of
$\widetilde\sigma_\alpha$ under the nonlinear anchor with the field variation
classified by $\rho_{\mathcal G}^{\mathrm{rep}}(\xi)$.  The ordered side rectangle
was defined by applying that anchor to the first row and using a degenerate unit
row at the second field variation.  Hence its ordered staircase maps to the same
ambient relation triangle as the corresponding pair of field-side variations.
This identity is natural before the Boolean total fibre is taken.  Taking the
reduced two-variable cross-effect and stabilizing in both variables therefore
preserves it and gives the displayed exact first-order equality.
\end{proof}

\subsection{Critical rectangular mixed calculus}
\label{subsec:symmetry-noether-critical-rectangular-mixed-calculus}

The critical formal groupoid changes the object space from $\overline A$ to
$C_L$.  Its mixed calculus is therefore reconstructed from the maximal
infinitesimal relation of $C_L/Q$ rather than obtained by declaring an arbitrary
object-space replacement to be a base change of the ambient rectangle.

Put
\[
G_C:=Q_{C_L/Q},
\qquad
V^C_\bullet:=R_{C_L/Q,\bullet}.
\]
Thus
\[
V^C_p\simeq
\underbrace{
C_L\times_{G_C}\cdots\times_{G_C}C_L
}_{p+1\ \mathrm{factors}}.
\]
For $p,r\ge0$ define
\[
\operatorname{Rect}_{V^C}(p,r)
:=
\underbrace{
C_L\times_{G_C}\cdots\times_{G_C}C_L
}_{(p+1)(r+1)\ \mathrm{factors}}
\]
and use the canonical critical anchor
\[
\phi_{\mathcal K,p}:
\mathcal K_{\mathcal G,L,p}\longrightarrow V^C_p
\]
to set
\[
\boxed{
M_{\mathcal K,V^C}^{p,r}
:=
\operatorname{Rect}_{V^C}(p,r)
\times_{(V^C_p)^{\times_{G_C}(r+1)}}
(\mathcal K_{\mathcal G,L,p})^{\times_{G_C}(r+1)}.
}
\]

\begin{proposition}[Critical mixed-grid functoriality]
\label{prop:symmetry-noether-critical-mixed-grid-functoriality}
The assignment
\[
([p],[r])\longmapsto M_{\mathcal K,V^C}^{p,r}
\]
is bisimplicial, with axes
\[
M_{\mathcal K,V^C}^{p,0}\simeq\mathcal K_{\mathcal G,L,p},
\qquad
M_{\mathcal K,V^C}^{0,r}\simeq V^C_r.
\]
Moreover the formal-groupoid morphism
\[
(\mathcal K_{\mathcal G,L,\bullet},C_L)
\longrightarrow
(\mathcal G_\bullet,\overline A)
\]
and functoriality of the relative de Rham effective quotients induce a canonical map
of bisimplicial mixed grids
\[
\boxed{
M_{\mathcal K,V^C}^{\bullet,\bullet}
\longrightarrow
M_{\mathcal G,V}^{\bullet,\bullet}.
}
\]
\end{proposition}

\begin{proof}
The bisimplicial structure is the same row-and-column deletion/repetition
construction as above.  The morphism $C_L\to\overline A$ induces
$G_C\to G_A$ and therefore sends every critical rectangle to the corresponding
ambient rectangle.  The formal-groupoid functor sends each lifted critical row to
a lifted ambient row.  Compatibility of the canonical anchors identifies the two
maps to the ambient row product, so the pullback universal property gives the
mixed-grid morphism with all simplicial coherence.
\end{proof}

For $p,r,q\ge0$ put
\[
R_{\mathcal K,V^C,H}^{p,r,q}(A)
:=
H_q\times_QM_{\mathcal K,V^C}^{p,r}.
\]
Stable restriction along the preceding grid morphism gives a contravariant map from
the ambient mixed tricochain object to the critical one.


\begin{theorem}[Critical mixed grid is the symmetry--\v{C}ech resolution of the reduced infinitesimal relation]
\label{thm:symmetry-noether-critical-mixed-grid-is-cech-resolution}
Put
\[
R_L:=\operatorname{Crit}(L)//\mathcal G,
\qquad
q_L:=q_{\mathcal G,L}:C_L\twoheadrightarrow R_L,
\]
and let
\[
G_R:=Q_{R_L/Q},
\qquad
V^R_r:=R_{R_L/Q,r}.
\]
Then the formal quotient map $q_L$ induces a canonical equivalence
\[
\boxed{G_C\xrightarrow{\ \sim\ }G_R.}
\]
Under this equivalence, for every $r\ge0$ the vertexwise quotient map
\[
q_r:V^C_r\longrightarrow V^R_r
\]
is an effective epimorphism and there is a canonical bisimplicial equivalence
\[
\boxed{
M_{\mathcal K,V^C}^{p,r}
\simeq
\check C_p(q_r).
}
\]
Consequently
\[
\left|M_{\mathcal K,V^C}^{\bullet,r}\right|
\simeq
V^R_r.
\]
\end{theorem}

\begin{proof}
Because $q_L$ is the effective quotient of a formal Lie groupoid, the formal-quotient
recognition theorem gives a relative de Rham equivalence
\[
(C_L)_{\mathrm{dR}/Q}
\xrightarrow{\ \sim\ }
(R_L)_{\mathrm{dR}/Q}.
\]
The effective-image factorizations of the two relative de Rham units therefore fit
into the diagram
\begin{tikzcd}[column sep=large,row sep=large]
C_L
  \arrow[r,two heads,"q_L"]
  \arrow[d,two heads,"\epsilon_C"'] &
R_L
  \arrow[d,two heads,"\epsilon_R"] \\
G_C
  \arrow[r,"\sim"']
  \arrow[d,hook] &
G_R
  \arrow[d,hook] \\
(C_L)_{\mathrm{dR}/Q}
  \arrow[r,"\sim"'] &
(R_L)_{\mathrm{dR}/Q}.
\end{tikzcd}
Uniqueness of effective-epimorphism--monomorphism factorization gives the middle
equivalence.  Hence
\[
V^C_r\simeq C_L^{\times_{G_C}(r+1)},
\qquad
V^R_r\simeq R_L^{\times_{G_C}(r+1)}.
\]
The map $q_r$ is an effective epimorphism by pullback and finite-product stability.
Its $p$th \v{C}ech term is an array of $(p+1)(r+1)$ points of $C_L$ with a common
$G_C$-image and with each fixed vertical vertex-column having a common image in
$R_L$.  Reindexing the same finite fibre product by rows identifies this array
exactly with the defining pullback for $M_{\mathcal K,V^C}^{p,r}$.  Deletion and
repetition of rows and columns agree under this reindexing, so the equivalence is
bisimplicial.  Effectivity of $q_r$ gives the realization formula.
\end{proof}

\begin{theorem}[Descent of the complete critical vertical--horizontal calculus]
\label{thm:symmetry-noether-descent-of-complete-critical-vertical-horizontal-calculus}
For $r,s\ge0$ define the reduced stable cochain coefficient
\[
C_{\mathrm{red}}^{r,s}(A;M)
:=
\Pi_{H_s\times_QV^R_r\to Q}
M|_{H_s\times_QV^R_r}.
\]
Then there is a canonical equivalence of bicosimplicial stable coefficient objects
\[
\boxed{
C_{\mathrm{red}}^{\bullet,\bullet}(A;M)
\simeq
\operatorname{Tot}_{\mathcal K}
C_{\mathcal K,V^C,H}^{\bullet,\bullet,\bullet}(A;M),
}
\]
where totalization is taken in the symmetry direction before vertical and
horizontal normalization.  The equivalence is compatible with vertical and
horizontal normalization, complete and partial horizontal totalizations,
horizontal tails, transgression fibres, and the coherent vertical differential.
In particular the whole critical presymplectic calculus, including vertical
closedness, descends to the reduced critical quotient.
\end{theorem}

\begin{proof}
For fixed $(r,s)$ the map
\[
H_s\times_QV^C_r
\longrightarrow
H_s\times_QV^R_r
\]
is an effective epimorphism, and the preceding theorem identifies its \v{C}ech nerve
with
\[
p\longmapsto H_s\times_QM_{\mathcal K,V^C}^{p,r}.
\]
Effective descent for parameterized spectra identifies the stable coefficient on
the target with the totalization of its pullbacks along this nerve.  Applying the
structural dependent product, which preserves limits, gives
\[
C_{\mathrm{red}}^{r,s}(A;M)
\simeq
\operatorname{Tot}_p
C_{\mathcal K,V^C,H}^{p,r,s}(A;M).
\]
Naturality in $r$ and $s$ assembles these equivalences bicosimplicially.  Every
subsequent construction named in the statement is a finite stable limit, a
normalization fibre, or a partial or complete totalization.  Fubini for limits and
exactness therefore transport the equivalence through the entire calculus.
\end{proof}

\begin{construction}[Critical ordered presymplectic contraction]
\label{constr:symmetry-noether-critical-ordered-presymplectic-contraction}
Project the critical rectangle $\operatorname{Rect}_{V^C}(1,1)$ to the staircase
$(0,0)\to(1,0)\to(1,1)$ and then pull the ambient presymplectic transgression
through
\[
R_{\mathcal K,V^C,H}^{\bullet,\bullet,\bullet}
\longrightarrow
R_{\mathcal G,V,H}^{\bullet,\bullet,\bullet}.
\]
Normalization and horizontal transgression give a canonical critical mixed section
\[
\boxed{
\iota_{\mathcal K}\omega_{L,\mathrm{crit}}^{(q)}
:
\mathbf1
\longrightarrow
\operatorname{Trans}_{\mathcal K,V^C,H}^{1,1,q}(A;M)[n].
}
\]
It is the restriction of the ambient ordered contraction through the constructed
critical formal-groupoid morphism; it is not obtained from an unproved identity
between $R_{C_L/Q,\bullet}$ and a pullback of
$R_{\overline A/Q,\bullet}$.
\end{construction}

\begin{theorem}[Represented critical contraction recognition]
\label{thm:symmetry-noether-represented-critical-contraction-recognition}
On a represented smooth scalar critical chart $C_{L,T}\to T$, let $\xi$ be a
represented first-order direction of the critical formal groupoid and let
\[
\eta\in T_{C_{L,T}/T}
\]
be a represented critical vertical direction.  The exact two-side first-order
extraction of the critical mixed contraction satisfies
\[
\boxed{
\lambda_{\mathcal K,V^C}^{(1,1)}
\left(\iota_{\mathcal K}\omega_{L,\mathrm{crit}}^{(q)}\right)(\xi,\eta)
\simeq
\lambda_{V^C}^{(1,1)}
\left(\omega_{L,\mathrm{crit}}^{(q)}\right)
\left(\rho_{\mathcal K,T}^{\mathrm{rep}}(\xi),\eta\right).
}
\]
No unlinearized tangent-pair evaluation is asserted.
\end{theorem}

\begin{proof}
Apply the ordered side-evaluation construction to the critical formal groupoid.
The first critical symmetry row maps under the critical anchor to the field
variation classified by $\rho_{\mathcal K,T}^{\mathrm{rep}}(\xi)$, while the
terminal row is the unit at the second critical field variation.  The map of
critical rectangles to ambient rectangles preserves this staircase.  The equality
therefore holds before reduced two-variable stabilization and remains an equality
after taking the Boolean cross-effect and exact $(1,1)$-linearization.
\end{proof}


\section{The Hamiltonian defect}
\label{sec:symmetry-noether-the-hamiltonian-defect}

The mixed calculus makes both terms of the Hamiltonian comparison live in
one coefficient.  Through the axis equivalence, regard the universal Noether
current as
\[
J_{\mathcal G,L}^{(q)}
:
\mathbf1
\longrightarrow
F_H^q\operatorname{Sec}_{\mathcal G,V}^{1,0}(A;M)[n].
\]
Apply the mixed vertical differential:
\[
d_V^{\mathrm{mix}}J_{\mathcal G,L}^{(q)}
:
\mathbf1
\longrightarrow
F_H^q\operatorname{Sec}_{\mathcal G,V}^{1,1}(A;M)[n].
\]
Let
\[
\jmath_{\mathrm{tr}}:
\operatorname{Trans}_{\mathcal G,V,H}^{1,1,q}
\longrightarrow
F_H^q\operatorname{Sec}_{\mathcal G,V}^{1,1}
\]
be the fibre inclusion.

\begin{definition}[Hamiltonian defect]
\label{def:symmetry-noether-hamiltonian-defect}
Define the canonical mixed section
\[
\boxed{
\mathfrak H_{\mathcal G,L}^{(q)}
:=
d_V^{\mathrm{mix}}J_{\mathcal G,L}^{(q)}
-
\jmath_{\mathrm{tr}}
\bigl(\iota_{\mathcal G}\omega_L^{(q)}\bigr).
}
\]
It is a section of
$F_H^q\operatorname{Sec}_{\mathcal G,V}^{1,1}(A;M)[n]$.
\end{definition}

\begin{proposition}[Horizontal image of the Hamiltonian defect]
\label{prop:symmetry-noether-horizontal-image-of-the-hamiltonian-defect}
There is a canonical coherent equivalence
\[
\boxed{
\iota_H^q
\mathfrak H_{\mathcal G,L}^{(q)}
\simeq
-
d_V^{\mathrm{mix}}\mathscr E_{\mathcal G,L}^{[q]}.
}
\]
\end{proposition}

\begin{proof}
The mixed vertical differential is a morphism of the complete horizontal
towers.  Applying it to the First Noether homotopy gives
\[
\iota_H^q d_V^{\mathrm{mix}}J_{\mathcal G,L}^{(q)}
\simeq
-d_V^{\mathrm{mix}}\mathscr E_{\mathcal G,L}^{[q]}.
\]
By definition,
$\iota_{\mathcal G}\omega_L^{(q)}$ lies in the transgression fibre, so the
composite of its fibre inclusion with
$F_H^q\to\operatorname{Sec}$ is canonically nullhomotopic.  Subtracting this
zero image gives the formula.
\end{proof}

\begin{remark}[The defect is not declared zero]
\label{rem:symmetry-noether-the-defect-is-not-declared-zero}
Neither the First Noether theorem nor the presymplectic recurrence forces
$\mathfrak H_{\mathcal G,L}^{(q)}$ to vanish.  The defect records precisely
the additional compatibility required for a Hamiltonian realization.
\end{remark}

\section{Hamiltonian realization, residual transgression, and presymplectic radical}
\label{sec:symmetry-noether-hamiltonian-realization-residual-transgression-and-presymplectic-radical}

\begin{definition}[Hamiltonian realization object]
\label{def:symmetry-noether-hamiltonian-realization-object}
Over the common Noether--Lepage parameter object define the geometric zero locus
\[
\boxed{
\operatorname{HamReal}_{\mathcal G,L}^{q}
:=
Z\!\left(\mathfrak H_{\mathcal G,L}^{(q)}\right).
}
\]
A point is exactly a specified coherent homotopy
\[
d_V^{\mathrm{mix}}J_{\mathcal G,L}^{(q)}
\simeq
\jmath_{\mathrm{tr}}
\bigl(
\iota_{\mathcal G}\omega_L^{(q)}
\bigr),
\]
and the object is formed in the geometric parameter category by the parameterized
stable zero-locus construction.
\end{definition}

\begin{definition}[Current-differential nullification]
\label{def:symmetry-noether-current-differential-nullification}
Define
\[
\boxed{
\operatorname{dJNull}_{\mathcal G,L}^{q}
:=
Z\!\left(d_V^{\mathrm{mix}}J_{\mathcal G,L}^{(q)}\right).
}
\]
A point is a specified coherent nullhomotopy of the mixed vertical differential of
the universal current.
\end{definition}

Put
\[
P_{\mathcal G,L}^{q,\mathrm{pre}}
:=
\operatorname{HamReal}_{\mathcal G,L}^{q}
\times
\operatorname{dJNull}_{\mathcal G,L}^{q}
\]
over their common parameter object.

The two nullhomotopies on this fibre product do not yet nullify the complete
presymplectic transgression itself.  They nullify only its image under the fibre
inclusion.  The residual ambiguity is a canonical loop-class, which we now retain
as an obstruction.

Abbreviate, after pullback to
$P_{\mathcal G,L}^{q,\mathrm{pre}}$,
\[
F_q
:=
F_H^q
\operatorname{Sec}_{\mathcal G,V}^{1,1}(A;M)[n],
\qquad
S_q
:=
\operatorname{Sec}_{\mathcal G,V}^{1,1}(A;M)[n],
\]
and
\[
T_q
:=
\operatorname{Trans}_{\mathcal G,V,H}^{1,1,q}(A;M)[n].
\]
By definition there is a fibre sequence
\[
T_q
\xrightarrow{\jmath_{\mathrm{tr}}}
F_q
\longrightarrow
S_q
\]
and hence the rotated fibre sequence
\[
\Omega F_q
\longrightarrow
\Omega S_q
\longrightarrow
T_q
\xrightarrow{\jmath_{\mathrm{tr}}}
F_q.
\]

\begin{construction}[Residual presymplectic transgression]
\label{constr:symmetry-noether-residual-presymplectic-transgression}
Let
\[
w_{\mathcal G,L}^{(q)}
:=
\iota_{\mathcal G}\omega_L^{(q)}
:
\mathbf1\longrightarrow T_q.
\]
Over $P_{\mathcal G,L}^{q,\mathrm{pre}}$, Hamiltonianity identifies
$\jmath_{\mathrm{tr}}w_{\mathcal G,L}^{(q)}$ with
$d_V^{\mathrm{mix}}J_{\mathcal G,L}^{(q)}$, while the
$\operatorname{dJNull}$ factor supplies a nullhomotopy of the latter.
Their composite is therefore a specified nullhomotopy
\[
\eta_{\mathcal G,L}^{(q)}:
\jmath_{\mathrm{tr}}w_{\mathcal G,L}^{(q)}\simeq0.
\]
Equivalently, it is a section
\[
\eta_{\mathcal G,L}^{(q)}:
P_{\mathcal G,L}^{q,\mathrm{pre}}
\longrightarrow
Z\!\left(\jmath_{\mathrm{tr}}w_{\mathcal G,L}^{(q)}\right).
\]
By the universal property of
\[
\Omega S_q\simeq\operatorname{fib}
\left(
T_q\to F_q
\right),
\]
the pair consisting of $w_{\mathcal G,L}^{(q)}$ and this nullhomotopy determines
a canonical residual section
\[
\boxed{
\mathfrak r_{\mathcal G,L}^{(q)}:
\mathbf1
\longrightarrow
\Omega S_q.
}
\]
Its image under $\Omega S_q\to T_q$ is
$w_{\mathcal G,L}^{(q)}$.
\end{construction}

\begin{definition}[Presymplectic-radical realization]
\label{def:symmetry-noether-presymplectic-radical-realization}
Define
\[
\boxed{
\operatorname{RadReal}_{\mathcal G,L}^{q}
}
\]
over $P_{\mathcal G,L}^{q,\mathrm{pre}}$ to be the homotopy fibre over
$\mathfrak r_{\mathcal G,L}^{(q)}$ of
\[
\Omega^\infty\Omega F_q
\longrightarrow
\Omega^\infty\Omega S_q.
\]
Thus a point is a lift
\[
\widetilde{\mathfrak r}:
\mathbf1\longrightarrow\Omega F_q
\]
whose image is the residual transgression.
\end{definition}

\begin{theorem}[Radical realization is the compatible missing nullification]
\label{thm:symmetry-noether-radical-realization-is-the-compatible-missing-nullification}
Over $P_{\mathcal G,L}^{q,\mathrm{pre}}$ there is a canonical equivalence
\[
\boxed{
\operatorname{RadReal}_{\mathcal G,L}^{q}
\simeq
Z_{P_{\mathcal G,L}^{q,\mathrm{pre}}}
\left(
 w_{\mathcal G,L}^{(q)};
 \jmath_{\mathrm{tr}},
 \eta_{\mathcal G,L}^{(q)}
\right).
}
\]
Thus $\operatorname{RadReal}$ is not merely the unconstrained zero locus of the
presymplectic transgression: it is the space of its nullhomotopies whose image in
$F_q$ is the already selected nullhomotopy determined by Hamiltonianity and
$dJ$-nullification.  In particular, over
$\operatorname{RadReal}_{\mathcal G,L}^{q}$ there is a canonical specified
homotopy
\[
\boxed{
\iota_{\mathcal G}\omega_L^{(q)}\simeq0
}
\]
in the complete transgression coefficient, with the required compatibility with
the previously selected image nullhomotopy.
\end{theorem}

\begin{proof}
For the fibre sequence
\[
\Omega F_q\longrightarrow\Omega S_q\longrightarrow T_q
\xrightarrow{\jmath_{\mathrm{tr}}}F_q,
\]
the chosen Hamiltonian and $dJ$ homotopies give the section
$\eta_{\mathcal G,L}^{(q)}$ of the zero locus of
$\jmath_{\mathrm{tr}}w_{\mathcal G,L}^{(q)}$.  The same data identify
$w_{\mathcal G,L}^{(q)}$ with the image of
$\mathfrak r_{\mathcal G,L}^{(q)}$ under $\Omega S_q\to T_q$.

A lift
\[
\widetilde{\mathfrak r}:\mathbf1\to\Omega F_q
\]
of $\mathfrak r_{\mathcal G,L}^{(q)}$ is, by exactness of the rotated fibre
sequence, exactly a nullhomotopy of $w_{\mathcal G,L}^{(q)}$ whose image under
$\jmath_{\mathrm{tr}}$ is the prescribed nullhomotopy
$\eta_{\mathcal G,L}^{(q)}$.  The compatible-nullification construction
$Z_P(\sigma;u,\eta)$ represents precisely this lifting problem.  Fibre-pasting
therefore gives the displayed equivalence.
\end{proof}

\begin{remark}[Hamiltonianity alone is weaker]
\label{rem:symmetry-noether-hamiltonianity-alone-is-weaker}
A point of
$\operatorname{HamReal}_{\mathcal G,L}^{q}
\times\operatorname{dJNull}_{\mathcal G,L}^{q}$
canonically nullifies the image of the presymplectic transgression in $F_q$.
It does not in general nullify the point of $T_q$.  The residual class
$\mathfrak r_{\mathcal G,L}^{(q)}$ measures exactly this difference.
\end{remark}

\subsection{Critical presymplectic radical}
\label{subsec:symmetry-noether-critical-presymplectic-radical}

Pull the critical mixed calculus along
\[
\operatorname{RadReal}_{\mathcal G,L}^{q}\longrightarrow Q.
\]
The map from the critical mixed grid to the ambient mixed grid carries the universal
compatible nullhomotopy to a nullhomotopy
\[
\iota_{\mathcal K}\omega_{L,\mathrm{crit}}^{(q)}\simeq0
\]
in the complete critical mixed transgression coefficient.  To extract an ordinary
represented radical from this complete nullhomotopy, we first construct the exact
first-order contraction target rather than identifying the raw degree-two
coefficient with an exterior-form complex.

Work on a represented smooth scalar degree-zero critical chart
\[
C_{L,T}\longrightarrow T
\]
and let
\[
c:\operatorname{Spec}B\longrightarrow C_{L,T}
\]
be an affine critical point.  Put
\[
N_c:=\Gamma\!\left(c^*L_{C_{L,T}/T}\right),
\qquad
T_c:=N_c^\vee.
\]
The chartwise quasi-smoothness theorem makes $N_c$ perfect.

\begin{lemma}[Stable extension of the exact critical side coefficient]
\label{lem:symmetry-noether-stable-extension-of-the-exact-critical-side-coefficient}
For connective $B$-modules $P,Q$, the maximal-critical ordered side square and its
Boolean total fibre are defined without a perfectness hypothesis.  Successive
reduced stabilization in the two independent direction variables produces a
bifunctor on the stabilized direction categories which is exact separately in each
variable.  Restricting this stabilized bifunctor to compact objects gives a
canonical separately exact $B$-linear bifunctor
\[
\operatorname{Perf}(B)\times\operatorname{Perf}(B)
\longrightarrow
\operatorname{Mod}_B.
\]
\end{lemma}

\begin{proof}
Structured split square-zero extensions are available for every connective module,
so the four side coefficients and their Boolean total fibre are already defined on
the full connective direction category.  Reduced stabilization is the universal
exact extension of a reduced connective square-zero direction to its stabilization.
Applying it successively in the two independent variables therefore produces a
bifunctor on the two stabilized direction categories which is exact separately in
each variable.  At a represented affine centre, the compact objects of the
stabilized $B$-module direction category are precisely $\operatorname{Perf}(B)$.
Restricting to compact objects gives the displayed bifunctor.  Perfectness enters
only at this restriction and later when ordinary tensor duality is used; it is not
an input to the side-square construction itself.  All coefficient maps before
stabilization are $B$-linear, and stabilization preserves that linearity.
\end{proof}

\begin{construction}[Exact bilinear coefficient of the critical presymplectic transgression]
\label{constr:symmetry-noether-exact-bilinear-critical-presymplectic-coefficient}
For perfect $B$-modules $P,Q$, define
\[
\boxed{
\mathscr B_{\omega,c}^{q}(P,Q)
}
\]
to be the separately exact stable bifunctor supplied by the preceding lemma.  For
connective perfect $P,Q$, it is concretely the Boolean total fibre of the maximal
critical side-coefficient square followed by reduced stabilization in both
variables; a pair of actual critical first-order directions determines a section
of this coefficient by
Lemma~\ref{lem:symmetry-noether-side-coefficient-and-evaluated-section-separation}.
Define
\[
D_{\omega,c}^{q}
:=
\mathscr B_{\omega,c}^{q}(B,B).
\]
Then there are canonical natural equivalences
\[
\boxed{
\mathscr B_{\omega,c}^{q}(P,Q)
\simeq
P\otimes_BQ\otimes_BD_{\omega,c}^{q}
}
\]
for perfect $P,Q$.
\end{construction}

\begin{proof}
For fixed $Q$, both functors
\[
P\longmapsto
P\otimes_B\mathscr B_{\omega,c}^{q}(B,Q),
\qquad
P\longmapsto
\mathscr B_{\omega,c}^{q}(P,Q)
\]
are exact $B$-linear functors on $\operatorname{Perf}(B)$, and the canonical
comparison is an equivalence at the tensor unit $B$.  The full subcategory on which
it is an equivalence is therefore thick.  Since $B$ thickly generates
$\operatorname{Perf}(B)$, the comparison is an equivalence for every perfect $P$:
\[
\mathscr B_{\omega,c}^{q}(P,Q)
\simeq
P\otimes_B\mathscr B_{\omega,c}^{q}(B,Q).
\]
Apply the same argument to the exact $B$-linear functor
$Q\mapsto\mathscr B_{\omega,c}^{q}(B,Q)$ to obtain
\[
\mathscr B_{\omega,c}^{q}(B,Q)
\simeq
Q\otimes_B\mathscr B_{\omega,c}^{q}(B,B)
=
Q\otimes_BD_{\omega,c}^{q}.
\]
Composing the two canonical equivalences gives the displayed tensor formula with
naturality in both variables.
\end{proof}

\begin{construction}[Affine-point presymplectic contraction morphism]
\label{constr:symmetry-noether-affine-point-presymplectic-contraction-morphism}
Take the two universal critical first-order directions
\[
\operatorname{id}_{N_c}:N_c\longrightarrow N_c.
\]
The exact $(1,1)$-linearization of the critical presymplectic transgression gives a
section
\[
\widetilde\omega_{L,c}^{(q)}:
B
\longrightarrow
N_c\otimes_BN_c\otimes_BD_{\omega,c}^{q}.
\]
Since $N_c$ is perfect, duality in the first factor adjoints this section to
\[
\boxed{
\omega_{L,c}^{(q),\flat}:
T_c
\longrightarrow
N_c\otimes_BD_{\omega,c}^{q}.
}
\]
Put
\[
\mathcal W_{\omega,c}^{(q)}
:=
N_c\otimes_BD_{\omega,c}^{q}.
\]
\end{construction}

\begin{proposition}[Affine-point descent of the presymplectic contraction target]
\label{prop:symmetry-noether-affine-point-descent-presymplectic-contraction-target}
The modules $D_{\omega,c}^{q}$ and $\mathcal W_{\omega,c}^{(q)}$, together with the
maps $\omega_{L,c}^{(q),\flat}$, are Cartesian under represented affine restriction.
Consequently they assemble to
\[
\boxed{
\mathcal W_{\omega,T}^{(q)}\in QCoh(C_{L,T})
}
\]
and a quasi-coherent morphism
\[
\boxed{
\omega_{L,T}^{(q),\flat}:
T_{C_{L,T}/T}
\longrightarrow
\mathcal W_{\omega,T}^{(q)}.
}
\]
This morphism is the exact represented first-order contraction of the complete
critical presymplectic transgression.  No identification
$\mathcal W_{\omega,T}^{(q)}\simeq L_{C_{L,T}/T}[m]$ is used or asserted.
\end{proposition}

\begin{proof}
Affine restriction carries the critical cotangent coefficient $N_c$ to its base
change and transports the critical maximal relation, the critical formal groupoid,
the ordered side-evaluation diagrams, Boolean cross-effects, and both reduced
stabilizations by the represented Beck--Chevalley maps.  Hence
$D_{\omega,c}^{q}$ base changes to the corresponding coefficient on the restricted
point.  Tensoring with $N_c$ and dualizing the perfect first factor commute with the
same base change.  The affine-point right-Kan description of quasi-coherent modules
therefore glues the Cartesian family to the displayed object and morphism.
\end{proof}

\begin{definition}[Represented presymplectic radical]
\label{def:symmetry-noether-represented-presymplectic-radical}
Define
\[
\boxed{
\operatorname{Rad}(\omega_{L,T}^{(q)})
:=
\operatorname{fib}
\left(
\omega_{L,T}^{(q),\flat}:
T_{C_{L,T}/T}
\longrightarrow
\mathcal W_{\omega,T}^{(q)}
\right).
}
\]
This is the radical of the constructed exact first-order contraction, not a fibre
of an unconstructed exterior-form identification.
\end{definition}

\begin{theorem}[Universal critical symmetry-radical factorization]
\label{thm:symmetry-noether-universal-critical-symmetry-radical-factorization}
After pullback to
$\operatorname{RadReal}_{\mathcal G,L}^{q}$, the represented tangent anchor of the
critical formal groupoid factors canonically as
\[
\boxed{
a_{\mathcal K,T}
\longrightarrow
\operatorname{Rad}(\omega_{L,T}^{(q)})
\longrightarrow
T_{C_{L,T}/T}.
}
\]
\end{theorem}

\begin{proof}
The universal point of $\operatorname{RadReal}_{\mathcal G,L}^{q}$ gives an
actual nullhomotopy of the complete ambient mixed presymplectic transgression.
Restriction through the critical mixed-grid morphism gives a nullhomotopy of the
complete critical ordered contraction.  The Boolean cross-effect and reduced
stabilization used to construct the exact $(1,1)$-linearization are functorial
finite stable operations, so they carry this nullhomotopy to a nullhomotopy of its
exact first-order image.  By represented critical contraction recognition, that
image on a critical symmetry direction $\xi$ is
\[
\omega_{L,T}^{(q),\flat}
\left(\rho_{\mathcal K,T}^{\mathrm{rep}}(\xi)\right).
\]
Thus the composite
\[
a_{\mathcal K,T}
\xrightarrow{\rho_{\mathcal K,T}^{\mathrm{rep}}}
T_{C_{L,T}/T}
\xrightarrow{\omega_{L,T}^{(q),\flat}}
\mathcal W_{\omega,T}^{(q)}
\]
carries a canonical nullhomotopy.  The universal property of the homotopy fibre
produces the displayed factorization.
\end{proof}

\begin{definition}[Presymplectic-radical discrepancy]
\label{def:symmetry-noether-presymplectic-radical-discrepancy}
On the represented stable coefficient domain define
\[
\boxed{
C_{\mathcal G,L,T}^{\mathrm{rad},q}
:=
\operatorname{cofib}
\left(
a_{\mathcal K,T}
\longrightarrow
\operatorname{Rad}(\omega_{L,T}^{(q)})
\right).
}
\]
Its vanishing is equivalent to the symmetry-to-radical morphism being an
equivalence.  Thus exhaustion of represented first-order presymplectic null
directions by critical symmetry directions is measured by a constructed
discrepancy rather than inserted as reduction data.
\end{definition}


\section{Complete presymplectic descent, basicness, and parameterized reduction}
\label{sec:symmetry-noether-derived-presymplectic-basicness-and-parameterized-reduction}

Put
\[
P_\omega^q:=\operatorname{Lep}^q(L),
\qquad
C_\omega^q:=C_L\times_QP_\omega^q.
\]
Pull the universal presymplectic transgression to $C_\omega^q$ and denote its
stable coefficient and section by
\[
(E_\omega^{(q)},\omega_{\mathrm{crit}}^{(q)}).
\]
Use the parameterized critical formal groupoid
\[
\mathcal K_{\omega,\bullet}
:=
\mathcal K_{\mathcal G,L,P_\omega^q,\bullet}
\rightrightarrows
C_\omega^q
\]
and write $J_{\mathcal K,\omega}$ for its stable jet comonad.


Put
\[
R_\omega^q
:=
\left(\operatorname{Crit}(L)//\mathcal G\right)\times_QP_\omega^q,
\qquad
q_\omega:C_\omega^q\twoheadrightarrow R_\omega^q.
\]
Write
\[
V^{C_\omega}_r:=R_{C_\omega^q/P_\omega^q,r},
\qquad
V^{R_\omega}_r:=R_{R_\omega^q/P_\omega^q,r}.
\]
By the critical mixed-grid descent theorem and arbitrary parameterized base change,
the symmetry direction of the critical mixed grid is the \v{C}ech direction of the
maps from the maximal infinitesimal relation of $C_\omega^q/P_\omega^q$ to the
maximal infinitesimal relation of $R_\omega^q/P_\omega^q$.

\begin{construction}[Cosimplicial coherently closed presymplectic transgression object]
\label{constr:symmetry-noether-cosimplicial-closed-presymplectic-transgression-object}
Leave both the symmetry degree and the vertical cosimplicial degree unnormalized.
Write
\[
C_{\mathcal K,\omega,V,H}^{p,r,s}(A;M)
\]
for the base change to $P_\omega^q$ of the parameterized critical mixed coefficient
in symmetry degree $p$, critical-vertical degree $r$, and horizontal degree $s$.
For each $p\ge0$ form the vertical cosimplicial stable objects
\[
\mathscr S_{\mathcal K,\omega}^{p,\bullet}
:=
\operatorname{Tot}_H
C_{\mathcal K,\omega,V,H}^{p,\bullet,\bullet}(A;M),
\]
\[
\mathscr S_{\mathcal K,\omega}^{p,\bullet,<q}
:=
\operatorname{Tot}_{H,\le q-1}
C_{\mathcal K,\omega,V,H}^{p,\bullet,\bullet}(A;M),
\]
and the horizontal tail
\[
F_H^q\mathscr S_{\mathcal K,\omega}^{p,\bullet}
:=
\operatorname{fib}
\left(
\mathscr S_{\mathcal K,\omega}^{p,\bullet}
\longrightarrow
\mathscr S_{\mathcal K,\omega}^{p,\bullet,<q}
\right).
\]
Define the vertical cosimplicial stage-$q$ transgression object by
\[
\boxed{
\mathscr T_{\mathcal K,\omega}^{p,\bullet,q}
:=
\operatorname{fib}
\left(
F_H^q\mathscr S_{\mathcal K,\omega}^{p,\bullet}
\longrightarrow
\mathscr S_{\mathcal K,\omega}^{p,\bullet}
\right).
}
\]
All these constructions are limits in stable functor categories and are
cosimplicial in the unnormalized symmetry degree $p$.

Vertical normalization is a finite matching fibre and therefore commutes with the
horizontal totalizations, horizontal tail, and transgression fibre.  Consequently,
for every $r\ge0$, there is a canonical equivalence
\[
\boxed{
N_V^r
\mathscr T_{\mathcal K,\omega}^{p,\bullet,q}
\simeq
\widetilde{\operatorname{Trans}}_{\mathcal K,\omega}^{p;r,q}.
}
\]
In particular the previously used degree-two and degree-three transgression
coefficients are the normalized layers $r=2,3$ of this one vertical cosimplicial
object, and the unary morphism
\[
d_V^{\mathrm{tr}}:
\widetilde{\operatorname{Trans}}_{\mathcal K,\omega}^{p;2,q}
\longrightarrow
\widetilde{\operatorname{Trans}}_{\mathcal K,\omega}^{p;3,q}
\]
is only the first boundary shadow of its complete vertical coherent structure.

Form the complete vertical totalization filtration
\[
F_V^r\operatorname{Tot}_V
\mathscr T_{\mathcal K,\omega}^{p,\bullet,q}
:=
\operatorname{fib}
\left(
\operatorname{Tot}_V\mathscr T_{\mathcal K,\omega}^{p,\bullet,q}
\longrightarrow
\operatorname{Tot}_{V,\le r-1}
\mathscr T_{\mathcal K,\omega}^{p,\bullet,q}
\right).
\]
Define the stable coefficient of coherently closed shifted degree-two
transgressions by
\[
\boxed{
\widetilde{\operatorname{Trans}}_{\mathcal K,\omega}^{p;2,q,\mathrm{cl}}(A;M;n)
:=
F_V^2\operatorname{Tot}_V
\mathscr T_{\mathcal K,\omega}^{p,\bullet,q}[n+2].
}
\]
Stable Dold--Kan identifies
\[
\operatorname{gr}_V^2
\left(
F_V^\bullet\operatorname{Tot}_V
\mathscr T_{\mathcal K,\omega}^{p,\bullet,q}
\right)
\simeq
N_V^2\mathscr T_{\mathcal K,\omega}^{p,\bullet,q}[-2],
\]
so the associated-graded map gives the underlying-transgression morphism
\[
\boxed{
u_{V,2}^{p,q}:
\widetilde{\operatorname{Trans}}_{\mathcal K,\omega}^{p;2,q,\mathrm{cl}}(A;M;n)
\longrightarrow
\widetilde{\operatorname{Trans}}_{\mathcal K,\omega}^{p;2,q}[n].
}
\]
Define
\[
\boxed{
\operatorname{ClTrans}_{\mathcal K,\omega}^{p,q}
:=
\Omega_{P_\omega^q}^\infty
\widetilde{\operatorname{Trans}}_{\mathcal K,\omega}^{p;2,q,\mathrm{cl}}(A;M;n).
}
\]
A point is therefore a degree-two transgression together with a specified
nullhomotopy of its first vertical boundary and the full hierarchy of higher
compatibility data encoded by the complete vertical totalization filtration.  In
particular $\operatorname{ClTrans}$ is not defined as the zero locus of the unary
map $d_V^{\mathrm{tr}}$.

The universal critical presymplectic transgression already carries complete
coherent vertical closedness in the variational calculus.  Restriction through the
critical mixed-grid morphism is a morphism of the complete vertical coherent
systems.  It therefore gives a canonical section
\[
\boxed{
\omega_{\mathrm{crit}}^{(q),\mathrm{cl}}:
P_\omega^q
\longrightarrow
\operatorname{ClTrans}_{\mathcal K,\omega}^{0,q}
}
\]
whose image under $u_{V,2}^{0,q}$ is the underlying critical presymplectic
transgression $\omega_{\mathrm{crit}}^{(q)}$.
\end{construction}

\begin{definition}[Complete presymplectic descent realization]
\label{def:symmetry-noether-complete-presymplectic-descent-realization}
Restriction from the complete symmetry totalization to symmetry degree zero gives
\[
\operatorname{Tot}_p
\operatorname{ClTrans}_{\mathcal K,\omega}^{p,q}
\longrightarrow
\operatorname{ClTrans}_{\mathcal K,\omega}^{0,q}.
\]
Define
\[
\boxed{
\operatorname{PreSymDesc}_{\omega,\mathcal G}^{q}(L)
:=
P_\omega^q
\times_{\operatorname{ClTrans}_{\mathcal K,\omega}^{0,q}}
\operatorname{Tot}_p
\operatorname{ClTrans}_{\mathcal K,\omega}^{p,q}.
}
\]
A point is the complete symmetry-\v{C}ech descent datum of the coherently closed
critical presymplectic transgression, including its entire vertical closedness
hierarchy.  No Cartan coaction, invariance datum, or basicness nullhomotopy is added
independently.
\end{definition}

\begin{theorem}[Complete presymplectic descent to the reduced critical calculus]
\label{thm:symmetry-noether-complete-presymplectic-descent-to-reduced-critical-calculus}
Over
\[
B_\omega^q:=\operatorname{PreSymDesc}_{\omega,\mathcal G}^{q}(L)
\]
let
\[
\mathscr T_{\mathrm{red},\omega}^{\bullet,q}
\]
denote the vertical cosimplicial stage-$q$ horizontal transgression object formed
from the reduced coefficient
$C_{\mathrm{red}}^{\bullet,\bullet}(A;M)$ after base change to $P_\omega^q$, by the
same construction used above before vertical normalization.  Define its coherently
closed shifted degree-two coefficient by
\[
\boxed{
\operatorname{Trans}_{V^{R_\omega},H}^{2,q,\mathrm{cl}}(A;M;n)
:=
F_V^2\operatorname{Tot}_V
\mathscr T_{\mathrm{red},\omega}^{\bullet,q}[n+2].
}
\]
Its associated-graded map gives
\[
\boxed{
u_{V,2}^{\mathrm{red}}:
\operatorname{Trans}_{V^{R_\omega},H}^{2,q,\mathrm{cl}}(A;M;n)
\longrightarrow
\operatorname{Trans}_{V^{R_\omega},H}^{2,q}(A;M)[n].
}
\]

The mixed-grid \v{C}ech theorem gives a canonical equivalence of vertical
cosimplicial transgression objects
\[
\boxed{
\operatorname{Tot}_p
\mathscr T_{\mathcal K,\omega}^{p,\bullet,q}
\simeq
\mathscr T_{\mathrm{red},\omega}^{\bullet,q}.
}
\]
Consequently Fubini for limits gives a canonical equivalence of complete closed
coefficients
\[
\boxed{
\operatorname{Tot}_p
\widetilde{\operatorname{Trans}}_{\mathcal K,\omega}^{p;2,q,\mathrm{cl}}(A;M;n)
\simeq
\operatorname{Trans}_{V^{R_\omega},H}^{2,q,\mathrm{cl}}(A;M;n).
}
\]
The universal point of $B_\omega^q$ therefore constructs a canonically descended
coherently closed section
\[
\boxed{
\omega_{L,\mathrm{red}}^{(q),\mathrm{cl}}:
\mathbf1
\longrightarrow
\operatorname{Trans}_{V^{R_\omega},H}^{2,q,\mathrm{cl}}(A;M;n).
}
\]
Define its underlying reduced presymplectic transgression by
\[
\boxed{
\omega_{L,\mathrm{red}}^{(q)}
:=
u_{V,2}^{\mathrm{red}}
\omega_{L,\mathrm{red}}^{(q),\mathrm{cl}}:
\mathbf1
\longrightarrow
\operatorname{Trans}_{V^{R_\omega},H}^{2,q}(A;M)[n].
}
\]
This section lives on
\[
R_{\omega,\mathrm{red}}^q
:=
\left(\operatorname{Crit}(L)//\mathcal G\right)\times_QB_\omega^q,
\qquad
\pi_{\mathrm{red}}:R_{\omega,\mathrm{red}}^q\to B_\omega^q,
\]
with effective quotient map
\[
q_{\omega,B}:
C_L\times_QB_\omega^q
\twoheadrightarrow
R_{\omega,\mathrm{red}}^q.
\]
Its pullback to $C_L\times_QB_\omega^q$ is the original critical presymplectic
transgression equipped with the complete selected symmetry descent datum and the
full inherited coherent vertical closedness structure.
\end{theorem}

\begin{proof}
Apply
Theorem~\ref{thm:symmetry-noether-descent-of-complete-critical-vertical-horizontal-calculus}
after base change from $Q$ to $P_\omega^q$.  Before vertical normalization, that
theorem identifies the symmetry totalization of the critical mixed coefficient with
the reduced vertical--horizontal coefficient.  Horizontal totalization, horizontal
partial totalization, the horizontal-tail fibre, and the transgression fibre are
all limits in stable functor categories.  They therefore commute with symmetry
totalization and give
\[
\operatorname{Tot}_p
\mathscr T_{\mathcal K,\omega}^{p,\bullet,q}
\simeq
\mathscr T_{\mathrm{red},\omega}^{\bullet,q}
\]
as vertical cosimplicial stable objects.

Now apply the complete vertical closedness construction itself, rather than only
its unary boundary.  Since
\[
F_V^2\operatorname{Tot}_V(-)
=
\operatorname{fib}
\left(
\operatorname{Tot}_V(-)
\longrightarrow
\operatorname{Tot}_{V,\le1}(-)
\right)
\]
is again a limit construction, Fubini yields
\[
\begin{aligned}
\operatorname{Tot}_p
F_V^2\operatorname{Tot}_V
\mathscr T_{\mathcal K,\omega}^{p,\bullet,q}[n+2]
&\simeq
F_V^2\operatorname{Tot}_V
\operatorname{Tot}_p
\mathscr T_{\mathcal K,\omega}^{p,\bullet,q}[n+2]\\
&\simeq
F_V^2\operatorname{Tot}_V
\mathscr T_{\mathrm{red},\omega}^{\bullet,q}[n+2].
\end{aligned}
\]
This is the displayed equivalence of coherently closed degree-two transgression
coefficients.  Relative $\Omega^\infty$ preserves limits, so the same equivalence
holds for the corresponding parameterized closed-transgression objects.

The defining pullback of $\operatorname{PreSymDesc}$ selects the point whose
symmetry degree-zero restriction is
$\omega_{\mathrm{crit}}^{(q),\mathrm{cl}}$.  Under the preceding equivalence this
point is precisely the reduced closed section
$\omega_{L,\mathrm{red}}^{(q),\mathrm{cl}}$.  Applying its
underlying-transgression map gives $\omega_{L,\mathrm{red}}^{(q)}$.  Because the
source is the complete vertical totalization filtration, the descended datum
contains the specified first vertical-boundary nullhomotopy and every higher
closedness compatibility, not merely the unary equation
$d_V^{\mathrm{tr}}\omega\simeq0$.  Effective descent identifies its pullback with
the original critical closed transgression carrying the selected complete
symmetry-\v{C}ech descent datum.
\end{proof}

\begin{construction}[Internal Cartan lift of the presymplectic pair]
\label{constr:symmetry-noether-internal-cartan-lift-of-the-presymplectic-pair}
For a test $T\to P_\omega^q$, let
\[
\mathsf{Pair}_{\mathrm{Cart}}(T)
\]
be the maximal $\infty$-groupoid of pairs consisting of a
$J_{\mathcal K,\omega,T}$-coalgebra coefficient on
$C_\omega^q\times_{P_\omega^q}T$ and a coalgebra section of that coefficient.
Let
\[
\mathsf{Pair}_{\mathrm{und}}(T)
\]
be the corresponding maximal $\infty$-groupoid of underlying coefficient--section
pairs.  These assignments satisfy descent.  Define
\[
\boxed{
\operatorname{CartLift}_{\omega,\mathcal G}^{q}(L)
\longrightarrow
P_\omega^q
}
\]
to be the internal homotopy fibre of
\[
\mathsf{Pair}_{\mathrm{Cart}}
\longrightarrow
\mathsf{Pair}_{\mathrm{und}}
\]
over the parameterized pair
$(E_\omega^{(q)},\omega_{\mathrm{crit}}^{(q)})$.
\end{construction}

A $T$-point is therefore an actual Cartan coaction on the pulled-back
presymplectic coefficient together with a lift of the pulled-back
presymplectic section as a Cartan section.

\begin{proposition}[Complete descent supplies the Cartan presymplectic pair]
\label{prop:symmetry-noether-complete-descent-supplies-cartan-presymplectic-pair}
Let
\[
B_\omega^q
=
\operatorname{PreSymDesc}_{\omega,\mathcal G}^{q}(L).
\]
After pullback to $B_\omega^q$, the reduced coefficient and reduced section
constructed by
Theorem~\ref{thm:symmetry-noether-complete-presymplectic-descent-to-reduced-critical-calculus}
determine canonically a Cartan coalgebra structure on the pulled-back
presymplectic coefficient $E_\omega^{(q)}$ together with a Cartan lift of
$\omega_{\mathrm{crit}}^{(q)}$.  Consequently there is a canonical morphism
\[
\boxed{
\operatorname{PreSymDesc}_{\omega,\mathcal G}^{q}(L)
\longrightarrow
\operatorname{CartLift}_{\omega,\mathcal G}^{q}(L).
}
\]
\end{proposition}

\begin{proof}
Write $E_{\omega,\mathrm{red}}^{(q)}$ for the reduced stable coefficient carrying
$\omega_{L,\mathrm{red}}^{(q)}$ on
\[
R_{\omega,\mathrm{red}}^q
=
\left(\operatorname{Crit}(L)//\mathcal G\right)\times_QB_\omega^q.
\]
The complete presymplectic descent theorem constructs a canonical identification
\[
q_{\omega,B}^*E_{\omega,\mathrm{red}}^{(q)}
\simeq
E_\omega^{(q)}|_{B_\omega^q}
\]
whose associated complete symmetry-\v{C}ech descent datum is the universal point of
$B_\omega^q$, and under this identification the pullback of
$\omega_{L,\mathrm{red}}^{(q)}$ is
$\omega_{\mathrm{crit}}^{(q)}$ with its selected complete descent coherence.

Effective descent for the base-changed critical formal groupoid is an equivalence
between stable coefficients on $R_{\omega,\mathrm{red}}^q$ and coalgebras for the
base-changed jet comonad on $C_L\times_QB_\omega^q$.  It therefore transports
$E_{\omega,\mathrm{red}}^{(q)}$ to a canonical Cartan coalgebra structure on the
identified pullback coefficient.  Full faithfulness on mapping spaces transports
the reduced section to a morphism of Cartan coalgebras from the tensor unit, whose
underlying section is exactly $\omega_{\mathrm{crit}}^{(q)}$.  This coefficient--section
pair is a point of the defining homotopy fibre of
$\operatorname{CartLift}_{\omega,\mathcal G}^{q}(L)$, naturally in the parameter,
which gives the displayed morphism.
\end{proof}

\begin{proposition}[Cartan equivariance of mixed contraction]
\label{prop:symmetry-noether-cartan-equivariance-of-mixed-contraction}
After pullback to
$\operatorname{CartLift}_{\omega,\mathcal G}^{q}(L)$, the critical mixed
contraction
\[
\iota_{\mathcal K}\omega_{\mathrm{crit}}^{(q)}
\]
is canonically a morphism in the stable Cartan coefficient category of the
corresponding base-changed critical formal groupoid.
\end{proposition}

\begin{proof}
Every stage of the critical mixed grid is formed by limits from the critical
formal-groupoid relation, the maximal critical vertical relation
$R_{C_L/Q,\bullet}$, and the horizontal Cartan relation.  Its comparison with the
ambient mixed grid is induced by the formal-groupoid morphism
$\mathcal K_{\mathcal G,L,\bullet}\to\mathcal G_\bullet$ and automatic anchor
compatibility.  The ordered staircase is natural for this map.  Stable pullback,
normalization, and horizontal transgression are functorial, so the selected Cartan
transport commutes with the critical mixed contraction in every simplicial degree
and with its higher cocycle coherence.
\end{proof}

\begin{definition}[Internal presymplectic basicness object]
\label{def:symmetry-noether-internal-presymplectic-basicness-object}
Over $\operatorname{CartLift}_{\omega,\mathcal G}^{q}(L)$, the preceding
proposition makes
\[
\iota_{\mathcal K}\omega_{\mathrm{crit}}^{(q)}
\]
a section in the stable Cartan coefficient system.  Apply the parameterized stable
zero-locus construction in that coefficient system and define
\[
\boxed{
\operatorname{BasicReal}_{\omega,\mathcal G}^{q}(L)
:=
Z^{\mathrm{Cart}}_{\operatorname{CartLift}_{\omega,\mathcal G}^{q}(L)}
\left(
\iota_{\mathcal K}\omega_{\mathrm{crit}}^{(q)}
\right)
\longrightarrow
P_\omega^q.
}
\]
A point is a Cartan lift of the presymplectic coefficient--section pair together
with a specified Cartan nullhomotopy of the symmetry contraction.  The nullification
is therefore a geometric parameter object, not a stable fibre mistaken for one.
\end{definition}


\begin{lemma}[Quotient degeneracy of the critical ordered staircase]
\label{lem:symmetry-noether-quotient-degeneracy-of-critical-ordered-staircase}
Let
\[
\tau_{\mathcal K,V^C}:
M_{\mathcal K,V^C}^{1,1}
\longrightarrow
V^C_2
\]
denote the critical ordered staircase map used in
Construction~\ref{constr:symmetry-noether-critical-ordered-presymplectic-contraction},
and let
\[
q_r:V^C_r\longrightarrow V^R_r
\]
be the vertexwise quotient maps of
Theorem~\ref{thm:symmetry-noether-critical-mixed-grid-is-cech-resolution}.  There is
a canonical morphism
\[
\beta:
M_{\mathcal K,V^C}^{1,1}
\longrightarrow
V^R_1
\]
and a canonical coherent equivalence
\[
\boxed{
q_2\tau_{\mathcal K,V^C}
\simeq
s_0\beta.
}
\]
The same identity holds after arbitrary parameterized base change, in particular
over $B_\omega^q$.  Consequently the pullback along the critical ordered staircase
of every normalized reduced vertical degree-two cochain carries the canonical
nullhomotopy induced by the degeneracy $s_0$, coherently in the symmetry-\v{C}ech
direction.
\end{lemma}

\begin{proof}
A point of $M_{\mathcal K,V^C}^{1,1}$ is a critical rectangle with vertices
\[
x_{00},\ x_{10},\ x_{01},\ x_{11}
\]
whose first row is lifted to an arrow of
$\mathcal K_{\mathcal G,L,1}$.  The ordered staircase has vertices
\[
x_{00},\ x_{10},\ x_{11}.
\]
Because
\[
q_L:C_L\twoheadrightarrow R_L
\]
is the effective quotient of $\mathcal K_{\mathcal G,L,\bullet}$, the endpoints of
every critical symmetry arrow have canonically the same image in $R_L$.  Hence
\[
q_L(x_{00})\simeq q_L(x_{10}).
\]
The common infinitesimal image of the rectangle gives a reduced relation edge from
$q_L(x_{00})$ to $q_L(x_{11})$; define $\beta$ by this edge.  The image of the
ordered staircase in $V^R_2$ is therefore
\[
\bigl(q_L(x_{00}),q_L(x_{10}),q_L(x_{11})\bigr)
\simeq
\bigl(q_L(x_{00}),q_L(x_{00}),q_L(x_{11})\bigr),
\]
which is exactly the degeneracy $s_0\beta$.  The quotient, staircase, and vertex
maps are simplicial and functorial under parameterized pullback, so the equivalence
carries all higher coherence and survives every such base change.

A normalized vertical degree-two cochain lies in the fibre of its codegeneracy
maps and therefore carries the canonical nullhomotopy after restriction along
$s_0$.  Pullback through the displayed factorization gives the asserted
nullhomotopy on the critical ordered staircase.  Naturality of the factorization in
the symmetry-\v{C}ech direction supplies all higher compatibility.
\end{proof}

\begin{theorem}[Complete presymplectic descent implies Cartan basicness]
\label{thm:symmetry-noether-complete-presymplectic-descent-implies-cartan-basicness}
There is a canonical morphism
\[
\boxed{
\operatorname{PreSymDesc}_{\omega,\mathcal G}^{q}(L)
\longrightarrow
\operatorname{BasicReal}_{\omega,\mathcal G}^{q}(L).
}
\]
In particular, complete presymplectic descent constructs both the Cartan transport
of the presymplectic coefficient--section pair and the coherent nullhomotopy
\[
\boxed{
\iota_{\mathcal K}\omega_{\mathrm{crit}}^{(q)}\simeq0.
}
\]
Thus presymplectic basicness is a consequence of complete descent rather than an
additional reduction hypothesis.
\end{theorem}

\begin{proof}
Proposition~\ref{prop:symmetry-noether-complete-descent-supplies-cartan-presymplectic-pair}
constructs the canonical morphism
\[
\operatorname{PreSymDesc}_{\omega,\mathcal G}^{q}(L)
\longrightarrow
\operatorname{CartLift}_{\omega,\mathcal G}^{q}(L),
\]
so the coefficient and the critical presymplectic section already carry their
Cartan structures on the complete-descent parameter object.

Pull the reduced vertically normalized degree-two transgression to the same base.
By
Lemma~\ref{lem:symmetry-noether-quotient-degeneracy-of-critical-ordered-staircase},
the quotient of the critical ordered staircase factors through the degeneracy
$s_0$.  Vertical normalization therefore supplies a canonical nullhomotopy of its
restriction.  Under the complete mixed-grid descent equivalence this restriction is
exactly
\[
\iota_{\mathcal K}\omega_{\mathrm{crit}}^{(q)}.
\]
The quotient-degeneracy factorization is simplicially natural in the
symmetry-\v{C}ech direction, so the nullhomotopy commutes with the selected Cartan
transport and carries all higher cocycle coherence.  It is therefore a
nullhomotopy in the stable Cartan coefficient category.  The parameterized Cartan
zero-locus universal property now lifts the preceding morphism through
$\operatorname{BasicReal}_{\omega,\mathcal G}^{q}(L)$, giving the displayed
canonical map.
\end{proof}

\begin{theorem}[Universal parameterized presymplectic reduction]
\label{thm:symmetry-noether-universal-parameterized-presymplectic-reduction}
Let
\[
B_\omega^q
:=
\operatorname{PreSymDesc}_{\omega,\mathcal G}^{q}(L).
\]
The structural map
\[
B_\omega^q\longrightarrow P_\omega^q\longrightarrow Q
\]
makes the base change
\[
\left(
\operatorname{Crit}(L)//\mathcal G
\right)
\times_QB_\omega^q
\]
well typed.  The reduced section
\[
\boxed{
\omega_{L,\mathrm{red}}^{(q)}
}
\]
and its coherent vertical closedness are the objects constructed by the complete
presymplectic descent theorem.  Its pullback along
\[
C_L\times_QB_\omega^q
\twoheadrightarrow
\left(
\operatorname{Crit}(L)//\mathcal G
\right)\times_QB_\omega^q
\]
is the selected critical presymplectic section with its complete symmetry-\v{C}ech
descent datum; the Cartan transport and null-contraction homotopy are the
consequences constructed in
Theorem~\ref{thm:symmetry-noether-complete-presymplectic-descent-implies-cartan-basicness}.
\end{theorem}

\begin{proof}
This is the parameter-base reformulation of
Theorem~\ref{thm:symmetry-noether-complete-presymplectic-descent-to-reduced-critical-calculus}.
The descent equivalence constructs the reduced coefficient, section, vertical
differential, and closedness simultaneously.  The basicness theorem records the
Cartan and degeneracy shadows of this complete descent.
\end{proof}


The reduction geometry is summarized by the Cartesian/effective-quotient diagram
\begin{tikzcd}[column sep=huge,row sep=large]
C_L\times_QB_\omega^q
  \arrow[r,two heads,"q_{\mathcal G,L}\times_QB_\omega^q"]
  \arrow[d] &
R_{\omega,\mathrm{red}}^q
  \arrow[d,"\pi_{\mathrm{red}}"] \\
B_\omega^q
  \arrow[r,equal] &
B_\omega^q,
\end{tikzcd}
where the closed presymplectic calculus on the upper left is the pullback of the
constructed reduced calculus on the upper right.

\begin{construction}[Intrinsic formal-leaf reduced tangent coefficient]
\label{constr:symmetry-noether-intrinsic-formal-leaf-reduced-tangent-coefficient}
On the locally coherent formalized coefficient domain, base-change the critical
formal groupoid to $B_\omega^q$ and let
\[
L_{\mathcal K,\omega}^{\mathrm{fm},sp}
\]
denote its constructed global formal leaf.  Define the formal-leaf ambient tangent
coefficient on $C_L\times_QB_\omega^q$ by
\[
\boxed{
T_{\omega,\mathrm{red}}^{pc}
:=
\operatorname{cofib}
\left(
T^{pc}_{(C_L\times_QB_\omega^q)/L_{\mathcal K,\omega}^{\mathrm{fm},sp}}
\longrightarrow
T^{pc}_{(C_L\times_QB_\omega^q)/B_\omega^q}
\right).
}
\]
The spectral anchor-fibre transitivity theorem canonically identifies this object
with the underlying homotopy anchor fibre of the differentiated base-changed
critical formal groupoid.  If the reduced formal leaf is represented by the
parameterized quotient $R_{\omega,\mathrm{red}}^q\to B_\omega^q$, ordinary tangent
transitivity identifies
\[
T_{\omega,\mathrm{red}}^{pc}
\simeq
q_{\omega,B}^*T^{pc}_{R_{\omega,\mathrm{red}}^q/B_\omega^q}.
\]
Thus a reduced tangent coefficient exists intrinsically on the formalized branch;
the affine-point ordinary tangent construction below is its represented
comparison, not a chosen quotient presentation.
\end{construction}


\begin{theorem}[Presymplectic descent forces critical symmetry directions into the represented radical]
\label{thm:symmetry-noether-presymplectic-descent-forces-critical-symmetry-radical}
On a represented smooth scalar critical chart, after pullback to
$\operatorname{PreSymDesc}_{\omega,\mathcal G}^{q}(L)$ there is a canonical
factorization
\[
\boxed{
a_{\mathcal K,T}
\longrightarrow
\operatorname{Rad}(\omega_{L,T}^{(q)})
\longrightarrow
T_{C_{L,T}/T}.
}
\]
This factorization is induced by complete quotient descent and is independent of the
Hamiltonian residual-nullification construction defining
$\operatorname{RadReal}_{\mathcal G,L}^{q}$.
\end{theorem}

\begin{proof}
The complete presymplectic descent theorem gives the canonical symmetry-contraction
nullhomotopy by quotient degeneracy.  The Boolean total fibre and two-variable
reduced stabilization defining the exact $(1,1)$ contraction are functorial finite
stable operations, so they carry this complete nullhomotopy to a nullhomotopy of
\[
a_{\mathcal K,T}
\xrightarrow{\rho_{\mathcal K,T}^{\mathrm{rep}}}
T_{C_{L,T}/T}
\xrightarrow{\omega_{L,T}^{(q),\flat}}
\mathcal W_{\omega,T}^{(q)}.
\]
The homotopy-fibre universal property gives the factorization.
\end{proof}

\begin{construction}[Hamiltonian--descent compatibility locus]
\label{constr:symmetry-noether-hamiltonian-descent-compatibility-locus}
Form the common parameter base
\[
P_{\mathrm{HamBasic}}^q
:=
\operatorname{PreSymDesc}_{\omega,\mathcal G}^{q}(L)
\times_{P_\omega^q}
\operatorname{RadReal}_{\mathcal G,L}^{q},
\]
using the structural map from $\operatorname{RadReal}$ through the common
Noether--Lepage parameter to $P_\omega^q$.  After pullback to this base, complete
descent and $\operatorname{RadReal}$ supply two specified nullhomotopies of the
same critical mixed contraction.  Let
\[
Z_\omega:=
Z\!\left(\iota_{\mathcal K}\omega_{\mathrm{crit}}^{(q)}\right)
\]
be its parameterized zero-locus object.  The two nullhomotopies are sections
$s_{\mathrm{desc}},s_{\mathrm{Ham}}:P_{\mathrm{HamBasic}}^q\rightrightarrows Z_\omega$.
Define
\[
\boxed{
\operatorname{HamBasicReal}_{\mathcal G,L}^{q}
:=
P_{\mathrm{HamBasic}}^q
\times_{Z_\omega\times Z_\omega}
Z_\omega^{\Delta^1},
}
\]
where $Z_\omega^{\Delta^1}\to Z_\omega\times Z_\omega$ is endpoint evaluation and
the first map is $(s_{\mathrm{desc}},s_{\mathrm{Ham}})$.  A point therefore
constructs the coherent identification between the quotient-descent nullhomotopy
and the Hamiltonian residual nullhomotopy instead of assuming that they coincide.
\end{construction}

\begin{convention}[Affine-point theorem domain for reduced first-order geometry]
\label{conv:symmetry-noether-affine-point-theorem-domain-for-reduced-first-order-geometry}
Retain the already constructed parameterized quotient and projection
\[
\pi_{\mathrm{red}}:
R_{\omega,\mathrm{red}}^{q}
\longrightarrow B_\omega^q.
\]
On the represented reduced scalar theorem domain, assume that for every affine
point
\[
t:T=\operatorname{Spec}C\longrightarrow B_\omega^q
\]
the base change
\[
R_T:=
T\times_{B_\omega^q}R_{\omega,\mathrm{red}}^{q}
\longrightarrow T
\]
is represented by spectral schemes and has perfect relative cotangent complex.
This is a theorem-domain property of the already constructed parameterized
quotient.  No global represented presentation of
$R_{\omega,\mathrm{red}}^{q}$ is chosen or retained.
\end{convention}

\begin{construction}[Affine-point reduced cotangent and tangent coefficients]
\label{constr:symmetry-noether-affine-point-reduced-cotangent-and-tangent-coefficients}
Let
\[
c:U=\operatorname{Spec}C\longrightarrow
R_{\omega,\mathrm{red}}^{q}
\]
be an affine point and put
\[
t_c:=\pi_{\mathrm{red}}c:
U\longrightarrow B_\omega^q,
\qquad
R_c:=
U\times_{B_\omega^q}R_{\omega,\mathrm{red}}^{q}.
\]
The point $c$ and the identity of $U$ determine a section
\[
\bar c:U\longrightarrow R_c.
\]
Define
\[
\boxed{
\left(
L_{R_{\omega,\mathrm{red}}^{q}/B_\omega^q}^{\mathrm{lin}}
\right)_c
:=
\bar c^*L_{R_c/U}.
}
\]
For a morphism of affine points
\[
g:(U',c')\longrightarrow(U,c),
\]
the canonical identification
\[
R_{c'}\simeq U'\times_UR_c
\]
and cotangent base change give
\[
\mathcal O_{U'}\otimes_{\mathcal O_U}
\bar c^*L_{R_c/U}
\simeq
\bar c'{}^*L_{R_{c'}/U'}.
\]
Hence these modules form a Cartesian affine-point family and therefore define
\[
\boxed{
L_{R_{\omega,\mathrm{red}}^{q}/B_\omega^q}^{\mathrm{lin}}
\in
QCoh(R_{\omega,\mathrm{red}}^{q}).
}
\]
They are pointwise perfect by the theorem-domain convention.  Define
\[
\boxed{
T_{R_{\omega,\mathrm{red}}^{q}/B_\omega^q}^{\mathrm{lin}}
:=
\left(
L_{R_{\omega,\mathrm{red}}^{q}/B_\omega^q}^{\mathrm{lin}}
\right)^\vee.
}
\]
For every affine test $T\to B_\omega^q$ in the stated theorem domain, affine-point
recognition identifies the pullbacks with the ordinary represented coefficients
\[
\left.
L_{R_{\omega,\mathrm{red}}^{q}/B_\omega^q}^{\mathrm{lin}}
\right|_{R_T}
\simeq
L_{R_T/T},
\qquad
\left.
T_{R_{\omega,\mathrm{red}}^{q}/B_\omega^q}^{\mathrm{lin}}
\right|_{R_T}
\simeq
T_{R_T/T}.
\]
\end{construction}

\begin{construction}[Affine-point reduced presymplectic contraction]
\label{constr:symmetry-noether-affine-point-reduced-presymplectic-contraction}
Apply the exact two-side first-order construction of
Construction~\ref{constr:symmetry-noether-exact-bilinear-critical-presymplectic-coefficient}
to the descended section $\omega_{L,\mathrm{red}}^{(q)}$ after every represented
affine base change $R_T\to T$.  At an affine point
$c:U=\operatorname{Spec}C\to R_{\omega,\mathrm{red}}^{q}$ put
\[
N_c:=
\left(
L_{R_{\omega,\mathrm{red}}^{q}/B_\omega^q}^{\mathrm{lin}}
\right)_c,
\qquad
T_c:=N_c^\vee.
\]
For connective $C$-modules the reduced presymplectic side-coefficient square and
its Boolean total fibre are defined first.  Successive reduced stabilization in the
two independent variables gives a separately exact functor on the stabilized
direction categories.  For perfect $C$-modules $P,Q$, let
\[
\mathscr B_{\omega,\mathrm{red},c}^{q}(P,Q)
\]
be the restriction of this separately exact $C$-linear bifunctor to compact
objects.  Put
\[
D_{\omega,\mathrm{red},c}^{q}
:=
\mathscr B_{\omega,\mathrm{red},c}^{q}(C,C).
\]
The same thick-generation argument as in the critical calculation gives
canonical natural equivalences
\[
\mathscr B_{\omega,\mathrm{red},c}^{q}(P,Q)
\simeq
P\otimes_CQ\otimes_CD_{\omega,\mathrm{red},c}^{q}.
\]

Evaluate the reduced presymplectic section on the two universal first variations
classified by
\[
\operatorname{id}_{N_c}:N_c\longrightarrow N_c.
\]
Its exact $(1,1)$ image is a section
\[
C
\longrightarrow
N_c\otimes_CN_c\otimes_CD_{\omega,\mathrm{red},c}^{q}.
\]
Perfect duality in the first factor adjoints this section to
\[
\boxed{
\omega_{L,\mathrm{red},c}^{(q),\flat}:
T_c
\longrightarrow
N_c\otimes_CD_{\omega,\mathrm{red},c}^{q}.
}
\]
Put
\[
\mathcal W_{\omega,\mathrm{red},c}^{(q)}
:=
N_c\otimes_CD_{\omega,\mathrm{red},c}^{q}.
\]
Represented cotangent base change, functoriality of the side-coefficient square,
Boolean total fibre, reduced stabilization, tensor product, and perfect duality
make these modules and morphisms Cartesian under affine restriction.  Affine-point
descent therefore constructs
\[
\boxed{
\mathcal W_{\omega,\mathrm{red}}^{(q)}
\in
QCoh(R_{\omega,\mathrm{red}}^{q})
}
\]
and the canonical morphism
\[
\boxed{
\omega_{L,\mathrm{red}}^{(q),\flat}:
T_{R_{\omega,\mathrm{red}}^{q}/B_\omega^q}^{\mathrm{lin}}
\longrightarrow
\mathcal W_{\omega,\mathrm{red}}^{(q)}.
}
\]
On every represented affine test $T\to B_\omega^q$, this is the exact two-side
first-order contraction
\[
T_{R_T/T}
\longrightarrow
\mathcal W_{\omega,\mathrm{red},T}^{(q)}
\]
of the pulled-back reduced presymplectic section.
\end{construction}

\begin{definition}[Reduced exact first-order nondegeneracy locus]
\label{def:symmetry-noether-reduced-nondegeneracy-locus}
For every test $T\to B_\omega^q$, use the contravariant descent coefficient system
\[
\mathcal C_{\mathrm{red}}(T)
:=
QCoh\!\left(
R_{\omega,\mathrm{red}}^{q}
\times_{B_\omega^q}T
\right).
\]
The pullback of $\omega_{L,\mathrm{red}}^{(q),\flat}$ is a morphism in
$\mathcal C_{\mathrm{red}}(T)$.  Define the derived first-order nondegeneracy locus
by applying the parameterized walking-equivalence construction with this coefficient
system:
\[
\boxed{
\operatorname{SympReal}_{\omega,\mathcal G}^{q}(L)
:=
\operatorname{Inv}_{B_\omega^q}
\left(
\omega_{L,\mathrm{red}}^{(q),\flat}
\right).
}
\]
Thus a $T$-point exists precisely when the complete base-changed reduced
contraction
\[
T_{R_T/T}
\longrightarrow
\mathcal W_{\omega,\mathrm{red},T}^{(q)}
\]
over
$R_T=
R_{\omega,\mathrm{red}}^{q}\times_{B_\omega^q}T$
is an equivalence, and when it exists its space is contractible.  Reduced
nondegeneracy is therefore represented by invertibility of the constructed exact
first-order contraction itself.  No identification of the raw presymplectic
transgression, or of its linearized target, with an exterior power or shifted
cotangent complex is inserted into the definition.
\end{definition}

\section{Reduced current classes and Cartan-diagram realization}
\label{sec:symmetry-noether-reduced-current-classes-and-cartan-morphism-realization}

Put
\[
P_J^q:=\mathcal N_{\mathcal G}^{q}(L),
\qquad
C_J^q:=C_L\times_QP_J^q.
\]
Pull the universal Noether-current coefficient and section to this family:
\[
(E_J^{(q)},J_{\mathcal G,L,\mathrm{crit}}^{(q)}).
\]
Use the parameterized critical groupoid
\[
\mathcal K_{J,\bullet}
:=
\mathcal K_{\mathcal G,L,P_J^q,\bullet}
\rightrightarrows C_J^q
\]
with jet comonad $J_{\mathcal K,J}$.

\begin{construction}[Internal current Cartan-lift object]
\label{constr:symmetry-noether-internal-current-cartan-lift-object}
Define
\[
\boxed{
\operatorname{CartLift}_{J,\mathcal G}^{q}(L)
\longrightarrow
P_J^q
}
\]
to be the internal homotopy fibre, over
$(E_J^{(q)},J_{\mathcal G,L,\mathrm{crit}}^{(q)})$, of the parameterized
forgetful morphism from
$J_{\mathcal K,J}$-coalgebra coefficient--section pairs to underlying stable
coefficient--section pairs.
\end{construction}

\begin{theorem}[Universal reduced current]
\label{thm:symmetry-noether-universal-reduced-current}
Let
\[
B_J^q:=\operatorname{CartLift}_{J,\mathcal G}^{q}(L).
\]
Its structural map
\[
B_J^q\to P_J^q\to Q
\]
makes the quotient
\[
\left(
\operatorname{Crit}(L)//\mathcal G
\right)\times_QB_J^q
\]
well typed.  Effective descent gives a canonical reduced current coefficient and
section
\[
\boxed{
J_{\mathcal G,L,\mathrm{red}}^{(q)}
}
\]
on this quotient.
\end{theorem}

\begin{proof}
Base-change the parameterized critical formal groupoid to $B_J^q$ and apply the
effective descent equivalence to the universal selected Cartan
coefficient--section pair.
\end{proof}

The native complete-horizontal transgression
$\mathcal U_{\mathcal G,L}^{(q)}$ is defined over the full
Euler-nullification object.  Compatibility of this particular morphism with a
selected Cartan structure is a second internal realization problem.

\begin{construction}[Common parameter object for current transgression]
\label{constr:symmetry-noether-common-parameter-object-for-current-transgression}
The current Cartan-lift object lies over $P_J^q$, and
$\operatorname{EulNull}_{\mathcal G}^{q}(L)$ carries its canonical map to the same
Noether--Lepage parameter object.  Define
\[
\boxed{
P_U^q
:=
\operatorname{EulNull}_{\mathcal G}^{q}(L)
\times_{P_J^q}
B_J^q.
}
\]
This fibre product is formed between internal parameter objects.  The
Euler-nullification factor gives a canonical structural map
\[
P_U^q\longrightarrow C_L.
\]
Base-change the current parameterized critical groupoid to $B_J^q$ and restrict its
object space along the induced section
\[
P_U^q\longrightarrow C_L\times_QB_J^q.
\]
Write
\[
\mathcal H_{U,\bullet}\rightrightarrows P_U^q
\]
for this raw object-space restriction and $J_{\mathcal H,U}$ for its jet comonad.
Apply the formalization coreflection and denote the resulting formal groupoid by
\[
\mathcal K_{U,\bullet}
:=
\widehat{\mathcal H_{U,\bullet}}^{\,\mathrm{dR}/Q}
\rightrightarrows P_U^q
\]
with jet comonad $J_{\mathcal K,U}$.  The formalization counit gives a canonical
fixed-object groupoid morphism
\[
\chi_U:
\mathcal K_{U,\bullet}
\longrightarrow
\mathcal H_{U,\bullet}.
\]
\end{construction}

\begin{construction}[Internal Cartan-diagram realization of the current transgression]
\label{def:symmetry-noether-internal-cartan-morphism-realization-object}
After pullback to $P_U^q$, abbreviate
\[
F_q^J
:=
F_H^q\operatorname{Sec}_{\mathcal G}^{1}(A;M)[n],
\qquad
S_q^J
:=
\operatorname{Sec}_{\mathcal G}^{1}(A;M)[n],
\]
and let
\[
u_q:F_q^J\longrightarrow S_q^J
\]
be the horizontal-tail inclusion.  The pulled-back current is a section
\[
J:\mathbf1\longrightarrow F_q^J,
\]
and the Euler-nullification factor of $P_U^q$ supplies the specified coherent
homotopy
\[
\eta_U:u_qJ\simeq0.
\]
Equivalently these data form an underlying $2$-simplex
\[
D_U:\Delta^2\longrightarrow E_{P_U^q}
\]
whose edges $0\to1$, $1\to2$, and $0\to2$ are respectively $J$, $u_q$, and the
zero map, and whose $2$-cell is $\eta_U$.

For a test $T\to P_U^q$, put
\[
\mathsf{Diag}_{\mathrm{Cart}}(T)
:=
\operatorname{Fun}
\left(
\Delta^2,
\operatorname{Coalg}_{J_{\mathcal K,U,T}}(E_T)
\right)^\simeq
\]
and
\[
\mathsf{Diag}_{\mathrm{und}}(T)
:=
\operatorname{Fun}(\Delta^2,E_T)^\simeq.
\]
Let $\mathsf{Pair}_{\mathrm{Cart}}(T)$ and
$\mathsf{Pair}_{\mathrm{und}}(T)$ be the analogous functor groupoids on the edge
$\Delta^{\{0,1\}}\simeq\Delta^1$.  Restriction to that edge and forgetting
coalgebra structures give a natural morphism
\[
\mathsf{Diag}_{\mathrm{Cart}}
\longrightarrow
\mathsf{Diag}_{\mathrm{und}}
\times_{\mathsf{Pair}_{\mathrm{und}}}
\mathsf{Pair}_{\mathrm{Cart}}.
\]
The factor $B_J^q$ of $P_U^q$ supplies a universal selected Cartan lift
$b_J$ of the current coefficient--section pair for the base-changed current
groupoid.  Raw object-space restriction along
$P_U^q\to C_L\times_QB_J^q$ and the effective-image disk--jet
Beck--Chevalley comparison first restrict $b_J$ to a
$J_{\mathcal H,U}$-coalgebra coefficient--section pair on the edge $0\to1$; write
this restricted pair as $b_J|_{P_U^q}$.  The fixed-object groupoid morphism
$\chi_U$ induces the contravariant morphism of jet comonads
\[
\chi_U^J:
J_{\mathcal H,U}
\longrightarrow
J_{\mathcal K,U}.
\]
Hence there is a canonical restriction functor on coalgebras
\[
\chi_U^{\sharp}:
\operatorname{Coalg}_{J_{\mathcal H,U}}(E_{P_U^q})
\longrightarrow
\operatorname{Coalg}_{J_{\mathcal K,U}}(E_{P_U^q}),
\qquad
(E,\rho)\longmapsto
\bigl(E,\chi_{U,E}^J\circ\rho\bigr).
\]
Because $\chi_U^J$ is a morphism of comonads, the counit and coassociativity
identities, together with all higher coherence, are preserved by
$\chi_U^{\sharp}$.  Define the canonically induced $\mathcal K_U$-Cartan lift by
\[
\boxed{
 b_{J,U}
 :=
 \chi_U^{\sharp}\bigl(b_J|_{P_U^q}\bigr).
}
\]
No new Cartan datum is introduced by this restriction.

Define
\[
\boxed{
\operatorname{CartMorReal}_{\mathcal G,L}^{q}
\longrightarrow
P_U^q
}
\]
to be the internal homotopy fibre of the preceding morphism over the pair
$(D_U,b_{J,U})$.  Thus a point extends the already selected, canonically restricted
Cartan structure on the current to a Cartan structure on the complete triangle
\[
\mathbf1\xrightarrow{J}F_q^J\xrightarrow{u_q}S_q^J,
\qquad
u_qJ\simeq0,
\]
including a Cartan lift of $u_q$, a Cartan lift of the chosen nullhomotopy, and all
higher compatibility data.  No Cartan structure on $S_q^J$ or on the
transgression fibre is named before this realization is formed.
\end{construction}

\begin{proposition}[The native transgression becomes a Cartan morphism]
\label{prop:symmetry-noether-native-transgression-becomes-cartan-morphism}
Over $\operatorname{CartMorReal}_{\mathcal G,L}^{q}$, the fibre
\[
\operatorname{Trans}_{\mathcal G,H}^{1,q}(A;M)[n]
=
\operatorname{fib}(F_q^J\to S_q^J)
\]
has the induced Cartan coalgebra structure, and the native Noether transgression
\[
\mathcal U_{\mathcal G,L}^{(q)}:
\mathbf1\longrightarrow
\operatorname{Trans}_{\mathcal G,H}^{1,q}(A;M)[n]
\]
is canonically a morphism in the stable Cartan coefficient category.
\end{proposition}

\begin{proof}
The stable jet comonad is exact, so its coalgebra category is stable and its
forgetful functor creates finite limits.  The universal Cartan $2$-simplex gives a
Cartan morphism $F_q^J\to S_q^J$ and a Cartan current whose composite carries the
selected Cartan nullhomotopy to zero.  Taking the fibre in the coalgebra category
therefore gives a Cartan structure on the underlying stable transgression fibre.
Its universal fibre map from the current is exactly
$\mathcal U_{\mathcal G,L}^{(q)}$ because forgetting coalgebra structures recovers
the underlying triangle $D_U$.
\end{proof}

\begin{theorem}[Descent of the constructed current transgression]
\label{thm:symmetry-noether-descent-of-the-constructed-current-transgression}
Over $\operatorname{CartMorReal}_{\mathcal G,L}^{q}$, the complete Cartan triangle,
its transgression fibre, and the native Noether transgression descend uniquely
through the effective quotient of the base-changed critical formal groupoid.  The
descended fibre and morphism are the corresponding complete-horizontal
transgression of the already selected reduced current.
\end{theorem}

\begin{proof}
The universal point supplies an actual finite diagram in the stable Cartan
category.  Effective formal-groupoid descent is an exact equivalence of stable
$\infty$-categories and is fully faithful on mapping spaces.  It therefore
transports the Cartan $2$-simplex, its fibre sequence, the selected current, and the
selected nullhomotopy to the quotient.  Exactness identifies the descended fibre
with the fibre of the descended horizontal-tail morphism, while full faithfulness
identifies the descended section with the unique transgression of the reduced
current carrying the descended nullhomotopy.  All higher coherences are transported
by the same equivalence.
\end{proof}

\section{Functoriality in the symmetry groupoid}
\label{sec:symmetry-noether-functoriality-in-the-symmetry-groupoid}

Let
\[
F:\mathcal G_\bullet\longrightarrow\mathcal H_\bullet
\]
be a fixed-object morphism of formal field-symmetry groupoids over
$\overline A/Q$.

\begin{theorem}[Fixed-field variance]
\label{thm:symmetry-noether-fixed-field-variance}
The morphism $F$ induces coherently:
\begin{enumerate}
\item covariant functors
      \[
      \mathcal V_{\mathcal G}^{(k)}
      \longrightarrow
      \mathcal V_{\mathcal H}^{(k)}
      \]
      on all structured symmetry-variation categories;
\item contravariant restriction morphisms
      \[
      V_{\mathcal H,M}^{(k)}(A)
      \longrightarrow
      V_{\mathcal G,M}^{(k)}(A)
      \]
      on their internal right-Kan coefficients;
\item contravariant restriction on symmetry--horizontal cochains and hence on
      complete symmetry realizations, Noether currents, native
      symmetry-coherence obstruction towers, Euler-nullifications, and native
      current transgressions;
\item on the formalized coherent spectral coefficient domain, the induced maps of
      formal leaves, Artinian complete-symmetry presheaves, their relative unstable
      Schlessinger reflections, stable reflected Noether coefficients, retractive
      differentiated algebroids, differentiated symmetry parameters, First Noether
      derivatives, Taylor classes, syzygies, morphism-latching systems, the
      coefficient-connection hierarchy, the intrinsic pro-coherent anchor-splitting
      curvature/Bianchi realization, and its represented finite-locally-free
      comparison;
\item covariant morphisms of raw critical restriction groupoids, their canonical
      formalizations, parameterized base changes, and effective critical quotients;
\item covariant morphisms of the ambient and critical rectangular mixed grids;
      stable cochains therefore vary contravariantly, carrying the ordered
      presymplectic contractions, Hamiltonian defects, residual transgressions, and
      radical-realization objects;
\item functorial maps of the complete presymplectic-descent, internal Cartan-lift,
      Cartan-diagram, basicness, reduced-current, and reduced-presymplectic
      realization objects.
\end{enumerate}
Every comparison morphism respects the variance inherited from its constituent
construction.
\end{theorem}

\begin{proof}
A $\mathcal G$-lift postcomposes with $F$ to an $\mathcal H$-lift, giving the
covariant functor on variation categories.  Restriction of right-Kan cones has the
opposite variance.  Relation maps pull stable cochains back contravariantly, and
normalization, partial totalization, connecting morphisms, complete limits, and
homotopy fibres preserve those maps.

On the formalized branch, quotient formalization and formal leaves are functorial
for fixed object space.  Artinian restriction, unstable and stable Schlessinger
localization, relative infinite-loop formation, partition-Lie differentiation,
finite limits, and the morphism corolla--latching construction are all functorial,
so they transport the differentiated Noether package.

Critical raw restriction is a degreewise finite-limit pullback of $F$ and
formalization is functorial.  Automatic compatibility of canonical anchors then
induces the maps of critical maximal-relation rectangles.  The Hamiltonian and
reduction objects are obtained from the resulting stable diagrams by zero loci,
rotated fibre sequences, Cartan mapping objects, and effective descent.
\end{proof}
\section{Arbitrary raw spectral base change}
\label{sec:symmetry-noether-arbitrary-raw-spectral-base-change}

Let
\[
S'\longrightarrow S
\]
be a morphism of spectral bases and put primes on raw pullbacks.

\begin{theorem}[Arbitrary base change for internal native Noether geometry]
\label{thm:symmetry-noether-arbitrary-base-change-for-internal-native-noether-geometry}
There are canonical coherent base-change equivalences for:
\begin{enumerate}
\item the symmetry--horizontal relation, stable cochains, complete bifiltration, and
      every internal one- and finite-cubical right-Kan symmetry coefficient;
\item the nondegenerate-chain relation-simplex diagrams and their finite symmetry-lift
      homotopy limits at each pulled-back affine centre;
\item the partial and complete symmetry-realization objects, the native First
      Noether current, global and local one-stage Euler restrictions, and the
      complete native symmetry-coherence obstruction tower;
\item Euler-nullification and complete-horizontal current transgression;
\item the raw critical restriction, its formalization, effective quotient, and every
      parameterized critical formal groupoid obtained by base change;
\item the ambient and critical rectangular mixed relations, mixed coherent
      tricomplexes, ordered presymplectic contractions, Hamiltonian defects,
      Hamiltonian realizations, current-differential nullifications, residual
      transgressions, and $\operatorname{RadReal}$;
\item the complete presymplectic-descent, internal Cartan-lift, Cartan-diagram,
      basicness, reduced-current, and reduced-presymplectic constructions after
      base-changing their parameter objects.
\end{enumerate}

For the external variation categories $\mathcal V_{\mathcal G}^{(k)}$ there are
canonical base-change restriction functors.  No arbitrary base-change equivalence
of those entire external categories is asserted.
\end{theorem}

\begin{proof}
Raw relation nerves, canonical anchors, formalization, quotients, disk--jet
calculus, and effective descent have the arbitrary-base-change comparisons
constructed in the preceding chapters.  The variational chapter gives arbitrary
base change for internal structured square-zero and cubical coefficients and
identifies the comma categories computing right-Kan values after pullback.

The cubical symmetry-lift condition is a finite limit of base-change-compatible
mapping problems.  Thus the external categories restrict functorially, while the
pointwise comma-category equivalences and stable Beck--Chevalley maps identify the
internal right-Kan coefficients.  Parameterized pullback preserves the limits used
in normalization, totalization, mapping fibres, ambient and critical finite \v{C}ech rectangles,
Hamiltonian defects, residual fibre sequences, and internal realization objects.
\end{proof}

\begin{proposition}[Base-change strength of the formalized algebraic branch]
\label{prop:symmetry-noether-base-change-strength-of-the-formalized-algebraic-branch}
For the quotient-formalized Artinian branch, an arbitrary spectral base change
induces the canonical mate comparisons supplied by quotient formalization,
Artinian restriction, unstable and stable Schlessinger reflection, relative
infinite-loop formation, partition-Lie differentiation, and tangent
representations.

On the coherent varying-object spectral-\'etale domain of Chapters 14--15 these
comparisons are equivalences.  For the unstable reflected symmetry parameter this is
Proposition~\ref{prop:symmetry-noether-spectral-etale-exchange-for-relative-unstable-schlessinger-reflector};
the stable and differentiated comparisons are the corresponding exchange and
naturality theorems of Chapters 14--15.  They transport
\[
\widehat{\operatorname{Sym}}_{\mathcal G,L}^{\mathrm{Sch},q},\qquad
\widehat{\mathcal O}_{\mathcal G,L}^{1,q},\qquad
\mathbb O_{\mathcal G,L}^{1,q},\qquad
\mathfrak b_{\mathcal G,L}^{1,q},\qquad
\mathfrak s_{\mathcal G,L}^{q},
\]
the zero and Noether splittings, the differentiated Second Noether homotopy,
$d_0\mathbf n_{\mathcal G,L}^{1,q}$, the field-action factorization
$\partial_{\mathrm{fld}}\mathbf n_{\mathcal G,L}^{1,q}$, the Taylor classes,
syzygy maps, morphism-latching tower, the intrinsic $\operatorname{SpCocyc}$,
tangent representation-bar constructions, the canonical Noether coefficient
connection and its flatness/Bianchi hierarchy, the intrinsic pro-coherent
anchor-splitting realization with curvature/Bianchi hierarchy, and its represented
finite-locally-free comparison.  For a general spectral base change, only the
canonical mate comparisons are asserted on this formalized branch.
\end{proposition}

\begin{proposition}[Base change of the affine-point critical and presymplectic linearizations]
\label{prop:symmetry-noether-base-change-of-the-affine-point-critical-linearization}
On the representably smooth scalar degree-zero sector, the affine-point linear
coefficient
\[
\mathcal E=L_{\overline A/Q}^{\mathrm{lin}},
\]
the native Euler zero locus, the affine-point derivative $de_L^{\mathrm{lin}}$, and
$T_{C_L/Q}^{\mathrm{lin}}$ commute with arbitrary pullback in the representably
smooth theorem domain.  The ordered side-evaluation diagrams, their Boolean
cross-effects, the exact $(1,1)$ reduced stabilization, the affine-point coefficients
$D_{\omega,c}^{q}$ and $\mathcal W_{\omega,c}^{(q)}$, the descended contraction
\[
\omega_{L,T}^{(q),\flat}:
T_{C_{L,T}/T}\longrightarrow\mathcal W_{\omega,T}^{(q)},
\]
and its radical commute with the same represented affine pullbacks.  After a
represented affine base change, the chartwise quasi-smooth zero-locus, ordinary
tangent/Jacobi fibre theorem, and exact first-order presymplectic-radical
factorization are therefore the corresponding pullbacks of the represented chart
calculations.
\end{proposition}
\section{Formal-\'etale locality}
\label{sec:symmetry-noether-formal-etale-locality}

\begin{theorem}[Formal-\'etale locality of the native Noether geometry]
\label{thm:symmetry-noether-formal-etale-locality-of-the-native-noether-geometry}
Let
\[
u:U\longrightarrow X
\]
be formally \'etale over $S$.  Pullback along $Q_{U/S}\to Q_{X/S}$ transports by
canonical equivalences all internal native objects of the symmetry--horizontal
calculus, including the right-Kan cubical coefficients, complete symmetry
realizations, native symmetry-coherence obstruction tower, First Noether current,
Euler-nullification, critical formalization and quotient, ambient and critical
mixed calculi, complete presymplectic descent, Hamiltonian defect,
$\operatorname{RadReal}$, and internal Cartan reduction objects.

The external symmetry-variation categories carry the induced restriction functors.
Their global spaces of choices are not asserted to be equivalent merely from
formal-\'etale restriction.

Whenever the representably smooth scalar critical/Jacobi theorem domain is
preserved, the affine-point linear critical tangent coefficient and the exact
$(1,1)$ presymplectic linearization are transported as well.  On represented affine
charts the ordinary critical tangent, Jacobi zero-mode, critical-Hessian radical,
the contraction target $\mathcal W_{\omega,T}^{(q)}$, and the represented
first-order presymplectic-radical factorization are carried to the corresponding
restricted chart calculations.
\end{theorem}

\begin{proof}
Formal-\'etale locality identifies the horizontal and ambient relative de Rham
relation objects with their pullbacks.  The formal symmetry groupoid and its
canonical anchor pull back accordingly.  Every internal native construction is
obtained from these relation objects by stable limits, structured square-zero
pullback, pointwise right Kan extension, formalization, realization, finite \v{C}ech rectangle, or
internal mapping fibre.

For the critical mixed calculus, first transport the critical object and its
formalization; formal-\'etale locality then identifies $R_{C_L/Q,\bullet}$ with the
corresponding pullback maximal relation.  Hence the critical rectangle and its map
to the ambient rectangle are transported canonically.  Residual transgression and
$\operatorname{RadReal}$ are formed from the resulting stable fibre sequences.
\end{proof}

\begin{proposition}[Spectral-\'etale transport of the formalized branch]
\label{prop:symmetry-noether-formal-etale-transport-of-the-formalized-branch}
On the locally coherent spectral coefficient sector, suppose the varying-object
restriction lies in the represented coherent spectral-\'etale domain of
Chapters 14--15.  Then the coherent spectral-\'etale transition equivalences carry
the Artinian complete-symmetry presheaf and, by
Proposition~\ref{prop:symmetry-noether-spectral-etale-exchange-for-relative-unstable-schlessinger-reflector},
its relative unstable Schlessinger reflection, together with the reflected Noether
coefficient, retractive differentiated algebroid,
differentiated reflected-symmetry parameter, base-changed Noether and zero
splittings, first derivative, its field-action factorization through the
differentiated anchor, directional Taylor classes, syzygies, morphism-latching
obstruction system, Noether--Spencer splitting object, and tangent-representation
chain object to the corresponding constructions on the restricted chart.

A merely formally \'etale varying-object map outside this spectral-\'etale theorem
domain receives only the native locality statements of the preceding theorem; no
formalized partition-Lie equivalence is asserted there.
\end{proposition}
\section{Finite spectral Nisnevich descent}
\label{sec:symmetry-noether-finite-spectral-nisnevich-descent}

\begin{theorem}[Finite spectral Nisnevich descent]
\label{thm:symmetry-noether-finite-spectral-nisnevich-descent}
For a finite spectral Nisnevich cover of a represented affine field centre, the
following internal objects are recovered as totalizations of their \v{C}ech
restrictions:
\begin{enumerate}
\item the symmetry cubical coefficients, symmetry--horizontal stable coefficient
      objects, partial and complete symmetry realizations, and native
      symmetry-coherence obstruction objects;
\item every finite stage of the ambient and critical mixed
      symmetry--vertical--horizontal coefficient systems, the Hamiltonian defect,
      residual transgression, and radical-realization object; on the represented
      scalar sector, the exact $(1,1)$ side-linearization diagrams and their
      affine-point contraction targets also descend;
\item the raw critical restriction, its formalization, and the corresponding
      parameterized quotient diagrams;
\item on the coherent formalized sector, the Artinian complete-symmetry presheaves,
      relative unstable Schlessinger reflections, stable reflected Noether
      representations, retractive differentiated algebroids, reflected-symmetry
      algebroids, base-changed Second Noether splittings, first derivatives, Taylor
      classes, syzygies, morphism-latching objects, tangent-representation
      Spencer chains, coefficient connections and their flatness/Bianchi fillers,
      intrinsic pro-coherent anchor-splitting curvature/Bianchi realizations, and
      their represented finite-locally-free comparisons;
\item complete presymplectic-descent, internal Cartan-lift, Cartan-diagram, and
      basicness realization objects; the associated reduced coefficients and
      sections are then reconstructed by effective Cartan descent.
\end{enumerate}
External variation categories and global point spaces carry the corresponding
restriction/descent maps; no stronger equivalence is inferred merely by applying
global points.
\end{theorem}

\begin{proof}
The underlying stable coefficient systems and structured cubical coefficients
satisfy finite spectral Nisnevich descent.  Symmetry lifting is a finite-limit
condition and right-Kan descent follows from the Cartesian-centre comma-category
lemma.  Both rectangular mixed grids are limits of relation degrees, and the
critical relation itself is reconstructed after descent from the critical
formalization.  Internal mapping objects, normalizations, partial totalizations,
rotated fibre sequences, and homotopy fibres preserve the resulting limit diagrams.

On the coherent formalized sector, a Nisnevich cover is spectral \'etale.  The
relative unstable exchange theorem proved above, together with the stable
Schlessinger, pro-coherent, formal-leaf, and formal-moduli descent theorems,
transports the reflected symmetry parameter, reflected
representations, partition-Lie differentiation, tangent representations, and
morphism-latching objects.  Once a Cartan lift or Cartan diagram has been constructed locally, effective
descent reconstructs the quotient coefficient, diagram, fibre, and morphism.
\end{proof}
\section{Audit of constructions and theorem scopes}
\label{sec:symmetry-noether-audit-of-constructions-and-theorem-scopes}

The chapter has three layers, and no fourth ``classical gauge-theory'' layer is
silently present.

\subsection{Unconditional native constructions}
\label{subsec:symmetry-noether-unconditional-native-constructions}

The formal symmetry anchor and effective quotient, symmetry--horizontal double
relation, stable cochains, complete bifiltration, structured symmetry-variation
categories, right-Kan coefficients, complete symmetry realizations, relative
Lepage restriction, universal Noether current, First Noether homotopy, native
symmetry-coherence obstruction tower, raw critical restriction followed by
formalization, parameterized critical quotient, mixed rectangular calculus and its
\v{C}ech identification with the reduced maximal relation, complete presymplectic
descent, Hamiltonian defect, residual transgression, radical realization, and the
internal Cartan realization objects are native constructions.  They require no represented
parameter bundle, no finite-order differential operator, and no integration by
parts.

\subsection{Formalized Artinian and differentiated constructions}
\label{subsec:symmetry-noether-formalized-artinian-and-differentiated-constructions}

On the locally coherent spectral coefficient sector, quotient formalization gives
the formalized relation and its Artinian leaves.  The first formalized Noether
coefficient is assembled on actual spectral Artinian tests and reflected by the
exact stable Schlessinger localization.  Complete formalized symmetry parameters
are restricted to the same Artinian leaf and reflected by the relative unstable
Schlessinger localization, whose coherent spectral-\'etale exchange was proved
above.  Their map to the zero locus of the reflected Noether
section is differentiated to the retractive partition-Lie algebroid, the zero and
Noether splittings, the first Noether derivative, the first syzygy, directional
Taylor classes, and the morphism corolla--latching hierarchy.  This is the domain of
the differentiated Brave-New Second Noether theorem.  The intrinsic
Noether--Spencer splitting object, tangent representation bar, canonical coefficient
connection with its flatness/Bianchi tower, and the intrinsic pro-coherent
anchor-splitting realization with its curvature/Bianchi hierarchy are refinements of
this same differentiated package.  The ordinary represented anchor-splitting space
is identified with that intrinsic realization only on the finite-locally-free
comparison locus.

\subsection{Represented geometric refinements}
\label{subsec:symmetry-noether-represented-geometric-refinements}

Representability is used only where the preceding chapters have supplied
quasi-coherent cotangent or tangent objects.  The first-order symmetry-lift theorem
then identifies the represented tangent shadow of the formal anchor.  On the
representably smooth scalar degree-zero critical sector, the affine-point linear
Euler zero locus gives
\[
T_{C_L/Q}^{\mathrm{lin}}
\simeq
\operatorname{fib}(\operatorname{Jac}^{\mathrm{lin}}_L),
\]
and on represented affine charts this becomes the ordinary tangent/Jacobi fibre.
The same represented first-order interface identifies critical symmetry directions
with Jacobi zero-modes.  The complete mixed presymplectic contraction remains a
relation-level coefficient; its separately constructed exact $(1,1)$
linearization produces the quasi-coherent contraction target
$\mathcal W_{\omega,T}^{(q)}$ and the represented presymplectic radical.  No
finite-order symmetry operator and no all-degree exterior-form identification is
inferred from these first-order statements.

\subsection{Explicitly deferred material}
\label{subsec:symmetry-noether-explicitly-deferred-material}

The chapter does not construct a general brave-new gauge field, principal gauge-field
connection, curvature of a gauge field, finite-order differential parameter action,
formal adjoint, integration-by-parts complex, superpotential, gauge-fixing datum,
an externally chosen reducibility complex, BRST differential, Koszul--Tate
resolution, BV complex, or a specific gauge theory such as Maxwell or Yang--Mills.
The native symmetry tower already constructs its obstruction/ambiguity/higher-loop
hierarchy, and the differentiated controller already constructs coefficient
connections together with intrinsic pro-coherent anchor-splitting curvature and
Bianchi coherence; on the finite-locally-free represented overlap it also identifies
the ordinary represented splitting realization with that intrinsic construction.
These are not substitutes for the additional gauge-field structures just listed.  The
remaining topics require additional geometric input and belong to later chapters.
Their future comparison theorems may use the infinitesimal symmetry and Noether
structures constructed here without changing the present definitions.
\section{The brave-new infinitesimal symmetry and Noether package}
\label{sec:symmetry-noether-the-brave-new-infinitesimal-symmetry-and-noether-package}

\begin{theorem}[Brave-new infinitesimal symmetry and Noether package]
\label{thm:symmetry-noether-brave-new-infinitesimal-symmetry-and-noether-package}
Let
\[
x:X\to S,\qquad Q=Q_{X/S},\qquad a:\overline A\to Q,
\]
let $M\in E_Q$, let
\[
L:\mathbf1_Q\to\operatorname{Sec}^0(A;M)[n],
\]
and let
\[
\mathcal G_\bullet\in\operatorname{FormLieGpd}(\overline A/Q).
\]
Then the constructions of this chapter canonically provide the following package.

\medskip
\noindent\emph{Formal symmetry and variational restriction.}
The canonical formal anchor
\[
\phi_{\mathcal G}:\mathcal G_\bullet\to R_{\overline A/Q,\bullet}
\]
and effective quotient
\[
q_{\mathcal G}:\overline A\twoheadrightarrow Q_{\mathcal G}
\]
together with the constructed effective arrow-image subrelation produce the
symmetry--horizontal double relation and the complete bifiltered stable
cochain object $\operatorname{Var}_{\mathcal G}(A;M)$.  Structured one- and
finite-cubical symmetry lifts form the categories
$\mathcal V_{\mathcal G}^{(k)}$ and their right-Kan coefficients
$V_{\mathcal G,M}^{(k)}(A)$, with canonical restriction maps from unrestricted
variation.  On represented first-order charts the lift space is classified by the
cotangent source coefficient and its first-order anchor; this comparison is not a
finite-order differential-operator realization.

\medskip
\noindent\emph{First Noether theory.}
For every horizontal stage $q\ge1$, the complete symmetry realization and the
relative Lepage realization meet over
\[
\mathcal N_{\mathcal G}^{q}(L)
=
\operatorname{Sym}_{\mathcal G}^{q}(L)
\times_Q
\operatorname{Lep}^{q}(L).
\]
Their difference defines the universal Noether current
\[
J_{\mathcal G,L}^{(q)}
=
\Theta_{\mathcal G,L}^{(q)}-B_{\mathcal G,L}^{(q)},
\]
and there is a canonical off-shell homotopy
\[
\boxed{
\iota_H^qJ_{\mathcal G,L}^{(q)}
\simeq
-\mathscr E_{\mathcal G,L}^{[q]}.
}
\]
The stagewise Dold--Kan fibres of the symmetry totalization give the complete native
symmetry-coherence obstruction tower.  The full Euler-nullification object promotes
the First Noether homotopy to the native complete-horizontal current transgression.

\medskip
\noindent\emph{Second Noether theory.}
On every coherent formal leaf $F$, the formalized first Noether coefficient has its
stable Schlessinger reflection
\[
\widehat{\mathcal O}_{\mathcal G,L}^{1,q}
\in\operatorname{Rep}^{\mathrm{fm}}(F),
\]
and complete formalized symmetry has its relative unstable reflection
\[
\widehat{\operatorname{Sym}}_{\mathcal G,L}^{\mathrm{Sch},q}\to F.
\]
The relative unstable reflector satisfies the coherent spectral-\'etale exchange
proved above.  The reflected symmetry parameter maps canonically to the zero locus of the reflected
Noether section.  Differentiation therefore gives a homotopy between the Noether
and zero splittings after base change to
$\mathfrak s_{\mathcal G,L}^{q}$ and hence the canonical nullhomotopy
\[
\boxed{
d_0\mathbf n_{\mathcal G,L}^{1,q}
\circ
U\vartheta_{\mathfrak s}
\simeq
0
}
\]
of the composite
$U\mathfrak s_{\mathcal G,L}^{q}\to
U\mathbb O_{\mathcal G,L}^{1,q}[n]$.
Independently, the first Noether derivative carries the canonical field-action
factorization
\[
d_0\mathbf n_{\mathcal G,L}^{1,q}
\simeq
\partial_{\mathrm{fld}}\mathbf n_{\mathcal G,L}^{1,q}
\rho_{\mathcal G}^{\pi,E_\infty},
\]
so the homotopy anchor fibre maps canonically to the first syzygy coefficient and
actual isotropy maps to it through the retained isotropy comparison.  Over the
reflected complete-symmetry parameter the Second Noether nullhomotopy coherently
nullifies every pulled-back exact labelled Taylor class while transporting the
complete Noether morphism-latching hierarchy to the zero hierarchy.  The universal
section-realization space gives the strong form on the entire controller whenever a
reflected complete-symmetry section is selected.  The differentiated Noether
representation also carries its canonical first-weight coefficient connection, with
curvature, Bianchi, and all higher flatness fillers supplied by its full tangent
representation action.  Independently, the differentiated anchor has an intrinsic
pro-coherent splitting realization with its own curvature and Bianchi hierarchy; on
the represented finite-locally-free comparison locus this is canonically identified
with the ordinary represented anchor-splitting realization.

\medskip
\noindent\emph{Critical geometry.}
Raw restriction of $\mathcal G_\bullet$ to the one-stage critical object is followed
by canonical unitwise formalization.  Its realization defines the formal
symmetry-reduced critical quotient.  On the representably smooth scalar degree-zero
sector,
\[
T_{C_L/Q}^{\mathrm{lin}}
\simeq
\operatorname{fib}(\operatorname{Jac}^{\mathrm{lin}}_L),
\]
and after a represented affine test
\[
T_{C_{L,T}/T}
\simeq
\operatorname{fib}(\operatorname{Jac}_{L,T}).
\]
Represented first-order critical symmetry directions therefore factor through the
Jacobi fibre and the critical-Hessian radical.

\medskip
\noindent\emph{Presymplectic and Hamiltonian geometry.}
The rectangular symmetry--vertical grid places the Noether current and the ambient
presymplectic transgression in one mixed coefficient system.  It gives the ordered
mixed contraction and the Hamiltonian defect
\[
\mathfrak H_{\mathcal G,L}^{(q)}
=
 d_V^{\mathrm{mix}}J_{\mathcal G,L}^{(q)}
-
\jmath_{\mathrm{tr}}\bigl(\iota_{\mathcal G}\omega_L^{(q)}\bigr).
\]
Hamiltonianity together with $d_VJ$-nullification determines a residual
transgression $\mathfrak r_{\mathcal G,L}^{(q)}$.  The derived object
$\operatorname{RadReal}_{\mathcal G,L}^{q}$ is precisely the compatible
nullification of that residual class and hence of the complete mixed
presymplectic transgression.  The complete mixed coefficient is not identified with
an ordinary bilinear form.  On the represented smooth scalar critical sector, the
ordered side-evaluation square and exact two-variable stabilization construct
\[
\omega_{L,T}^{(q),\flat}:
T_{C_{L,T}/T}\longrightarrow\mathcal W_{\omega,T}^{(q)},
\]
whose fibre is the represented first-order presymplectic radical.  Complete
symmetry-\v{C}ech descent of the closed critical presymplectic transgression constructs
$\omega_{L,\mathrm{red}}^{(q)}$ on the reduced maximal infinitesimal calculus and
forces critical symmetry directions into the represented radical; Cartan basicness
and null contraction are consequences of that descent.  Independently,
$\operatorname{RadReal}$ supplies the Hamiltonian residual nullification, and
$\operatorname{HamBasicReal}$ constructs the compatibility between the two
nullhomotopies.  Current Cartan-lift and Cartan-diagram realizations descend the
Noether current and its compatible transgression to the corresponding parameterized
effective quotient.

\medskip
\noindent\emph{Functoriality and locality.}
All native internal coefficient and moduli objects have the raw base-change,
formal-\'etale locality, and finite Nisnevich descent strength proved above.
External variation categories carry restriction functors.  The formalized
pro-coherent and Schlessinger-reflected branch has the separately scoped
spectral-\'etale varying-object equivalences of Chapters 14--15.
\end{theorem}

\begin{proof}
Every item has been constructed in the preceding sections from the canonical
formal-groupoid anchor, stable coefficient restriction, structured split
square-zero geometry, finite or complete limits, right Kan extension, stable
Dold--Kan, unstable and stable Schlessinger localization, relative infinite-loop
formation, partition-Lie differentiation, the morphism corolla--latching theorem,
parameterized zero and inverse loci, rotated stable fibre sequences, formalization
coreflection, and effective descent.  No finite-order gauge-theory operator or
classical integration-by-parts package occurs in the proof.
\end{proof}
\section{Handoff}
\label{sec:symmetry-noether-handoff}

The variational geometry is now connected to formal infinitesimal symmetry at the
level required by Noether theory.  Finite structured symmetry variations are actual
coherent lifting problems on the nondegenerate face-chain category.  The First
Noether theorem is intrinsically a complete-filtration current identity, and the
native symmetry-coherence tower records successive cosimplicial extension
obstructions together with their ambiguity torsors and higher automorphism spaces,
without introducing an external reducibility complex.

The differentiated branch is constructional at a different level.  The first
Artinian Noether coefficient is assembled on actual formal-leaf tests and reflected
by stable Schlessinger localization.  Complete formalized symmetries are restricted
to the same Artinian leaf and reflected independently by the relative unstable
Schlessinger localization.  Their canonical map to the zero locus of the reflected
Noether section differentiates to the homotopy between the Noether and zero
partition-Lie splittings.  The first derivative is separately shown to factor
through the differentiated field anchor by constructing its canonical nullhomotopy
on the homotopy anchor fibre and using the stable cofiber universal property.  This
is the Brave-New Second Noether theorem together with its intrinsic first syzygy;
its Taylor and morphism-latching consequences are inherited from the actual
differentiated morphism rather than imposed as a separate hierarchy of identities.

Critical symmetry is handled by raw restriction followed by canonical
formalization.  On the representably smooth scalar degree-zero sector, the Chapter
12 affine-point linear cotangent refinement constructs the critical tangent/Jacobi
fibre and hence the represented Hessian-degeneracy statement.  The mixed
symmetry--vertical calculus is rebuilt from explicit \v{C}ech rectangles of the
relevant maximal relations and supplies the ordered presymplectic contraction.
Hamiltonianity and $d_VJ$ nullification leave a residual class, and
$\operatorname{RadReal}$ is exactly the derived compatible nullification of that
residual class.  A separate exact $(1,1)$ first-order linearization constructs the
represented contraction target and radical without collapsing the raw degree-two
coefficient to exterior forms.  The critical mixed grid is the \v{C}ech resolution of the reduced maximal
infinitesimal relation, so the closed presymplectic transgression descends as a
complete calculus; basicness is its degeneracy shadow rather than an additional
hypothesis.  Current Cartan-lift objects, together with the Cartan-diagram
realization of the full current/nullhomotopy triangle, provide the corresponding
current descent data.

No theory of gauge fields has been smuggled into these constructions.  A future
brave-new gauge-theory development may introduce principal or higher gauge fields,
connections and curvature, finite-order parameter actions, reducibility,
formal-adjoint and superpotential technology, BRST/Koszul--Tate/BV structures, and
specific models such as Maxwell or Yang--Mills.  Those later chapters may take the
formal symmetry controller, Noether current, differentiated Noether splitting,
critical quotient, Hamiltonian defect, radical realization, and reduced
presymplectic/current families constructed here as inputs without reopening the
foundations of the present infinitesimal-symmetry theory.
